% =====================================================================
%  «Аштавакра-ґіта» — РОБОЧЕ ПІЛОТНЕ ВИДАННЯ
%  Шаблон верстки власного PDF через XeLaTeX.
%
%  Це НЕ завершена книга. Джерело тексту — власні редакторські картки
%  editorial/cards/*.json зі статусом «чернетка». Читань цифрових
%  свідків (GRETIL, SARIT) тут не відтворено: паспорт джерела дозволяє
%  PDF лише для нашого власного внеску.
%
%  Складання:
%    python3 scripts/build_pilot_book.py
%    cd books/ashtavakra-hita-pilot && latexmk -xelatex main.tex
%  або (без latexmk):
%    xelatex main.tex && makeindex main.idx && xelatex main.tex && xelatex main.tex
%
%  Потрібні пакети: fontspec, polyglossia, geometry, hyperref, makeidx,
%  longtable; шрифт DejaVu Serif (пакет fonts-dejavu-core).
% =====================================================================
\documentclass[10pt,a5paper,openany]{book}

\usepackage{fontspec}
\setmainfont{DejaVu Serif}[Ligatures=TeX]
\setsansfont{DejaVu Sans}[Ligatures=TeX,Scale=MatchLowercase]
\setmonofont{DejaVu Sans Mono}[Scale=MatchLowercase]

\usepackage{polyglossia}
\setdefaultlanguage{ukrainian}
\setotherlanguage{english}

\usepackage[a5paper,top=17mm,bottom=19mm,inner=17mm,outer=15mm,headsep=5mm]{geometry}
\usepackage{longtable}
\usepackage{makeidx}
% hyperref вмикаємо лише тоді, коли в TeX Live є TFM pzdr (пакет
% texlive-fonts-recommended). Без нього XeTeX-драйвер hyperref падає на
% \XeTeXLinkBox, тому краще скласти PDF без закладок, ніж не скласти зовсім.
\IfFileExists{pzdr.tfm}{\usepackage[hidelinks]{hyperref}}{}

\makeindex

% --- дрібні типографські зручності, без залежності від latex-extra -----
\setlength{\parindent}{1.1em}
\setlength{\parskip}{0pt}
\renewcommand{\baselinestretch}{1.04}
\clubpenalty=9000
\widowpenalty=9000
\brokenpenalty=3000
% A5 з довгими IAST-композитами: без запасу на розтяг рядок «лопається».
\tolerance=1500
\emergencystretch=3em
\hbadness=10000
\vbadness=10000
% Видання електронне й одностороннє: порожні сторінки від \cleardoublepage
% (у \mainmatter і \backmatter) тут не потрібні.
\let\cleardoublepage\clearpage

% --- макроси одного вірша ---------------------------------------------
% Верстка навмисно проста: жодного пакета поза texlive-latex-base,
% texlive-latex-recommended і texlive-xetex, щоб збірка в CI була
% відтворюваною без «магії».
\newcommand{\virhead}[1]{%
  \par\addvspace{1.5em}%
  {\sffamily\bfseries\large Вірш~#1\par}%
  \nopagebreak\addvspace{0.25em}}

\long\def\virs#1#2{\virhead{#1}\noindent\textbf{Переклад.}~#2\par}
\long\def\virlit#1{\addvspace{0.35em}\noindent\textit{Дослівно.}~#1\par}
\long\def\virgram#1{%
  \addvspace{0.4em}\noindent
  {\sffamily\footnotesize\textbf{Граматичний розбір.}}~%
  {\footnotesize\raggedright #1\par}}
\long\def\vircom#1{\addvspace{0.4em}\noindent\textbf{Коментар.}~#1\par}
\long\def\virvars#1{%
  \addvspace{0.4em}\noindent
  {\sffamily\footnotesize\textbf{Різночитання.}}\par
  \nopagebreak\begin{quote}\footnotesize\raggedright
  \setlength{\parskip}{0.25em}#1\end{quote}}
\long\def\virvaritem#1#2{\noindent\textit{#1}: #2\par}
\long\def\virvarnone#1{%
  \addvspace{0.4em}\noindent
  {\sffamily\footnotesize\textbf{Різночитання.}}~{\footnotesize\raggedright #1\par}}
\long\def\virsources#1#2{%
  \addvspace{0.3em}\noindent
  {\sffamily\footnotesize\textbf{Свідки (лише покажчик, текст не відтворено).}}~%
  {\footnotesize GRETIL \texttt{#1}; SARIT \texttt{#2}.}\par
  \addvspace{0.3em}\hrule\relax\vspace{0.1em}}

\newcommand{\pilotchapter}[2]{%
  \chapter*{Глава~#1}%
  \addcontentsline{toc}{chapter}{Глава~#1}%
  \markboth{Глава~#1}{Глава~#1}%
  \noindent{\sffamily\small У пілоті зверстано #2 карток цієї глави.\par}%
  \addvspace{0.4em}}

\begin{document}

% =====================================================================
%  ТИТУЛ
% =====================================================================
\begin{titlepage}
\centering
\vspace*{2.0cm}

{\sffamily\Large АШТАВАКРА-ҐІТА\par}
\vspace{0.6em}
{\sffamily\small Aṣṭāvakragītā\par}

\vspace{2.2cm}
{\sffamily\bfseries\large РОБОЧЕ ПІЛОТНЕ ВИДАННЯ\par}
\vspace{0.5em}
{\sffamily Глави 2 і 3. Власні редакторські картки\par}

\vspace{2.4cm}
\begin{minipage}{0.86\textwidth}
\small\raggedright
\textbf{Це не завершена книга і не академічне видання.}

\vspace{0.5em}
Це \emph{робоче пілотне видання}: воно показує, що верстальна машина
працює і що з наявного матеріалу вже можна скласти PDF. Книги як такої
в репозиторії немає.

\vspace{0.5em}
Усі редакторські картки, з яких зібрано цей текст, мають статус
\textbf{«чернетка»}. Жодна з них не позначена як рецензована, жодна не
має підтвердженого аудиту прав. Санскритські читання цифрових свідків
(GRETIL, SARIT) тут \textbf{не відтворено}.

\vspace{0.5em}
Не цитувати як переклад «Аштавакра-ґіти».
\end{minipage}

\vfill
{\sffamily\footnotesize Електронне робоче видання. Не для продажу.\par}
\end{titlepage}

% =====================================================================
%  ВИХІДНІ ДАНІ
% =====================================================================
\thispagestyle{empty}
\vspace*{\fill}
\begingroup\small
\noindent\textbf{\sffamily Вихідні дані}

\vspace{0.8em}
\noindent
\begin{tabular}{@{}p{0.34\textwidth}p{0.58\textwidth}@{}}
\textsf{Назва} & Аштавакра-ґіта (Aṣṭāvakragītā). Робоче пілотне видання \\[0.3em]
\textsf{Зміст} & Глави 2 і 3: український переклад, дослівне значення,
граматичний розбір, коментар, примітки про різночитання \\[0.3em]
\textsf{Первісний твір} & санскритська поема; за наявними даними —
суспільне надбання через вік \\[0.3em]
\textsf{Джерело нашого тексту} & власні редакторські картки
\texttt{editorial/}\allowbreak\texttt{cards/*.json} \\[0.3em]
\textsf{Статус карток} & «чернетка» — усі, без винятків \\[0.3em]
\textsf{Укладач} & редакція проєкту \texttt{juv4uk/ukrainian-books} \\[0.3em]
\textsf{Ліцензія власного внеску} & «ВОЛЬНІСТЬ» \\[0.3em]
\textsf{Тип видання} & робочий пілот, електронний, не для продажу \\[0.3em]
\textsf{Складання} & XeLaTeX, шрифт DejaVu Serif \\[0.3em]
\end{tabular}

\vspace{1.2em}
\noindent\footnotesize
Правовий статус матеріалів описано в \texttt{COPYRIGHT\_}\allowbreak\texttt{AND\_}\allowbreak\texttt{SOURCES.md},
видавничий висновок — у \texttt{docs/}\allowbreak\texttt{VYDAVNYCHYI\_}\allowbreak\texttt{VYSNOVOK.md}. Цей PDF не є
юридичним висновком і не є доказом правової чистоти.
\endgroup
\vspace*{\fill}

% =====================================================================
%  ПЕРЕДМОВА
% =====================================================================
\frontmatter
\pagestyle{headings}
\tableofcontents
\clearpage
\chapter*{Передмова}
\addcontentsline{toc}{chapter}{Передмова}
\markboth{Передмова}{Передмова}

Це \textbf{робоче пілотне видання}, а не книга. Мета одна: довести, що
верстальна машина справді складає український PDF із кирилицею та діакритикою
IAST, і показати, скільки власного матеріалу вже придатно до верстки.

\section*{Що тут зверстано}

Зверстано дві глави — другу й третю. Це єдині глави, де редакторські картки
покривають \emph{всі} вірші реєстру: 25 із 25 та 14 із 14. Разом 39 віршів.
Кожен вірш подано нашим власним внеском: українським перекладом, дослівним
значенням, граматичним розбором, коментарем і примітками про різночитання.

\section*{Чого тут немає}

Санскритського тексту віршів. Поля \texttt{sanskrit} і \texttt{iast} у картках
відтворюють читання цифрових транскрипцій GRETIL (CC BY-NC-SA 4.0) і SARIT
(CC BY-SA 3.0). Паспорт джерела дозволяє PDF лише для нашого власного внеску,
тому в пілоті залишено \emph{покажчик} на свідка (локатор вірша), а не його
текст. Це не технічне обмеження, а свідоме редакційне рішення; відкрите
правове питання описано у видавничому висновку.

Так само немає тут рецензії: усі 206 карток мають статус «чернетка»,
незалежного рецензента не залучено, санскритологічної вичитки не проведено.
Переклади й коментарі в цьому PDF — робочі, не вивірені.

\section*{Як це складається}

Текст генерується з карток скриптом \texttt{scripts/}\allowbreak\texttt{build\_}\allowbreak\texttt{pilot\_}\allowbreak\texttt{book.py},
а не пишеться руками. PDF складає XeLaTeX. Якщо картка зміниться, PDF
зміниться сам; якщо зміниться статус картки, скрипт зупиниться й попросить
переглянути пілот. Деталі — у \texttt{docs/}\allowbreak\texttt{BOOK\_}\allowbreak\texttt{PDF\_}\allowbreak\texttt{MACHINE.md}.

\mainmatter
% ЗГЕНЕРОВАНО scripts/build_pilot_book.py — не редагувати вручну.
% Джерело: editorial/cards/*.json, глави 2, 3, 4, 5, 6, 7, 8, 9, 10, 11, 12, 13, 14, 19.
% Містить лише власні поля карток: переклад, дослівне значення,
% граматичний розбір, коментар, примітки про різночитання.

\pilotchapter{2}{25}
\index{bodha@\textit{bodha}}\index{santa@\textit{śānta}}\index{niranjana@\textit{nirañjana}}\index{aho@\textit{aho}}
\virs{2.1}{О диво: я — незаплямлене, спокійне усвідомлення, вище за природу. Увесь цей час мене дурила сама лише омана.}
\virlit{«О диво: незаплямлене, спокійне усвідомлення — я, понад природу. Аж стільки часу я був одурений самою лише оманою.»}
\virgram{aho (вигук подиву) nirañjanaḥ (m. nom. sg. «незаплямлений») śāntaḥ (m. nom. sg. «спокійний») bodhaḥ (m. nom. sg. «усвідомлення») aham (nom. sg.) prakṛteḥ (f. abl./gen. sg. «від природи») paraḥ (m. nom. sg. «понад, вищий») | etāvantam (m. acc. sg. «аж стільки») aham kālam (m. acc. sg. «час») mohena (m. instr. sg. «оману») eva (підсилювальна частка) viḍambitaḥ (m. nom. sg. дієприкметник минулого від vi-ḍambh «обманювати, дурити»).}
\vircom{Розділ 2 відкривається вигуком подиву (aho) — це не аргумент, а зміна погляду. Три означення (nirañjana, śānta, bodha) ідуть як одне ціле: Я не «має» спокій, а є усвідомленням, яке не плямується. Viḍambitaḥ — не «я помилявся», а «мене дурили»: омана діяла, а не Я. Для українського перекладу важливо не зробити з цього самодокір: текст не звинувачує, а констатує.}
\virvarnone{власних приміток про різночитання немає; звіряння свідків не позначене як виконане}
\virsources{Avg\_2.1}{verse\_2.1}

\index{jagat@\textit{jagat}}\index{yatha@\textit{yathā}}\index{tatha@\textit{tathā}}\index{athava@\textit{athavā}}
\virs{2.2}{Як я один освітлюю це тіло, так освітлюю і світ. Тому весь світ — мій; або ж — аніщо.}
\virlit{«Як я, єдиний, освітлюю це тіло, так і світ. Тому весь світ — мій; або ж — аніщо.»}
\virgram{yathā … tathā (корелятивна пара «як … так») prakāśayāmi (1 sg. pres. pra-kāś «освітлювати, робити видимим») ekaḥ (m. nom. sg.) deham (m. acc. sg.) enam (вказівний займенник, acc. sg.) jagat (n. nom./acc. sg. «світ») | ataḥ (прислівник «тому») mama (gen. sg. «мій») jagat sarvam (n. nom. sg. «весь світ») athavā («або ж») na ca kiñcana («і аніщо»).}
\vircom{Вірш тримає дві відповіді одночасно: «світ мій» і «світ — аніщо». Це не суперечність, а два боки одного: світ не існує окремо від того, хто його освітлює, і саме тому не є чимось набутим. athavā («або ж») — сигнал, що текст свідомо лишає обидві формули, а не обирає одну. Перекладати його як «тобто» було б спрощенням.}
\virvarnone{власних приміток про різночитання немає; звіряння свідків не позначене як виконане}
\virsources{Avg\_2.2}{verse\_2.2}

\index{aho@\textit{aho}}
\virs{2.3}{О диво: тепер, відкинувши світ разом із тілом, я — самим лише вмінням, невідомо звідки — бачу вище Я.}
\virlit{«О диво: разом із тілом світ відкинувши мною тепер, якось самим умінням вище Я споглядається.»}
\virgram{saśarīram (прислівниковий прислівник «разом із тілом», sa- + śarīra) aho (вигук) viśvam (n. acc. sg.) parityajya (герундив від pari-tyaj «відкинувши») mayā (instr. sg. «мною») adhunā («тепер») | kutaścit («якось, звідкись») kauśālāt (n. abl. sg. «від уміння, вправності») eva (підсилювальна частка) paramātmā (m. nom. sg. «вище Я») vilokyate (3 sg. pass. vi-lok «споглядається, бачиться»). Пасивна конструкція: не «я споглядаю», а «Я бачиться».}
\vircom{Ключове тут — пасив vilokyate: бачення не є зусиллям того, хто бачить. kutaścit («невідомо звідки») підкреслює, що вміння не вироблене вправою. Для української редакції це найтонше місце: активне «я бачу» зраджує сенс, тому в перекладі додано «невідомо звідки» як відповідник kutaścit. saśarīram означає «не залишаючи тіла» — не фізичне відкидання, а відмова від ототожнення.}
\virvarnone{власних приміток про різночитання немає; звіряння свідків не позначене як виконане}
\virsources{Avg\_2.3}{verse\_2.3}

\index{atman@\textit{ātman}}\index{yatha@\textit{yathā}}\index{tatha@\textit{tathā}}
\virs{2.4}{Як хвилі, піна й бульбашки не відмінні від води, так і світ, що вийшов із Я, не відмінний від Я.}
\virlit{«Як не відмінні від води хвилі, піна й бульбашки, так не відмінний від Я світ, що вийшов із Я.»}
\virgram{yathā … tathā (корелятивна пара) na (заперечення) toyataḥ (n. abl. sg. від toya «вода»: «від води») bhinnāḥ (m. nom. pl. «відмінні») taraṅgāḥ (m. nom. pl. «хвилі») phena-budbudāḥ (m. nom. pl. «піна й бульбашки», dvandva) | ātmanaḥ (m. abl. sg. «від Я») na tathā bhinnam (n. nom. sg. «не відмінний») viśvam (n. nom. sg.) ātma-vinirgatam (n. nom. sg. «той, що вийшов із Я», дієприкметник від nir-gam з ātman-).}
\vircom{Образ води й хвиль працює на одному слові — bhinnam («відмінний»): хвиля не є водою в сенсі «іншою речовиною», але й не тотожна їй як окремий предмет. Це не тотожність і не різниця, а відмова від самої пари. ātmavinirgatam («той, що вийшов із Я») — не «створений», тому в перекладі немає слова «створений».}
\virvarnone{власних приміток про різночитання немає; звіряння свідків не позначене як виконане}
\virsources{Avg\_2.4}{verse\_2.4}


\virs{2.5}{Як тканина при дослідженні виявляється лише ниткою, так і цей світ, коли його дослідити, є лише Я.}
\virlit{«Лише ниткою виявилася б тканина, як при дослідженні; лише Я — це, так само світ досліджений.»}
\virgram{tantu-mātraḥ (m. nom. sg. «лише нитка», mātra — обмежувальний суфікс) bhavet (opt. 3 sg. bhū «була б, виявляється») eva paṭaḥ (m. nom. sg. «тканина») yadvat («як, подібно до того як») vicārataḥ (n. abl. sg. від vicāra «дослідження»: «при дослідженні») | ātma-tanmātram («лише те — Я», tanmātra — «саме те, тільки те») eva idam (n. nom. sg. «це») tadvat («так само») viśvam (n. nom. sg.) vicāritam (n. nom. sg. дієприкметник минулого від vi-car «досліджене»).}
\vircom{Це перший вірш книги, де прямо названо метод — vicāra, «дослідження». Тканина не зникає, коли ми бачимо нитку; зникає лише уявлення, що тканина — щось понад нитку. tanmātra («саме те, тільки те») — технічний термін, який у санкх'ї означає тонкі елементи; тут він ужитий у звичайному обмежувальному сенсі, і плутати ці два вживання не слід. Це спостереження про слововживання, а не встановлений факт.}
\virvars{\virvaritem{SARIT verse\_2.5 і GRETIL/TextGrid Avg\_2.5}{У кінці першого піввірша GRETIL/TextGrid подає vicāritaḥ (форма, що узгоджується з paṭaḥ), SARIT — vicārataḥ. Картка бере sarit-івське vicārataḥ і читає його як аблатив від vicāra («при дослідженні»); прочитання цієї форми лишається відкритим, друкованих свідків не звірено. Зіставлення GRETIL-боку зафіксоване в PR \#53.}\virvaritem{SARIT verse\_2.5 і GRETIL/TextGrid Avg\_2.5}{У другому піввірші GRETIL/TextGrid має ātmatanmatram (короткий a), SARIT — ātmatanmātram (довгий ā); картка бере довгий ā за SARIT. Розбіжність кількості голосного не вирішена третім свідком і не оголошується рукописним різночитанням. Зіставлення GRETIL-боку зафіксоване в PR \#53.}}
\virsources{Avg\_2.5}{verse\_2.5}

\index{yatha@\textit{yathā}}\index{tatha@\textit{tathā}}
\virs{2.6}{Як цукор уже вкладений у сік цукрової тростини й просякнутий ним, так і світ вкладений у мене й просякнутий мною безперервно.}
\virlit{«Як у соці тростини вкладений і ним просякнутий цукор, так світ у мені вкладений і мною просякнутий безперервно.»}
\virgram{yathā…tathā (корелятивна пара) iva (порівняльна частка) ikṣu-rase (m. loc. sg. «у соці цукрової тростини») kḷptā (f. nom. sg. «вкладена, влаштована») tena (n. instr. sg. «ним») vyāptā (f. nom. sg. «просякнута») iva śarkarā (f. nom. sg. «цукор») | tathā viśvam (n. nom. sg.) mayi (loc. sg. «у мені») kḷptam mayā (instr. sg.) vyāptam nirantaram (adv. «безперервно»).}
\vircom{Образ працює не на «створенні», а на присутності: цукор не доданий до соку — він у ньому вже є, і окремо його не видно. Ключове слово — kḷpta («вкладений, влаштований»), а не sṛṣṭa («створений»); підміна змінила б думку з «світ у мені влаштований» на «світ мною зроблений». Тому в перекладі немає слова «створений».}
\virvarnone{власних приміток про різночитання немає; звіряння свідків не позначене як виконане}
\virsources{Avg\_2.6}{verse\_2.6}

\index{jagat@\textit{jagat}}
\virs{2.7}{Від незнання Я світ сяє; від знання Я він не сяє. Від незнання мотузки сяє змія; від знання її вона не сяє.}
\virlit{«Від незнання Я світ сяє; від знання Я не сяє. Від незнання мотузки змія сяє; від знання її не сяє.»}
\virgram{ātma-ajñānāt (n. abl. sg. «від незнання Я») jagat (n. nom. sg.) bhāti (3 sg. pres. bhā «сяє») ha (підсилювальна частка) ātma-jñānāt (abl.) na bhāsate (3 sg. med. «не сяє») | rajju-ajñānāt (n. abl. sg. «від незнання мотузки») ahiḥ (m. nom. sg. «змія») bhāti taj-jñānāt (n. abl. sg. «від знання її») bhāsate na hi (hi — причинова частка).}
\vircom{Класичний приклад мотузки й змії, але з точністю, яку легко зіпсувати: текст каже na bhāsate — «не сяє», а не «зникає». Зникає не мотузка, а змія; світ не анульовується, а перестає сяяти як щось окреме. Українське «зникає» тут було б зручнішим і хибним.}
\virvarnone{власних приміток про різночитання немає; звіряння свідків не позначене як виконане}
\virsources{Avg\_2.7}{verse\_2.7}


\virs{2.8}{Світло — моя власна природа; я не більший за нього. Коли сяє світ, сяє лише сяйво Я.}
\virlit{«Світло — моя власна природа; я не більший за нього. Коли сяє світ, сяє саме сяйво Я.»}
\virgram{prakāśaḥ (m. nom. sg. «світло, сяйво») me (gen. sg. «моє») nijam rūpam (n. nom. sg. «власна природа») na atiriktaḥ (m. nom. sg. «не більший, не надлишковий») asmi aham tataḥ («я від того») | yadā … tadā (корелятивна пара) prakāśate (3 sg. med. «сяє») viśvam (n. nom. sg.) ahaṃ-bhāsaḥ (m. nom. sg. «сяйво Я») eva hi.}
\vircom{Тонкість у nātiriktaḥ — «не більший, не надлишковий»: світло не є чимось поза своїм проявом, і того, хто сяє, не можна відділити від сяйва. Складник ahaṃbhāsa рідкісний; перекладено «сяйво Я», а не «его» — его було б ahaṃkāra, інше слово.}
\virvarnone{власних приміток про різночитання немає; звіряння свідків не позначене як виконане}
\virsources{Avg\_2.8}{verse\_2.8}

\index{aho@\textit{aho}}\index{yatha@\textit{yathā}}
\virs{2.9}{О диво: вигаданий світ через незнання сяє в мені — як срібло в перламутрі, як змій у мотузці, як вода в мареві сонячного променя.}
\virlit{«О диво: вигаданий світ через незнання сяє в мені — як срібло в перламутрі, як змій у мотузці, як вода в сонячному промені.»}
\virgram{aho (вигук «о диво») vikalpitam (n. nom. sg. «вигадане, сконструйоване») viśvam (n. nom. sg.) ajñānāt (n. abl. sg. «через незнання») mayi (loc. sg. «у мені») bhāsate («сяє»). Далі три порівняння в одній конструкції yathā: rūpyam (n. «срібло») śuktau (f. loc. sg. «у перламутрі, в устричній черепашці»), phaṇī (m. «змій») rajjau (f. loc. sg. «у мотузці»), vāri (n. «вода») sūryakare (m. loc. sg. «у сонячному промені», тобто в мареві). Усі три — locativus loci: місце, де вигадане сяє.}
\vircom{Три приклади підібрані за одним принципом: у кожному є річ, що є, і річ, якої немає, але яка видима. Перламутр є, срібла немає; мотузка є, змія немає; промінь є, води немає. Спільне в них не «світ ілюзорний», а точніше: **помилка має місце в тому, що є, і не додає нічого нового до буття**. Слово vikalpitam («сконструйоване») вказує на роботу розуму, а не на онтологічний статус світу. Порівняй 2.10, де той самий світ описаний уже як такий, що вийшов із мене й повернеться: там акцент на походженні, тут — на способі явлення.}
\virvars{\virvaritem{SARIT і GRETIL (перевірено scripts/fetch\_witnesses.py)}{Читання збігаються; різниця лише в поділі слів (viśvamajñānān / viśvam ajñānān). Не різночитання.}}
\virsources{Avg\_2.9}{verse\_2.9}


\virs{2.10}{Світ, що вийшов із мене, у мені ж і розчиниться — як горщик у глині, як хвиля у воді, як браслет у золоті.}
\virlit{«Світ, що вийшов від мене, у мені ж прийде до розчинення — як горщик у глині, як хвиля у воді, як браслет у золоті.»}
\virgram{mattaḥ (abl. sg. «від мене») vinirgatam (n. nom. sg. «вийшло», дієприкметник від nir-gam) viśvam (n. nom. sg.) mayi (loc. sg.) eva («саме») layam eṣyati (fut. 3 sg. i- з layam «прийде до розчинення»). Три порівняння: mṛdi (f. loc. sg. «у глині») kumbhaḥ (m. «горщик»), jale (n. loc. sg. «у воді») vīciḥ (m./f. «хвиля»), kanake (n. loc. sg. «у золоті») kaṭakam (n. «браслет»). Усі три пари — причина й наслідок, де наслідок не є чимось іншим за причину.}
\vircom{Порівняння тут іншого типу, ніж у 2.9. Там ішлося про помилку сприйняття, тут — про **відношення причини й наслідку**. Горщик не є чимось додатковим до глині, браслет — до золота: змінилася форма, не речовина. Це класична для індійської думки модель, де наслідок уже міститься в причині. Слово laya («розчинення») не означає знищення: воно означає повернення до того, з чого виникло. Разом 2.9 і 2.10 дають два різні пояснення одного й того самого: світ і помилковий у способі явлення (2.9), і не відокремлений від джерела (2.10).}
\virvars{\virvaritem{SARIT і GRETIL (перевірено scripts/fetch\_witnesses.py)}{Читання збігаються; різниця лише в поділі слів (mayyeva / mayy eva, layameṣyati / layam eṣyati). Не різночитання.}}
\virsources{Avg\_2.10}{verse\_2.10}

\index{jagat@\textit{jagat}}\index{aho@\textit{aho}}\index{bahuvrihi@\textit{bahuvrīhi}}\index{namah@\textit{namaḥ}}
\virs{2.11}{О диво! Поклін мені самому — тому, в кого немає загибелі, хто стоїть, навіть коли гине світ від Брахми аж до стебла трави.}
\virlit{«О диво! Я — поклін собі! У мене, в кого немає загибелі, — того, хто стоїть, навіть коли гине світ від Брахми аж до стебла трави.»}
\virgram{aho (вигук) aham («я») namaḥ (n. nom. sg. «поклін») mahyam (dat. sg. «мені самому»; давальний при namaḥ). vināśaḥ (m. nom. sg. «загибель») yasya (gen. sg.) na asti me (gen. sg. «у мене»). brahmādistambaparyantam (bahuvrīhi-складник: brahma-ādi «починаючи з Брахми» + stamba «стебло трави» + paryanta «аж до»; узгоджене з jagat) jagat (n. nom. sg. «світ»). jagannāśe 'pi (loc. sg. «навіть коли світ гине») tiṣṭhataḥ (gen. sg. дієприкметника теперішнього часу sthā- «того, хто стоїть», узгодженого з me).}
\vircom{Зачин «aho ahaṃ namo mahyaṃ» — самопрославлення, і це навмисний парадокс у тексті, що весь будується на запереченні «я». Тут «я» не его, а те, що лишається, коли его знято: тому й поклін спрямований не собі як особі, а тому, що в собі не гине. Діапазон «brahmādistambaparyantam» («від Брахми до стебла трави») — це не поетичне перебільшення, а повний обсяг живих істот у традиційній картині: найвищий бог і найменша рослина. Твердження сильне саме тому, що охоплює обидва краї. Те саме «namo mahyaṃ» повторюється в 2.12 — це рефрен, а не випадковий збіг.}
\virvars{\virvaritem{SARIT і GRETIL (перевірено scripts/fetch\_witnesses.py)}{Читання збігаються; відмінності лише в аваґрасі (nāśe+api / nāśe 'pi) і в комі після mahyaṃ у GRETIL. Не різночитання.}}
\virsources{Avg\_2.11}{verse\_2.11}

\index{aho@\textit{aho}}\index{bahuvrihi@\textit{bahuvrīhi}}
\virs{2.12}{О диво! Поклін мені самому! Я один — і хоч маю тіло, нікуди не йду й не приходжу: охопивши все, перебуваю.}
\virlit{«О диво! Я — поклін собі! Я один, і хоч маю тіло, нікуди не йду й не приходжу — охопивши все, перебуваю.»}
\virgram{aho ahaṃ namo mahyaṃ — той самий зачин, що в 2.11. ekaḥ (m. nom. sg. «один») aham (nom. sg.). dehavān (bahuvrīhi «той, хто має тіло», від deha + -vant) api («навіть, хоч»). kvacit («кудись, десь») na gantā (m. nom. sg. дієприкметника майбутнього часу від gam- «той, хто піде») na āgantā («не той, хто прийде»). vyāpya (герунд vi-āp «наповнивши, охопивши») viśvam (n. acc. sg.) avasthitaḥ («розташований, перебуває»).}
\vircom{Парадокс у «dehavān api» («навіть маючи тіло»): тіло є, а руху немає, бо немає куди йти — усе вже охоплене. «na gantā nāgantā» («не той, хто йде, і не той, хто приходить») знімає не дію, а подорож як таку: обидва дієприкметники майбутнього часу, тобто йдеться не про те, що вже сталося, а про те, що не станеться. Разом із 2.11 це пара: там «не гину, коли гине світ» (час), тут «не йду, коли все охоплене» (простір). Два виміри, один висновок.}
\virvars{\virvaritem{SARIT і GRETIL (перевірено scripts/fetch\_witnesses.py)}{Читання збігаються; відмінності лише в аваґрасі (eko+ahaṃ / eko 'haṃ) і в поділі слів. Не різночитання.}}
\virsources{Avg\_2.12}{verse\_2.12}

\index{aho@\textit{aho}}\index{bahuvrihi@\textit{bahuvrīhi}}
\virs{2.13}{О диво! Поклін мені самому — тому, хто, не торкаючись тілом, довго втримує все; немає тут нікого, хто дорівнював би мені в цьому.}
\virlit{«О диво! Я — поклін собі! Немає тут нікого, хто дорівнював би мені спритністю — того, ким, не торкаючись тілом, довго втримується все.»}
\virgram{Зачин aho ahaṃ namo mahyaṃ — той самий рефрен, що в 2.11 і 2.12. dakṣaḥ (m. nom. sg. «спритний, здатний») na asti («немає») iha («тут») mat-samaḥ (bahuvrīhi «рівний мені», mat = abl. основа від aham). asaṃspṛśya (герунд a-saṃ-spṛś «не торкнувшись») śarīreṇa (n. instr. sg. «тілом»). yena (m. instr. sg. відносного займенника, узгодженого з «мною»: «ким [я], що...») viśvam (n. acc. sg.) ciram («довго») dhṛtam (n. nom./acc. sg. дієприкметника минулого часу dhṛ- «утримуване»).}
\vircom{Парадокс в asaṃspṛśya («не торкнувшись»): світ утримується, і водночас немає дотику. Тіло назване тут не як перешкода, а як те, чим дотик не відбувається — це не заперечення тіла, а заперечення його як інструменту. Слово dakṣa («спритний») ужите з іронією: спритність зазвичай означає вправність у дії, а тут вона полягає в тому, щоб діяти без дотику. Третій рефренний зачин поспіль (2.11, 2.12, 2.13) показує, що це не окремі вигуки, а серія: кожен наступний вірш додає новий вимір до того самого «поклону собі».}
\virvars{\virvaritem{SARIT і GRETIL (перевірено scripts/fetch\_witnesses.py)}{Читання збігаються; відмінності лише в аваґрасі та поділі слів. Не різночитання.}}
\virsources{Avg\_2.13}{verse\_2.13}

\index{aho@\textit{aho}}\index{bahuvrihi@\textit{bahuvrīhi}}\index{athava@\textit{athavā}}
\virs{2.14}{О диво! Поклін мені самому! У мене немає нічого — або ж у мене є все, що взагалі досяжне слову й розуму.}
\virlit{«О диво! Я — поклін собі! У мене немає анічого — або ж у мене все, що досяжне слову й розуму.»}
\virgram{aho ahaṃ namo mahyaṃ — рефрен. yasya (gen. sg.) me (gen. sg.) na asti («немає») kiñcana («чогось, анічого»). athavā («або ж») — сполучник альтернативи, що не обирає сторону. yasya me («у кого — у мене») sarvam (n. nom. sg. «все»). yad vāṅmanasagocaram (bahuvrīhi: vāc «слово» + manas «розум» + gocara «поле, досяжність»; «те, що є полем слова й розуму»).}
\vircom{Вірш ставить дві взаємовиключні відповіді й не обирає між ними: athavā («або ж») залишає обидві. Це не нерішучість, а точний хід: якщо «нічого» і «все» однаково хибні описи, то обидва виражають те саме — те, що не піддається опису ні як володіння, ні як відсутність. Та сама формула athavā na ca kiñcana є і в 2.2, тож це не одноразовий зворот, а повторюваний прийом тексту. Друга половина обмежує коло: sarvaṃ — це не «все взагалі», а лише те, що потрапляє в поле слова й розуму (vāṅmanasagocara). Те, що поза цим полем, не називається ні «всім», ні «нічим».}
\virvars{\virvaritem{SARIT і GRETIL (перевірено scripts/fetch\_witnesses.py)}{Читання збігаються; відмінності лише в аваґрасі та поділі слів. Не різночитання.}}
\virsources{Avg\_2.14}{verse\_2.14}

\index{niranjana@\textit{nirañjana}}\index{tatha@\textit{tathā}}
\virs{2.15}{Знання, пізнаване й пізнавач — ця трійця не є справжньою. Те, у чому вона сяє через незнання, — це я, незаплямований.}
\virlit{«Знання, пізнаване й пізнавач — ця трійця не є дійсною. Те, в чому вона сяє через незнання, — це я, незаплямований.»}
\virgram{jñānam (n. nom. sg. «знання») jñeyam (n. nom. sg. «пізнаване») tathā («так само») jñātā (m. nom. sg. «пізнавач») — tritayam (n. nom. sg. «трійця»). na asti («немає») vāstavam (n. nom. sg. «справжнього, дійсного»; присудкове до tritayam). ajñānāt (n. abl. sg. «через незнання») bhāti («сяє»). yatra («у чому, де») idam (n. nom. sg.). saḥ (m. nom. sg.) aham (nom. sg.) asmi («я є») nirañjanaḥ (m. nom. sg. «незаплямований»).}
\vircom{Заперечується не знання як таке, а **трійця** — поділ на того, хто знає, те, що знають, і саме знання. Це найточніше місце глави: доки є трійця, є й поділ, а поділ і є те, що тут названо незнанням. Зверни увагу на порядок: спершу сказано, що трійця не дійсна, і лише потім — що вона сяє. Тобто вона не зникає: вона продовжує являтися, але вже не як справжня. saḥ aham asmi («я є те») — пряма формула ототожнення, і вона стосується не трійці, а того, у чому трійця являється. nirañjana («незаплямований») стоїть останнім словом навмисно: те, у чому сяє незнання, саме ним не забруднюється.}
\virvars{\virvaritem{SARIT і GRETIL (перевірено scripts/fetch\_witnesses.py)}{Читання збігаються; відмінності лише в аваґрасі та поділі слів. Не різночитання.}}
\virsources{Avg\_2.15}{verse\_2.15}

\index{cit@\textit{cit}}\index{aho@\textit{aho}}\index{bahuvrihi@\textit{bahuvrīhi}}\index{duhkha@\textit{duḥkha}}
\virs{2.16}{О диво! Корінь страждання — двоїстість, і немає від нього інших ліків. Усе це видиме — хибне. Я один: смак свідомості, незаплямований.}
\virlit{«О диво! Страждання має корінь у двоїстості, і немає від нього інших ліків. Усе це видиме — хибне. Я один — смак свідомості, незаплямований.»}
\virgram{dvaitamūlam (bahuvrīhi: dvaia «двоїстість» + mūla «корінь»; «чий корінь — двоїстість») aho («о диво») duḥkham (n. nom. sg. «страждання»). na anyat («немає іншого») tasya (n. gen. sg. «від нього») asti bheṣajam (n. nom. sg. «ліки»). dṛśyam (n. nom. sg. «видиме») etat (n. nom. sg.) mṛṣā («хибне, марне»; присудкове) sarvam (n. nom. sg. «усе»). ekaḥ (m. nom. sg. «один») aham (nom. sg.) cidrasaḥ (bahuvrīhi: cit «свідомість» + rasa «смак, сутність»; «чий смак — cit») amalaḥ (m. nom. sg. «незаплямований»).}
\vircom{Вірш має три частини й будується як діагноз, рецепт і твердження. Діагноз: duḥkha має корінь у двоїстості, а не в обставинах. Рецепт: nānyat tasya asti bheṣajam — «немає інших ліків», тобто лікування не в усуненні окремих страждань, а в усуненні поділу. Твердження: dṛśyam etan mṛṣā — видиме хибне. Останнє найризикованіше, і його не слід читати як «світу немає»: mṛṣā означає «марне, несправжнє» в сенсі несамостійного буття, а не відсутності. cidrasaḥ («чий смак — cit») — рідкісний і точний складник: не «я є свідомість», а «я маю смак свідомості», тобто свідомість не річ, якою я володію, а те, чим я відчувається.}
\virvars{\virvaritem{SARIT і GRETIL (перевірено scripts/fetch\_witnesses.py)}{Читання збігаються; відмінності лише в аваґрасі та поділі слів. Не різночитання.}}
\virsources{Avg\_2.16}{verse\_2.16}

\index{bodha@\textit{bodha}}
\virs{2.17}{Я — лише усвідомлення; через незнання я домисив собі обмеження. Хто так постійно розмірковує, той стоїть у не-розрізненні — і це моє стояння.}
\virlit{«Я — лише усвідомлення; через незнання мною домислене обмеження. Так тому, хто постійно розмірковує, — моє стояння в не-розрізненні.»}
\virgram{bodha-mātraḥ (m. nom. sg. «лише усвідомлення», mātra — обмежувальний суфікс) aham | ajñānāt (n. abl. sg.) upādhiḥ (m. nom. sg. «обмеження, накладка») kalpitaḥ (m. nom. sg. «домислений, сконструйований») mayā (instr. sg.) | evam vimṛśataḥ (gen. sg. дієприкм. від vi-mṛś «того, хто так розмірковує») nityam (adv. «постійно») nirvikalpe (loc. sg. «у не-розрізненні») sthitiḥ mama («моє стояння»).}
\vircom{Тут названо механізм: upādhi («накладка, обмеження») не є чимось зовнішнім — він kalpita, «домислений», і домисил його саме той, хто тепер його знімає (mayā). Це не докір собі: у тексті немає оцінки, лише констатація авторства помилки. nirvikalpa — «без розрізнення», стан, у якому vikalpa (конструювання) не діє; це той самий корінь, що в 2.9 vikalpita.}
\virvarnone{власних приміток про різночитання немає; звіряння свідків не позначене як виконане}
\virsources{Avg\_2.17}{verse\_2.17}

\index{aho@\textit{aho}}
\virs{2.18}{О диво: світ ніби в мені — а насправді він у мені не вміщується. Немає в мене ні пут, ні визволення; омана вщухла, втративши опору.}
\virlit{«О диво: у мені розташований світ — а насправді він у мені не розташований. Немає в мене ні пут, ні визволення; омана заспокоєна, без опори.»}
\virgram{aho mayi (loc. sg.) sthitam (n. nom. sg. «розташоване») viśvam (n. nom. sg.) vastutaḥ (adv. «насправді, сутнісно») na mayi sthitam | na me bandhaḥ asti mokṣaḥ vā bhrāntiḥ (f. nom. sg. «омана») śāntā (f. nom. sg. «заспокоєна») nirāśrayā (f. nom. sg. «без опори»).}
\vircom{Парадокс навмисний: перший піввірш дає звичну формулу, другий її знімає через vastutaḥ («насправді»). Далі — те саме заперечення, що в 15.18: немає ні пут, ні визволення. nirāśrayā («без опори») вказує, чому омана вщухла: у неї не лишилося основи, а не тому, що її перемогли.}
\virvarnone{власних приміток про різночитання немає; звіряння свідків не позначене як виконане}
\virsources{Avg\_2.18}{verse\_2.18}


\virs{2.19}{Цей світ разом із тілом — ніщо, так установлено; а Я — лише чиста свідомість. У чому ж тоді ще домислювати?}
\virlit{«Цей світ разом із тілом — ніщо, так установлено; а Я — лише чиста свідомість. У чому ж тоді лишається домислювання?»}
\virgram{saśarīram (adv. «разом із тілом») idam viśvam (n. nom. sg.) na kiñcit («ніщо») iti niścitam («так установлено») | śuddha-cin-mātraḥ (m. nom. sg. «лише чиста свідомість») ātmā (m. nom. sg.) ca tat kasmin (loc. sg. «у чому ж тоді») kalpanā adhunā («домислювання тепер»).}
\vircom{Друга половина — риторичне питання, і воно важливіше за першу: після того як названо «ніщо» і «чиста свідомість», для домислювання не лишається місця (kasmin — «у чому?»). śuddha-cin-mātra зберігає потрійну будову: чисте — свідомість — лише; кожен складник несе вагу.}
\virvarnone{власних приміток про різночитання немає; звіряння свідків не позначене як виконане}
\virsources{Avg\_2.19}{verse\_2.19}

\index{cit@\textit{cit}}\index{tatha@\textit{tathā}}
\virs{2.20}{Тіло, небо й пекло, пута й визволення, страх — усе це лише домислювання. Що мені робити, коли моя природа — свідомість?}
\virlit{«Тіло, небо й пекло, пута й визволення, страх — усе це лише домислювання. Що мені, чиїй природі — свідомість, робити?»}
\virgram{śarīram (n. nom. sg.) svarga-narakau (m. du. «небо й пекло») bandha-mokṣau (m. du. «пута й визволення») bhayam (n. nom. sg.) tathā | kalpanā-mātram (n. nom. sg. «лише домислювання») eva etat («усе це») kim me kāryam («що мені робити») cit-ātmanaḥ (m. gen. sg. «тому, чия природа — свідомість»).}
\vircom{Два двоїни (dual) у переліку — svarga-narakau і bandha-mokṣau — тримають пари разом: текст не відкидає пекло й не обирає небо, він ставить обидва в один ряд із домислюванням. Питання kiṃ me kāryam не риторичне в сенсі «нічого»: воно справді питає, що лишається робити, і відповіді в цьому вірші немає.}
\virvarnone{власних приміток про різночитання немає; звіряння свідків не позначене як виконане}
\virsources{Avg\_2.20}{verse\_2.20}

\index{aho@\textit{aho}}
\virs{2.21}{О диво! Навіть у натовпі, коли я не бачу двоїстості, все стає як ліс. Де ж мені шукати втіхи?}
\virlit{«О диво! Навіть у натовпі, коли я не бачу двоїстості, [усе] стало як ліс. Де ж мені знайти втіху?»}
\virgram{aho (вигук) janasamūhe (m. loc. sg. «у скупченні людей, у натовпі») api («навіть»). na («не») dvaitam (n. acc. sg. «двоїстість») paśyataḥ (gen. sg. дієприкметника теперішнього часу dṛś- «того, хто бачить», узгодженого з mama) mama (gen. sg. «мене»). araṇyam (n. nom. sg. «ліс») iva («як») saṃvṛttam (n. nom. sg. дієприкметника минулого часу saṃ-vṛt «стало»). kva («де») ratim (f. acc. sg. «втіху, насолоду») karavāṇi (1 sg. opt./conj. kṛ- «мені робити»; питальна конструкція з kva). aham (nom. sg.).}
\vircom{Парадокс подвійний. Перший: натовп для того, хто не бачить двоїстості, — те саме, що ліс, тобто самотність не залежить від людей навколо. Другий, і він тонший: питання kva ratiṃ karavāṇi («де мені знайти втіху») поставлене там, де втіху вже ніде шукати, бо поділу немає. Це не скарга, а риторичне питання, що закриває можливість відповіді. Форма optativus karavāṇi («мені б робити») додає відтінку роздуму, а не бажання. Зверни увагу: слово rati («втіха») — те саме, що в 15.2 у складі vaiṣayiko rasaḥ («смак до чуттєвого»); тут воно вжите без засудження, але й без місця, куди його покласти.}
\virvars{\virvaritem{SARIT і GRETIL (перевірено scripts/fetch\_witnesses.py)}{Читання збігаються; відмінності лише в аваґрасі та поділі слів. Не різночитання.}}
\virsources{Avg\_2.21}{verse\_2.21}

\index{atman@\textit{ātman}}\index{cit@\textit{cit}}\index{jiva@\textit{jīva}}
\virs{2.22}{Я не тіло, і тіло не моє. Я не живчик — я свідомість. Ось у чому була моя уза: спрага до життя.}
\virlit{«Я не тіло, і тіло не моє; я не живчик — я ж свідомість. Ось моя уза: [то була] спрага до життя.»}
\virgram{na aham («я не») dehaḥ (m. nom. sg. «тіло»). na me dehaḥ («тіло не моє»; me — gen. sg.). jīvaḥ (m. nom. sg. «живчик, індивідуальна душа») na aham («я не [він]»). aham (nom. sg.) hi («ж») cit (n. nom. sg. «свідомість»; присудкове). ayam (m. nom. sg.) eva hi me (gen. sg. «моя») bandhaḥ (m. nom. sg. «уза»). āsīt (3 sg. imperfect as- «була»; минулий незавершений час). yā (f. nom. sg. відносного займенника) jīvite (n. loc. sg. «у житті») spṛhā (f. nom. sg. «спрага, прагнення»).}
\vircom{Найважливіше тут — **час дієслова**. āsīt — це минулий незавершений час, тобто «була», а не «є». Уза не описується як наявна: вона названа в минулому, бо вже побачена як те, чим вона була. Це не проголошення свободи, а звіт про те, що побачене. Другий точний момент: заперечуються три ототожнення поспіль (тіло, володіння тілом, jīva), і лише четверте — ствердне (aham cit). jīva тут не синонім ātman: це саме «живчик», індивідуальна істота, і її заперечення стоїть окремо від заперечення тіла. Причина узи названа одна — spṛhā, спрага до життя, а не до насолод. Це вужче й точніше, ніж звичне «бажання».}
\virvars{\virvaritem{SARIT і GRETIL (перевірено scripts/fetch\_witnesses.py)}{Читання збігаються; відмінності лише в аваґрасі та поділі слів. Не різночитання.}}
\virsources{Avg\_2.22}{verse\_2.22}

\index{aho@\textit{aho}}
\virs{2.23}{О диво! Різнобарвні хвилі світів раптом піднялися — бо в мені, безмежному великому океані, здіймається вітер розуму.}
\virlit{«О диво! Різнобарвні хвилі світів швидко піднялися — коли в мені, безмежному великому океані, здіймається вітер розуму.»}
\virgram{aho (вигук) bhuvanakallolaiḥ (m. instr. pl. «хвилями світів»: bhuvana «світ» + kallola «хвиля»). vicitraiḥ (m. instr. pl. «різнобарвними, різноманітними»). drāk («швидко, раптом»). samutthitam (n. nom. sg. дієприкметника минулого часу sam-ut-thā «піднялося»). mayi (loc. sg. «у мені») anantamahāmbhodhau (m. loc. sg. «у безмежному великому океані»: ananta «безмежний» + mahā «великий» + ambhodhi «океан»). cittavāte (m. loc. sg. «коли вітер розуму»: citta «розум, свідомість» + vāta «вітер») samudyate (3 sg. pr. sam-ud-i «здіймається»).}
\vircom{Початок триптиха (2.23–2.25), де той самий образ розгортається в три кроки: вітер здіймається, вітер вщухає, хвилі граються самі. Тут важлива **причинність**: хвилі не є океаном і не породжені ним безпосередньо — їх піднімає вітер, тобто citta. Океан лишається океаном і без хвиль. Слово ambhodhi буквально означає «той, що містить води» — це поетичне іменування океану, і воно тримає образ: води багато, хвиля — її форма, а не додаток. Порівняй 2.10, де були горщик у глині й хвиля у воді: там ішлося про невідокремленість наслідку, тут — про те, що рух має причину, і причина не в океані.}
\virvars{\virvaritem{SARIT і GRETIL (перевірено scripts/fetch\_witnesses.py)}{Читання збігаються; відмінності лише в аваґрасі та поділі слів. Не різночитання.}}
\virsources{Avg\_2.23}{verse\_2.23}

\index{jagat@\textit{jagat}}\index{jiva@\textit{jīva}}
\virs{2.24}{Коли в мені, безмежному океані, вітер розуму вщухає, — на лихо торговцеві-живчику — човен світу гине.}
\virlit{«Коли в мені, безмежному океані, вітер розуму вщухає — на лихо торговця-живчика, човен світу гине.»}
\virgram{mayi anantamahāmbhodhau (див. 2.23). cittavāte (m. loc. sg.) praśāmyati (3 sg. pr. pra-śam «вщухає, заспокоюється»; абсолютна конструкція: «коли вітер розуму вщухає»). abhāgyāt (n. abl. sg. «через нещастя, на лихо»). jīvavaṇijaḥ (m. gen. sg. «торговця-живчика»: jīva + vaṇija «торговець»). jagatpotaḥ (m. nom. sg. «човен світу»: jagat «світ» + pota «човен, судно»). vinaśvaraḥ (m. nom. sg. «той, що гине, тлінний»).}
\vircom{Іронія тут точна й гірка. abhāgyāt («на лихо») сказано з погляду **торговця**: для jīva, який плавав на цьому човні й мав із нього зиск, затишшя — катастрофа. Для океану — спокій. Одне й те саме явище описане двічі, і оцінки протилежні. Метафора подвійна: човен пливе **на** океані, тобто світ не є океаном, а лише тим, що по ньому рухається; і човен гине не від бурі, а від її відсутності. Це відповідає 15.18 («нема ні пут, ні визволення»): те, що зникає, не було справжнім, а лише плавало.}
\virvars{\virvaritem{SARIT і GRETIL (перевірено scripts/fetch\_witnesses.py)}{Читання збігаються; відмінності лише в аваґрасі та поділі слів. Не різночитання.}}
\virsources{Avg\_2.24}{verse\_2.24}

\index{jiva@\textit{jīva}}
\virs{2.25}{У мені, безмежному океані, диво: хвилі живих істот здіймаються, б'ються, граються й входять — самі собою, за власною природою.}
\virlit{«У мені, безмежному океані, — диво: хвилі живих істот здіймаються, б'ються, граються й входять — за власною природою.»}
\virgram{mayi anantamahāmbhodhau (див. 2.23; тут форма ambhodhau перед ā- наступного слова). āścaryam (n. nom. sg. «диво»; присудкове). jīvavīcayaḥ (m. nom. pl. «хвилі живих істот»: jīva + vīci «хвиля»). Чотири дієслова в теперішньому часі: udyanti («здіймаються»), ghnanti («б'ються, ударяють»), khelanti («граються»), praviśanti («входять»). svabhāvataḥ (прислівник «за власною природою, самі собою»).}
\vircom{Завершення триптиха, і воно знімає те, що було в 2.23–2.24. Там хвилі піднімав вітер, і човен гинув — тобто була причина й був той, кому зле. Тут **причини немає**: svabhāvataḥ («за власною природою») стоїть останнім словом і скасовує питання «хто це робить». Чотири дієслова йдуть без сполучників, підряд — це не опис послідовності, а перелік того, що відбувається водночас. ghnanti («б'ються») — єдине слово з відтінком болю, і воно стоїть у тому ж ряду, що khelanti («граються»), без переходу й без пояснення. Це і є āścaryam («диво»): не те, що все добре, а те, що все це відбувається саме, без господаря.}
\virvars{\virvaritem{SARIT і GRETIL (перевірено scripts/fetch\_witnesses.py)}{Читання збігаються; відмінності лише в аваґрасі та поділі слів. Не різночитання.}}
\virsources{Avg\_2.25}{verse\_2.25}

\pilotchapter{3}{14}
\index{katham@\textit{katham}}
\virs{3.1}{Пізнавши істинно Атман — незнищенний і єдиний — як у тебе, атмазнавця й мудрого, може бути втіха в набутті користі?}
\virlit{«Пізнавши істинно Атман — незнищенний, єдиний — як у тебе, атмазнавця, мудрого, можлива втіха в набутті користі?»}
\virgram{avināśinam (m. acc. sg. «незнищенний») ātmānam (m. acc. sg. «Атман») ekam (m. acc. sg. «єдиний») vijñāya (герунд vi-jñā «пізнавши») tattvataḥ (прислівник «істинно, по суті»). tava (gen. sg. «твого») ātmajñasya (m. gen. sg. «того, хто знає Атман») dhīrasya (m. gen. sg. «мудрого, стійкого»). katham («як?») arthārjane (n. loc. sg. «у набутті користі, в придбанні»: artha «користь» + arjana «набуття») ratiḥ (f. nom. sg. «втіха, прихильність»).}
\vircom{Питання katham («як?») риторичне й будується на несумісності двох занять: пізнання єдиного й накопичення множинного. Слово artha тут не «сенс», а саме «користь, вигода» — це той самий термін, що в 10.1 серед того, що відкидають. Зверни увагу на порядок означень: спершу те, що пізнано (незнищенне, єдине), потім той, хто пізнав (атмазнавець, мудрий). Втіха в набутті не заборонена — поставлено питання, чи вона взагалі можлива за такого пізнання.}
\virvars{\virvaritem{SARIT і GRETIL (перевірено scripts/fetch\_witnesses.py)}{ТУТ СПРАВЖНЄ РІЗНОЧИТАННЯ, не орфографія. SARIT подає один складений tavātmajñānadhīrasya (tava-ātma-jñāna-dhīra «твого знання-Атмана-мудрого»), GRETIL — два слова tavātmajñasya dhīrasya (tava-ātma-jñasya dhīrasya «твого, атмазнавця, мудрого»). Друге членування граматично чисте: ātmajña («той, хто знає Атман») — звичайний bahuvrīhi, а dhīra («мудрий») узгоджується з ним окремо. У SARIT доводиться припускати складене ātmajñāna («знання Атмана»), де dhīra стає третім членом — це можливо, але менш природно. Яке членування правильне, НЕ ВСТАНОВЛЕНО: третього свідка для цього вірша не перевірено. У перекладі взято членування GRETIL і це позначено тут, а не приховано.}}
\virsources{Avg\_3.1}{verse\_3.1}

\index{aho@\textit{aho}}\index{yatha@\textit{yathā}}
\virs{3.2}{О диво! Через незнання Атмана з'являється прихильність у полі омани чуттєвих предметів — як через незнання перламутру з'являється жадоба в омані срібла.}
\virlit{«О диво! Через незнання Атмана — прихильність у полі омани чуттєвих предметів; як через незнання перламутру — жадоба в омані срібла.»}
\virgram{ātmājñānāt (n. abl. sg. «через незнання Атмана») aho (вигук) prītiḥ (f. nom. sg. «прихильність, задоволення») viṣayabhramagocare (m. loc. sg. «у полі омани чуттєвих предметів»: viṣaya «чуттєвий предмет» + bhrama «омана» + gocara «поле»). śukteḥ (f. gen./abl. sg. «перламутру, устричної черепашки») ajñānataḥ (прислівник «через незнання»; тут причина, не джерело) lobhaḥ (m. nom. sg. «жадоба») yathā («як») rajatavibhrame (m. loc. sg. «в омані срібла»).}
\vircom{Вірш будується на точній паралелі з двох частин, і паралель тут не прикраса, а аргумент: prīti (прихильність) співвіднесена з lobha (жадобою), viṣayabhrama (омана чуттєвого) — з rajatavibhrama (омана срібла), ātmājñāna (незнання Атмана) — з śuktyajñāna (незнання перламутру). Причина в обох випадках одна — незнання, а не вада предмета: перламутр не винен, що його беруть за срібло. Це та сама паралель, що в 2.9, але там був перелік трьох прикладів, а тут — один, розгорнутий до кінця. Слово prīti м'якше за spṛhā (2.22) і за lobha: це радше «прихильність», ніж «спрага».}
\virvars{\virvaritem{SARIT і GRETIL (перевірено scripts/fetch\_witnesses.py)}{Читання збігаються; відмінності лише в аваґрасі та поділі слів. Не різночитання.}}
\virsources{Avg\_3.2}{verse\_3.2}


\virs{3.3}{Пізнавши: «я — те, у чому сяє весь цей світ, як хвилі в морі», — чого ти біжиш, наче вбогий?}
\virlit{«Пізнавши: „я — те, в чому сяє цей всесвіт, як хвилі в морі“, — чого ти біжиш, мов убогий?»}
\virgram{viśvam (n. nom. sg.) sphurati (3 sg. pr. sphur «сяє, тремтить, являється») yatra («де, у чому») idam (n. nom. sg.). taraṅgāḥ (m. nom. pl. «хвилі») iva sāgare (m. loc. sg. «у морі»). saḥ (m. nom. sg.) aham (nom. sg.) asmi («я є») iti («так»; частка прямої мови). vijñāya (герунд «пізнавши») kim («чого?») dīnaḥ (m. nom. sg. «убогий, жалюгідний») iva dhāvasi (2 sg. pr. dhāv «біжиш»).}
\vircom{Питання kim ... dhāvasi («чого біжиш») спрямоване не до дії взагалі, а до бігу як способу життя: дієслово dhāv означає саме фізичний біг, метушню. Порівняння з dīna («убогий») вказує на мотив, а не на стан: біжать ті, кому бракує. Формула so 'ham asmi («я є те») — та сама, що в 2.15, але там вона завершувала заперечення трійці, а тут відкриває питання. Образ хвиль у морі повторює 2.10 і 2.23–2.25: це наскрізний образ тексту, а не місцева метафора.}
\virvars{\virvaritem{SARIT і GRETIL (перевірено scripts/fetch\_witnesses.py)}{Читання збігаються; відмінності лише в аваґрасі та поділі слів. Не різночитання.}}
\virsources{Avg\_3.3}{verse\_3.3}


\virs{3.4}{Навіть почувши про Атман — чисту свідомість, надзвичайно прекрасну, — той, хто надмірно прив'язаний до того, що під рукою, набуває бруду.}
\virlit{«Навіть почувши [про] Атман — чисту свідомість, надзвичайно прекрасну, — надмірно прив'язаний до [того, що] під рукою, набуває бруду.»}
\virgram{śrutvā (герунд śru «почувши») api («навіть»). śuddhacaitanyam (m. acc. sg. «чиста свідомість»: śuddha «чистий» + caitanya «свідомість») ātmānam (m. acc. sg.) atisundaram (m. acc. sg. «надзвичайно прекрасний»). upasthe (n./m. loc. sg. «на лоні, під рукою, поблизу»; слово багатозначне) atyantasaṃsaktaḥ (m. nom. sg. «надмірно прив'язаний»: atyanta «надмірно» + saṃsakta «прив'язаний»). mālinyam (n. acc. sg. «бруд, нечистоту, занепад») adhigacchati (3 sg. pr. adhi-gam «набуває, доходить до»).}
\vircom{Слово upasthe навмисно не уточнене: воно означає і «на лоні», і «поблизу, під рукою», і традиційно вживається з плотським відтінком. Точний зміст залежить від контексту, якого вірш не дає, — і саме тому переклад зберігає двозначність, а не обирає одне значення. Структура вірша — контраст між почутим і робленим: śrutvā api («навіть почувши») підкреслює, що знання про Атман саме собою нічого не змінює. mālinya («бруд») — не моральна оцінка, а стан забрудненості; те саме слово вживане про одяг і про посуд. Тобто йдеться не про гріх, а про те, що річ втрачає чистоту від дотику.}
\virvars{\virvaritem{SARIT і GRETIL (перевірено scripts/fetch\_witnesses.py)}{Читання збігаються; відмінності лише в аваґрасі та поділі слів. Не різночитання.}}
\virsources{Avg\_3.4}{verse\_3.4}


\virs{3.5}{О диво: у мудреця, який знає Атман у всіх істотах і всі істоти — в Атмані, досі тримається «моє».}
\virlit{«О диво: у мудреця, який знає Атман у всіх істотах і всі істоти в Атмані, — [досі] тримається „моє“.»}
\virgram{sarvabhūteṣu (n. loc. pl. «у всіх істотах») ca ātmānam (m. acc. sg.) sarvabhūtāni (n. nom./acc. pl.) ca ātmani (m. loc. sg. «в Атмані»). muneḥ (m. gen. sg. «мудреця, відлюдника») jānataḥ (m. gen. sg. дієприкметника теперішнього часу jñā- «того, хто знає»). āścaryam (n. nom. sg. «диво»; присудкове) mamatvam (n. nom. sg. «„моє“ як стан, мійність») anuvartate (3 sg. pr. anu-vṛt «іде слідом, тримається»).}
\vircom{Двозмінна конструкція (Атман у всіх, усі в Атмані) — не повтор, а два напрямки одного знання: назовні й усередину. Диво в тому, що навіть за такого знання лишається mamatvam — не «моє» як майно, а мійність як звичка. Дієслово anuvartate («іде слідом») важливе: воно не каже, що мійність панує, а що вона тягнеться позаду. Це найточніше місце глави про межі знання: знання не вичищає звичку автоматично. Порівняй 3.4, де йшлося про того, хто лише почув: там бруд набувають, тут мійність лишається — різні ступені, і текст не згладжує різниці.}
\virvars{\virvaritem{SARIT і GRETIL (перевірено scripts/fetch\_witnesses.py)}{Читання збігаються; відмінності лише в аваґрасі та поділі слів. Не різночитання.}}
\virsources{Avg\_3.5}{verse\_3.5}

\index{bahuvrihi@\textit{bahuvrīhi}}
\virs{3.6}{Той, хто утвердився у вищій недвоїстості й стоїть навіть у справі звільнення, — о диво! — під владою бажання: безпорадний, мов у грі-вправлянні.}
\virlit{«Той, хто утвердився у вищій недвоїстості, хто стоїть навіть у справі звільнення, — о диво! — під владою бажання, безпорадний, наче [заради] вправи в грі.»}
\virgram{āsthitaḥ (m. nom. sg. «той, хто стоїть у, утвердився в») paramādvaitam (n. acc. sg. «вища недвоїстість»: parama «вищий» + advaita «недвоїстість»). mokṣārthe (m./n. loc. sg. «у справі звільнення»: mokṣa + artha «справа, мета») api («навіть»). vyavasthitaḥ (m. nom. sg. «утверджений, розміщений»). āścaryam (n. nom. sg. «диво»; присудкове) kāmavaśagaḥ (bahuvrīhi «той, хто пішов під владу бажання»: kāma «бажання» + vaśa «влада» + ga «той, що йде») vikalaḥ (m. nom. sg. «безпорадний, скалічений, позбавлений [сили]») keliśikṣayā (f. instr. sg. «як вправою-грою, задля забавки»: keli «гра, забава» + śikṣā «вправляння, наука»).}
\vircom{Вірш починає серію із чотирьох «див» (3.6–3.8, далі 3.10), і кожне будується за однією схемою: високий стан — і те, що в ньому несподівано лишається. Тут протиставлені два утвердження (āsthitaḥ, vyavasthitaḥ — майже синоніми, поставлені поруч навмисно) і один стан, що їх перекреслює. Слово vikala («безпорадний») сильніше за просте «слабкий»: воно означає позбавленість частини, скаліченість. keliśikṣayā («грою-вправлянням») знижує регістр: те, що мало бути шляхом, виявилося забавкою.}
\virvars{\virvaritem{SARIT і GRETIL (перевірено scripts/fetch\_witnesses.py)}{Читання збігаються; відмінності лише в аваґрасі та поділі слів. Не різночитання.}}
\virsources{Avg\_3.6}{verse\_3.6}

\index{jnana@\textit{jñāna}}
\virs{3.7}{О диво! Зрозумівши, що знання виявилося поганим другом, — надзвичайно слабкий — він усе одно прагнув би бажання, дійшовши до кінця часу.}
\virlit{«О диво! Зваживши, що постало погане знання-друг, — надзвичайно слабкий — він усе ж прагнув би бажання, дійшовши до кінця часу.»}
\virgram{udbhūtam (n. nom./acc. sg. «постале, яке виникло») jñānadurmitram (n. nom./acc. sg. «поганого друга — знання»: jñāna «знання» + durmitra «поганий друг»). avadhārya (герунд ava-dhṛ «зваживши, зрозумівши») atidurbalaḥ (m. nom. sg. «надзвичайно слабкий»). āścaryam (n. nom. sg. «диво»; присудкове) kāmam (m. acc. sg. «бажання») ākāṅkṣet (3 sg. opt. ā-kāṅkṣ «прагнув би»). kālam (m. acc. sg. «час») antam (m. acc. sg. «кінець») anuśritaḥ (m. nom. sg. «той, хто дійшов до, притулився до»: anu-śri).}
\vircom{Складене jñānadurmitram можна прочитати двома способами: «знання, що є поганим другом» і «поганий друг, [який полягає] у знанні». Обидва дають той самий сенс: знання, на яке покладалися, не врятувало. Найсильніше тут — час дієслова й обставина: ākāṅkṣet стоїть в optativus («прагнув би»), а kālam antam anuśritaḥ («дійшовши до кінця часу») означає останню межу, коли вже нічого не змінити. Тобто бажання не просто лишається — воно лишається там, де його вже неможливо задовольнити. Це не докір, а констатація абсурду, і слово āścaryam («диво») вжите саме в цьому сенсі.}
\virvars{\virvaritem{SARIT і GRETIL (перевірено scripts/fetch\_witnesses.py)}{Читання збігаються; відмінності лише в аваґрасі та поділі слів. Не різночитання.}}
\virsources{Avg\_3.7}{verse\_3.7}


\virs{3.8}{О диво: у того, хто відвернувся від цього й того світу, хто розрізняє вічне й невічне, хто бажає звільнення, — страх саме від звільнення.}
\virlit{«О диво: у того, хто відвернувся від цього й того [світу], хто розрізняє вічне й невічне, хто бажає звільнення, — страх саме від звільнення!»}
\virgram{iha-amutra (прислівники «тут» і «там»; разом «у цьому й у тому світі») viraktasya (m. gen. sg. «того, хто відвернувся, втратив прихильність»). nityānityavivekinaḥ (m. gen. sg. «того, хто розрізняє вічне й невічне»: nitya «вічне» + anitya «не вічне» + vivekin «той, хто розрізняє»). āścaryam (n. nom. sg. «диво»; присудкове) mokṣakāmasya (m. gen. sg. «того, хто бажає звільнення»). mokṣāt (m. abl. sg. «від звільнення») eva («саме») vibhīṣikā (f. nom. sg. «страх, жах; те, що лякає»).}
\vircom{Найгостріший вірш серії. Три ознаки поставлені поспіль і всі три «правильні»: відречення, розрізнення, бажання звільнення. І саме з третьої виростає страх, бо звільнення означає втрату того, хто звільняється. ablativus mokṣāt («від звільнення») з eva («саме») робить джерелом страху не світ і не смерть, а мету. vibhīṣikā — не «боязкість», а те, що лякає; слово вживане і про опудало. Тобто звільнення постає як те, чим лякають.}
\virvars{\virvaritem{SARIT і GRETIL (перевірено scripts/fetch\_witnesses.py)}{Читання збігаються; відмінності лише в аваґрасі та поділі слів. Не різночитання.}}
\virsources{Avg\_3.8}{verse\_3.8}


\virs{3.9}{Мудрий же — чи його голублять, чи його мучать — завжди, бачачи лише Атман, і не радіє, і не гнівається.}
\virlit{«Мудрий же — чи його годують, чи його мучать — завжди, бачачи лише самого Атман, не радіє й не гнівається.»}
\virgram{dhīraḥ (m. nom. sg. «мудрий, стійкий») tu («ж, але»). bhojyamānaḥ (m. nom. sg. пасивного дієприкметника теперішнього часу bhuj «той, кого годують, кому дають насолоду») api («хоч, і»). pīḍyamānaḥ (m. nom. sg. пасивного дієприкметника pīḍ «той, кого мучать, гнітять») api. sarvadā («завжди»). ātmānam (m. acc. sg.) kevalam (m. acc. sg. «лише, самого, єдиного») paśyan (m. nom. sg. дієприкметника теперішнього часу dṛś «бачачи»). na tuṣyati (3 sg. pr. tuṣ «не радіє, не задовольняється») na kupyati (3 sg. pr. kup «не гнівається»).}
\vircom{Два дієприкметники пасиву (bhojyamāna, pīḍyamāna) поставлені як рівноцінні умови: похвала й мука описані однією граматичною формою, тобто для мудрого вони вже не різні події. Це не байдужість: na tuṣyati na kupyati — два окремі заперечення, а не одне «не хвилюється». Слово kevalam («лише») робить ключову роботу: бачить не «Атман крізь усе», а лише Атман, і саме тому немає кому радіти. Форма дієприкметника paśyan (теперішній час) вказує на тривалість, а не на одноразове прозріння.}
\virvars{\virvaritem{SARIT і GRETIL (перевірено scripts/fetch\_witnesses.py)}{Читання збігаються; відмінності лише в аваґрасі та поділі слів. Не різночитання.}}
\virsources{Avg\_3.9}{verse\_3.9}

\index{katham@\textit{katham}}
\virs{3.10}{Своє тіло, що рухається, він бачить як тіло чуже. То як у хвалі чи в ганьбі збентежився б великодухий?}
\virlit{«Своє тіло, що ворушиться, він бачить як тіло чуже. У хвалі й у ганьбі — як міг би збентежитися великодухий?»}
\virgram{ceṣṭamānam (n. acc. sg. дієприкметника теперішнього часу ceṣṭ «той, що ворушиться, діє») śarīram (n. acc. sg. «тіло») svam (n. acc. sg. «своє») paśyati («бачить») anyaśarīravat (прислівник «як чуже тіло»: anya «інший» + śarīra + -vat). saṃstave (m. loc. sg. «у хвалі, у славослів'ї») ca api («і навіть») nindāyām (f. loc. sg. «у ганьбі»). katham («як?») kṣubhyet (3 sg. opt. kṣubh «збентежився б, захитався б») mahāśayaḥ (m. nom. sg. «великодухий, благородний»; буквально «той, у кого велике ложе/серце»: mahā + śaya).}
\vircom{Вірш закриває серію «див» і робить це тихо, без āścaryam: тут немає вигуку, є питання katham («як?»), на яке відповідь очевидна. ceṣṭamānam («те, що ворушиться») вибране навмисно: тіло описане не як річ, а як процес, і саме тому його легко побачити збоку. Пара saṃstava / nindā («хвала» / «ганьба») — та сама пара, що bhojyamāna / pīḍyamāna у 3.9, але там ішлося про дію на тіло, тут — про слово. mahāśayaḥ («великодухий») буквально означає «той, у кого велике ложе», і це не метафора розуму, а образ простору, у якому вміщається і хвала, і ганьба.}
\virvars{\virvaritem{SARIT і GRETIL (перевірено scripts/fetch\_witnesses.py)}{Розбіжність kṣubhyen (SARIT) / kṣubhyet (GRETIL) — це сандгі: t перед m переходить у n (kṣubhyet mahāśayaḥ → kṣubhyen mahāśayaḥ). Однакова форма, різне написання. Не різночитання.}}
\virsources{Avg\_3.10}{verse\_3.10}

\index{maya@\textit{māyā}}\index{bahuvrihi@\textit{bahuvrīhi}}\index{katham@\textit{katham}}\index{dhira@\textit{dhīra}}
\virs{3.11}{Бачачи, що весь цей світ — тільки мая, і втративши до нього охоту, — навіть коли смерть уже поруч, як злякався б той, чий розум стійкий?}
\virlit{«Бачачи, що весь цей світ — лише мая, [і] втративши цікавість, — навіть коли смерть уже прийшла, як злякався б той, чий розум стійкий?»}
\virgram{māyāmātram (n. nom. sg. «лише мая, не більше ніж мая»: māyā «мана, ілюзія» + mātra «лише») idam (n. nom. sg.) viśvam (n. nom. sg.). paśyan (m. nom. sg. дієприкметника теперішнього часу dṛś «бачачи») vigatakautukaḥ (bahuvrīhi «той, у кого минула цікавість»: vi-gata «минулий» + kautuka «цікавість, охота, святковість»). api («навіть») sannihite (m. loc. sg. дієприкметника sam-ni-dhā «присутній, близький») mṛtyau (m. loc. sg. «смерть»; locativus absolutus «коли смерть прийшла»). katham («як?») trasyati (3 sg. pr. tras «лякається, тремтить») dhīradhīḥ (bahuvrīhi «чий розум стійкий»: dhīra «стійкий» + dhī «розум»).}
\vircom{Слово kautuka тут важить більше, ніж здається: це не просто «цікавість», а охота до святкового, потяг до видовища. Тобто йдеться не про втрату знань, а про втрату інтересу до гри. locativus absolutus sannihite mṛtyau («коли смерть присутня») ставить умову максимально жорстко: не «якось колись», а тут і тепер. dhīradhīḥ («чий розум стійкий») — не «хоробрий»: стійкість розуму, а не безстрашність характеру. Питання katham знову риторичне, як у 3.3 і 3.10 — глава послідовно закриває питання питаннями.}
\virvars{\virvaritem{SARIT і GRETIL (перевірено scripts/fetch\_witnesses.py)}{Читання збігаються; відмінності лише в аваґрасі та поділі слів. Не різночитання.}}
\virsources{Avg\_3.11}{verse\_3.11}


\virs{3.12}{У кого розум без прагнень, хто навіть у безнадії лишається великим духом, хто вдоволений знанням Атмана, — з ким його й порівнювати?}
\virlit{«У кого розум без прагнень, хто навіть у безнадії великий духом, хто вдоволений знанням Атмана — з ким виникає порівняння [з ним]?»}
\virgram{niḥspṛham (n. nom. sg. «позбавлений прагнень»: niḥ «без» + spṛhā «спрага») mānasaṃ (n. nom. sg. «розум, мислення»; підмет підрядного) yasya (gen. sg.). nairāśye (n. loc. sg. «у безнадії, у відсутності сподівань»: nir-āśā «без сподівання») api («навіть»). mahātmanaḥ (m. gen. sg. «великого духом»). tasya (gen. sg.) ātmajñānatṛptasya (m. gen. sg. «вдоволеного знанням Атмана»: ātmajñāna + tṛpta «насичений, вдоволений»). tulanā (f. nom. sg. «порівняння, рівня, міра») kena (m. instr. sg. «з ким») jāyate (3 sg. pr. jan «виникає, народжується»).}
\vircom{nairāśye 'pi («навіть у безнадії») — найтонше місце. Безнадія тут не недолік, а умова: відсутність сподівань і є тим, що робить розум вільним від прагнень. Тому й сказано api («навіть») — ніби всупереч очікуванню, що без надії людина меншає. tulanā («порівняння, міра») вжите як річ, що не виникає: не «він не має рівних», а «порівняння не народжується», тобто зникає сама операція зіставлення. Форма jāyate (теперішній, «народжується») показує це як постійний стан, а не як разовий висновок.}
\virvars{\virvaritem{SARIT і GRETIL (перевірено scripts/fetch\_witnesses.py)}{Читання збігаються; відмінності лише в аваґрасі та поділі слів. Не різночитання.}}
\virsources{Avg\_3.12}{verse\_3.12}


\virs{3.13}{Знаючи самою природою, що все видиме — аніщо, — «це треба взяти, це треба відкинути» — то що ж бачить той, чий розум стійкий?}
\virlit{«Знаючи за самою природою, що це видиме — аніщо: „це слід узяти, це слід відкинути“ — що ж бачить той, чий розум стійкий?»}
\virgram{svabhāvāt (m. abl. sg. «за власною природою, самóю суттю») eva («саме»). jānānaḥ (m. nom. sg. дієприкметника теперішнього часу jñā «знаючи»). dṛśyam (n. nom. sg. «видиме») etat (n. nom. sg.) na kiñcana («ніщо, аніщо»). idam (n. nom. sg.) grāhyam (n. nom. sg. герундив від grah «те, що слід узяти») idam (n. nom. sg.) tyājyam (n. nom. sg. герундив від tyaj «те, що слід відкинути»). saḥ (m. nom. sg.) kim («що?») paśyati («бачить») dhīradhīḥ.}
\vircom{Дві пари герундивів (grāhyam / tyājyam) подані в лапках як чужа мова: це не позиція мудрого, а те, що він чує навколо або в собі. Заперечується не вибір між «взяти» й «відкинути», а сама наявність пари. svabhāvād eva («самою природою») знімає потребу в доказі: це не висновок, а спосіб бачення. Питання kim paśyati («що бачить») точніше за «чи бачить»: воно питає про предмет, а не про факт, бо коли видиме — аніщо, предмета не лишається.}
\virvars{\virvaritem{SARIT і GRETIL (перевірено scripts/fetch\_witnesses.py)}{Читання збігаються; відмінності лише в аваґрасі та поділі слів. Не різночитання.}}
\virsources{Avg\_3.13}{verse\_3.13}

\index{bahuvrihi@\textit{bahuvrīhi}}
\virs{3.14}{У того, в кого всередині відкинуто каламуть, хто без двоїстості й без сподівань, усе, що випадає само собою, не стає ні стражданням, ні задоволенням.}
\virlit{«У того, в кого всередині відкинуто каламуть, хто без двоїстості й без сподівань, — те, що випало, прийшовши само собою, не [йде] ні на страждання, ні на задоволення.»}
\virgram{antaḥ-tyakta-kaṣāyasya (bahuvrīhi «той, у кого всередині відкинуто каламуть»: antaḥ «усередині» + tyakta «відкинутий» + kaṣāya «каламуть, осад, пристрасть») nirdvandvasya (m. gen. sg. «той, хто без двоїстості»: nir «без» + dvandva «пара, двоїстість») nirāśiṣaḥ (m. gen. sg. «той, хто без сподівань, без побажань»: nir «без» + āśis «побажання, благословення»). yadṛcchayāgataḥ (m. nom. sg. «той, що прийшов сам собою»: yadṛcchā «власна воля, випадок» + āgata «прийшов») bhogaḥ (m. nom. sg. «досвід, насолода, те, що випадає»). na duḥkhāya (n. dat. sg. «на страждання») na tuṣṭaye (f. dat. sg. «на задоволення»).}
\vircom{Давальні duḥkhāya і tuṣṭaye — не «не приносить», а «не йде на»: досвід не зараховується ні туди, ні туди. Це тонше за байдужість: байдужість усе ще ставить оцінку «ніяк», а тут досвід просто не конвертується в жодну з двох категорій. kaṣāya («каламуть, осад») — слово з мови фарб і напоїв: це не гріх і не пристрасть у моральному сенсі, а те, що каламутить прозорість. antaḥ («усередині») стоїть першим членом навмисно: відкинуто не зовнішнє, а осад усередині. yadṛcchayāgataḥ («прийшов сам собою») — те саме yadṛcchā, що в 4.4 і далі; тут воно означає не випадковість, а відсутність замовлення.}
\virvars{\virvaritem{SARIT і GRETIL (перевірено scripts/fetch\_witnesses.py)}{Розбіжність yadṛcchayā+a+agato (SARIT) / yadṛcchayāgato (GRETIL) — це сандгі: подовжене ā на стику зникає у злитному написанні. Однакова форма. Не різночитання.}}
\virsources{Avg\_3.14}{verse\_3.14}

\pilotchapter{4}{6}
\index{samsara@\textit{saṃsāra}}
\virs{4.1}{О! Атмазнавцеві, мудрому, що грається грою досвідів, немає рівні з дурнями, які тягнуть сансару.}
\virlit{«О! Атмазнавця, мудрого, що грається грою досвідів, — немає рівності з дурнями, що тягнуть сансару.»}
\virgram{hanta (вигук «о!, ну ж!») ātmajñasya (m. gen. sg. «того, хто знає Атман») dhīrasya (m. gen. sg. «мудрого, стійкого»). khelataḥ (m. gen. sg. дієприкметника теперішнього часу khel «того, хто грається»). bhogalīlayā (f. instr. sg. «грою досвідів, грою в насолоди»: bhoga «досвід, насолода» + līlā «гра»). na hi («адже немає») saṃsāravāhīkaiḥ (m. instr. pl. «з тими, хто тягне сансару»: saṃsāra + vāhīka «той, хто везе, тягне, погонич») mūḍhaiḥ (m. instr. pl. «з дурнями, затьмареними») saha («разом з») samānatā (f. nom. sg. «рівність, однаковість»).}
\vircom{Слово vāhīka («погонич, той, хто везе») вибране точно: сансара тут не стан, у якому перебувають, а вантаж, який тягнуть. Дурість полягає не в незнанні, а в тому, що людина впряглася. Протиставлення будується на двох способах дії: khelataḥ («того, хто грається») і vāhīkaiḥ («тих, хто тягне»). Те саме діяння, різна якість. līlā («гра») — слово з релігійної мови, де гра не протистоїть серйозності, а означає дію без потреби. Вигук hanta на початку задає тон подиву, а не докору.}
\virvars{\virvaritem{SARIT verse\_4.1 і GRETIL/TextGrid Avg\_4.1}{Лексеми збігаються; відмінності — аваґраха, поділ слів і написання mūḍaiḥ (GRETIL/TextGrid) проти mūḍhaiḥ (SARIT). Аспірація ḍh — графічно-транслітераційна розбіжність, а не засвідчене рукописне різночитання; картка бере mūḍhaiḥ. Зіставлення GRETIL-боку зафіксоване в PR \#80.}\virvaritem{Атрибуція мовця перед главою 4: SARIT і GRETIL/TextGrid}{SARIT перед verse\_4.1 подає «janaka uvāca», GRETIL/TextGrid — «aṣṭāvakra uvāca». Це заголовковий матеріал електронних свідків, а не лексеми самого вірша; картка навмисно не приписує репліку жодному мовцеві. Перевірено лише sarit-івський бік; GRETIL-бік наведено за зіставленням у PR \#80.}}
\virsources{Avg\_4.1}{verse\_4.1}

\index{aho@\textit{aho}}
\virs{4.2}{О диво: те, до чого прагнуть убогі всі боги, починаючи з Шакри, — йог, що вже стоїть у цьому, не відчуває захоплення.}
\virlit{«О диво: стан, якого прагнуть убогі всі боги від Шакри, — йог, що стоїть у ньому, не відчуває захоплення.»}
\virgram{yat (n. nom./acc. sg. відносного займенника) padam (n. nom./acc. sg. «стопа, стан, місце»). prepsavaḥ (m. nom. pl. дезидератива pra-āp «ті, що прагнуть досягти») dīnāḥ (m. nom. pl. «убогі, жалюгідні») śakrādyāḥ (m. nom. pl. «починаючи з Шакри», тобто Індри) sarvadevatāḥ (f./m. nom. pl. «усі боги»). aho (вигук) tatra («там, у ньому») sthitaḥ (m. nom. sg. «той, хто стоїть») yogī (m. nom. sg. «йог») na harṣam (m. acc. sg. «радість, захоплення») upagacchati (3 sg. pr. upa-gam «доходить до, відчуває»).}
\vircom{Слово dīnāḥ («убогі») вжите про богів навмисно: убогість тут не в майні, а в прагненні — хто прагне, той має брак. Дезидератив prepsavaḥ («ті, що хочуть досягти») підкреслює саме бажання досягти, а не досягнення. Йог описаний дієприкметником sthitaḥ («той, хто стоїть») — це стан, а не подія; і саме тому harṣa («захоплення») не виникає: захоплення потребує новизни, а тут нічого нового не сталося. Вигук aho той самий, що в серії «див» глави 3, і тут він позначає межу між прагненням і перебуванням.}
\virvars{\virvaritem{SARIT і GRETIL (перевірено scripts/fetch\_witnesses.py)}{Читання збігаються; відмінності лише в аваґрасі та поділі слів. Не різночитання.}}
\virsources{Avg\_4.2}{verse\_4.2}


\virs{4.3}{Того, хто осягнув Істину, не зачіпають ані заслуги, ані провини. Так само й небо не зливається з димом, хоча на око вони ніби поєднані.}
\virlit{«У того, хто знає Те, не виникає внутрішнього дотику з заслугою та провиною — адже й у неба немає справжнього з'єднання з димом, навіть якщо здається, що таке з'єднання видно». Це авторський дослівно-смисловий виклад зі збереженням можливих неоднозначностей, а не затверджена редакція.}
\virgram{tajjñasya (tat-jña, родовий однини: «того, хто знає Те»; точний референт tat визначає контекст учення); puṇya-pāpābhyām (двійна форма іменників «заслуга» і «провина», орудний або відкладний за формою; у реченні зв'язана з «дотиком/впливом»); sparśaḥ → sparśo (називний однини: «дотик, зіткнення»); hi («адже»), antar («усередині»), na jāyate (заперечена форма третьої особи однини, ātmanepada, «не виникає»). Далі ākāśasya («неба, простору», родовий однини); dhūmena («димом», орудний однини); dṛśyamānā api («навіть коли видима/нібито помітна», дієприкметник жіночого роду з часткою); saṅgati / saṃgatiḥ («зв'язок, з'єднання», у контексті називний однини жіночого роду). SARIT має saṃgatiḥ, а TextGrid/GRETIL — saṅgati без кінцевого ḥ; морфологічний статус останньої форми потребує санскритологічної перевірки. Читання відновлено 2026-11-01 із SARIT (`verse\_4.3`, гілка master) і нормалізовано: у джерелі знак `+` позначає аваґраху, тому `dṛśyamānā+api` передано як `dṛśyamānā 'pi`. Поділ на слова звірено з джерелом; самостійної санскритологічної звірки за рукописом не проводилося.}
\vircom{У вірші 4.3 використано порівняння неба і диму, щоби розрізнити видимий зв'язок та сутнісне ототожнення. Зовнішнє спостереження диму на тлі неба ще не означає, що небо набуває властивостей диму. Аналогічно фігура «того, хто знає Те» в недвоїстому вченні не визначається заслугами чи провинами як власною сутністю. З цього не випливає юридичний або етичний дозвіл чинити шкоду: вірш формулює філософську тезу про природу свідка, а не правила поведінки. У робочому перекладі «осягнув Істину» — тлумачна експлікація займенника tat, можлива буквальніша редакція — «той, хто знає Те»; на рецензії слід обрати один варіант. Потрібно окремо звірити різницю «з'єднання» / «дотик» і мовну форму фінального іменника saṅgati / saṃgatiḥ. **Атрибуція мовця не усталена за звіреними електронними свідками:** GRETIL/TextGrid перед 4.1 подає «aṣṭāvakra uvāca», а SARIT в заголовку глави 4 — «janaka uvāca». Поки не залучено незалежне друковане видання чи рукописні свідчення, не слід беззастережно приписувати твердження Аштавакрі або Джанаці. Два електронні свідки не є критичним виданням; коментарі сучасних авторів тут не передруковано.}
\virvars{\virvaritem{GRETIL/TextGrid Avg\_4.3 та SARIT verse\_4.3}{GRETIL/TextGrid подає закінчення другого рядка як saṅgati (без кінцевого ḥ), SARIT — saṃgatiḥ (з ṃ та кінцевим вісарґа). Це реально засвідчена різниця двох електронних транскрипцій; без зіставлення друкованих видань і рукописного апарату не можна стверджувати, що це давнє рукописне різночитання. Особливо потрібно з'ясувати, чи відсутність ḥ у TextGrid є помилкою/редакційною конвенцією.}\virvaritem{GRETIL/TextGrid Avg\_4.3 та SARIT verse\_4.3}{GRETIL розділяє dṛśyamānāpi, а SARIT позначає відновлену голосну як dṛśyamānā+api. Те саме явище сандгі не є достатньою підставою для гіпотези про різні рукописи; до віршової бази цю різницю слід заносити як орфографічну/кодувальну.}\virvaritem{GRETIL/TextGrid — вступна репліка перед Avg\_4.1; SARIT TEI — <div n="4">, заголовок розділу}{ПЕРЕВІРЕНА СУПЕРЕЧНІСТЬ АТРИБУЦІЇ МОВЦЯ ВСІЄЇ ГЛАВИ: GRETIL/TextGrid безпосередньо перед 4.1 має «aṣṭāvakra uvāca», а SARIT у розмітці <div n="4"> має «janaka uvāca». Заголовки стоять перед 4.1, окремої зміни мовця всередині 4.1–4.3 не показано. Це конфлікт двох електронних джерел; не доведено рукописне різночитання і не встановлено авторитетнішу атрибуцію. На незалежній рецензії потрібно звірити походження заголовків з друкованими виданнями. Застереження стосується також сусідніх віршів 4.1, 4.2, 4.4–4.6, але жодні їхні картки тут не змінюються.}}
\virsources{Avg\_4.3}{verse\_4.3}

\index{jagat@\textit{jagat}}
\virs{4.4}{Хто здатен стримати того, хто пізнав увесь світ як єдиний Атман і живе так, як складаються обставини?}
\virlit{«Хто міг би стримати того великодушного, ким увесь цей світ пізнано як сам Атман, коли він живе так, як випадає?» Це власне смислове розгортання синтаксису, не буквальна копія стороннього перекладу; запитання має риторичний характер.}
\virgram{ātmaiva = ātmā + eva («сам Атман», із підсиленням eva); idaṃ jagat sarvaṃ — «увесь цей світ», середній рід, однина; jñātaṃ — дієприкметник минулого часу пасивного значення, середній рід однини, узгоджений із «світом». yena mahātmanā — обидві форми орудного однини («ким, яким великодушним/величним за духом»), що визначають того, ким це пізнано. yadṛcchayā — орудний способу/обставини: «як випадає, спонтанно, за обставинами», а не неодмінно «як заманеться». vartamānaṃ — дієприкметник теперішнього часу середнього стану від √vṛt, знахідний однини чоловічого роду, пов'язаний із taṃ («його, того»), а не обов'язково з jagat; TextGrid пише varttamānaṃ. niṣeddhuṃ — інфінітив «стримати, заборонити»; kṣameta — 3-тя особа однини потенційного способу, середній стан («зміг би»); kaḥ — питальний займенник «хто?». Граматичне віднесення vartamānaṃ до taṃ — робочий аналіз, який має окремо підтвердити санскритолог. Читання відновлено 2026-11-01 із SARIT (`verse\_4.4`, гілка master) і нормалізовано: у джерелі знак `+` позначає аваґраху; у цьому вірші такого знака немає, отже приклад із 4.3 тут не діє. Поділ на слова звірено з джерелом; самостійної санскритологічної звірки за рукописом не проводилося.}
\vircom{Вірш поєднує дві теми: пізнання недвоїстої природи світу та свободу від нав'язаної поведінкової схеми. Питання «хто здатен стримати?» є риторичним і характеризує свободу того, хто усвідомив світ як Атман. Слово yadṛcchayā не слід автоматично прирівнювати до примхливого бажання або безвідповідальності: у цьому контексті воно допускає читання «за збігом обставин, природним чином». Українське «єдиний Атман» — смислове відтворення підсилювального eva, а не окремий прикметник оригіналу; перед остаточним виданням перевірити, чи краще писати «сам Атман». Важливо, що vartamānaṃ taṃ граматично підтримує тлумачення про людину, яка живе спонтанно, а не тезу, ніби «світ робить усе, що йому хочеться». Цей коментар — редакторська інтерпретація, не незалежна наукова рецензія. Додатковий текстологічний блокер: GRETIL і SARIT не збігаються в атрибуції мовця всієї четвертої глави, тому цей вірш не приписано беззастережно жодному співрозмовникові.}
\virvars{\virvaritem{GRETIL/TextGrid Avg\_4.4 і SARIT verse\_4.4}{GRETIL/TextGrid має varttamānaṃ (подвоєне t), SARIT — vartamānaṃ (одне t). В обох контекст і морфологічна роль збігаються. Це підтверджена різниця написання двох електронних транскрипцій, але не доказ рукописного різночитання; потрібна незалежна звірка з друкованим виданням.}\virvaritem{GRETIL/TextGrid Avg\_4.4 і SARIT verse\_4.4}{GRETIL подає jagat sarvaṃ із пробілом; SARIT — jagatsarvaṃ. Ця відмінність сегментації не змінює переклад «увесь світ» і не є самостійним підтвердженим варіантом сенсу.}\virvaritem{Заголовок глави 4 у GRETIL/TextGrid та SARIT}{Перед Avg\_4.1 GRETIL/TextGrid подає aṣṭāvakra uvāca, а SARIT в заголовку <div n="4"> подає janaka uvāca. Окремої зміни мовця між 4.1 та 4.4 у цих електронних текстах не показано. Отже атрибуція голосу 4.4 між цифровими свідками суперечлива. Це ще не доведене рукописне різночитання; потрібна незалежна звірка історичних видань.}}
\virsources{Avg\_4.4}{verse\_4.4}


\virs{4.5}{Серед чотирьох родів істот — від Брахми до стеблини трави — лише той, хто пізнав, спроможний полишити і бажання, і небажання.}
\virlit{«У чотириродій сукупності істот, від Брахми аж до стеблини, спроможність полишити бажання й небажання належить, справді, саме тому, хто знає». Авторське дослівно-смислове розгортання іменної конструкції, не затверджений переклад.}
\virgram{ā-brahma-stamba-paryante — місцевий відмінок однини складного означення («аж до стеблини від Брахми»), синтаксично пов'язаний із bhūta-grāme; brahma тут в образі верхньої межі ряду може означати бога Брахму, а не автоматично безособове Брахман. bhūta-grāme — місцевий однини від bhūta-grāma («спільнота/сукупність істот»); catur-vidhe — місцевий однини, «чотириродій/чотирьох видів». Конкретний перелік чотирьох традиційних способів народження в цьому вірші не названий: без окремого джерела його не слід додавати як буквальний текст. vijñasya eva (в записі vijñasyaiva) — родовий однини «саме знавця / лише того, хто знає», з підсилювальною часткою eva; hi — «адже, справді»; sāmarthyam — середній рід, однина, «спроможність». icchā-anicchā-vivarjane — місцевий однини від vivarjana («відкинення, полишення») з парою icchā («бажання») / anicchā («небажання»; у цьому контексті допускає тлумачення «відраза»). У реченні немає вираженої особової дієслівної форми; український присудок «здатен» відновлено з іменного звороту. Питання, чи «відраза» достатньо точно передає anicchā, потребує рецензування. Читання відновлено 2026-11-01 із SARIT (`verse\_4.5`, гілка master) і нормалізовано: у джерелі знак `+` позначає аваґраху; у цьому вірші такий знак є, і його передано апострофом. Поділ на слова звірено з джерелом; самостійної санскритологічної звірки за рукописом не проводилося.}
\vircom{Вірш протиставляє всеосяжний образ світу істот — від найвищого міфологічного персонажа до найскромнішої рослини — стану знання. Ключовий вислів vijñasyaiva hi sāmarthyam указує не на соціальний ранг, а на спроможність «того, хто знає». Пара icchā / anicchā охоплює внутрішні схильності до прийняття й відкидання; переклад «бажання й небажання» ближчий до буквального заперечного словотвору, тоді як «бажання й відраза» є сильнішою філософською інтерпретацією, яку варто перевірити в незалежному коментарі. «Чотири роди» не слід подавати як сучасну біологічну систематику: вірш використовує традиційну космологічну категоризацію, але не перелічує її членів. У GRETIL/TextGrid перед 4.1 позначено «aṣṭāvakra uvāca», тоді як у SARIT заголовок четвертої глави має «janaka uvāca». Через невирішену суперечність заголовків коментар не приписує ці слова певному мовцеві. Це власне робоче пояснення за двома електронними транскрипціями, не критичне видання й не незалежна фахова рецензія.}
\virvars{\virvaritem{GRETIL/TextGrid Avg\_4.5 і SARIT verse\_4.5}{GRETIL/TextGrid: sāmarthyam icchānicchāvivarjane із пробілом і суцільним написанням icchānicchā; SARIT: sāmarthyamicchā+anicchāvivarjane із графічним розкриттям icchā+anicchā. Зміна межі слів та знака «+» у двох електронних наборах не доводить різного словникового читання чи рукописного різночитання; обидва передають протиставлення бажання/небажання. Потрібна перевірка друкованих джерел перед остаточним апаратом.}\virvaritem{GRETIL/TextGrid і SARIT: заголовок глави 4 перед віршем 4.1}{У GRETIL/TextGrid вказано aṣṭāvakra uvāca, у SARIT під заголовком 4-ї глави — janaka uvāca. Нова репліка між 4.1 та 4.5 не позначена. Тож атрибуція мовця для 4.5 за двома електронними свідками суперечлива; потрібне порівняння незалежних друкованих видань, а не самовільний вибір.}}
\virsources{Avg\_4.5}{verse\_4.5}

\index{jagat@\textit{jagat}}
\virs{4.6}{Хтось знає Атман — недвоїстий, володаря світу. Що він знає, те й робить, — і страху в нього немає ніде.}
\virlit{«Хтось знає Атман — недвоїстий, володаря світу. Що він знає, те й робить; страху в нього немає ніде.»}
\virgram{ātmānam (m. acc. sg.) advayam (m. acc. sg. «недвоїстий») kaścit («хтось, дехто») jānāti («знає»). jagat (n. nom./acc. sg. «світ») īśvaram (m. acc. sg. «володар»; разом jagadīśvara «володар світу», узгоджене з ātmānam). yad (n. acc. sg.) vetti («знає») tat (n. acc. sg.) saḥ (m. nom. sg.) kurute («робить»). na bhayam (n. nom. sg. «страх») tasya (m. gen. sg. «у нього») kutracit («будь-де, ніде»; з na — «ніде»).}
\vircom{kaścit («хтось») на початку — не випадкове слово: вірш не обіцяє, що це стосується кожного, і не називає його. Друга половина дає формулу yad vetti tat kurute («що знає, те й робить»), і вона сильніша за звичне «чини за знанням»: тут знання й дія не узгоджуються, а ототожнюються. kutracit з na («ніде») закриває простір повністю — страх не має місця, а не просто відсутній тут. Це останній вірш глави, і він повертає до її початку: 4.1 протиставляв того, хто грається, тому, хто тягне; тут сказано, чому — бо немає страху, а отже немає й потреби тягнути.}
\virvars{\virvaritem{SARIT і GRETIL (перевірено scripts/fetch\_witnesses.py)}{Читання збігаються; відмінності лише в аваґрасі та поділі слів. Не різночитання.}}
\virsources{Avg\_4.6}{verse\_4.6}

\pilotchapter{5}{4}

\virs{5.1}{У тебе немає зв'язку ні з чим — то чого ж ти, чистий, хочеш відкинути? Здійснюючи розчинення скупчення, так само йди до розчинення.}
\virlit{«У тебе немає зв'язку ні з чим — чого ти, чистий, хочеш відкинути? Здійснюючи розчинення скупчення, так само йди до розчинення.»}
\virgram{na te (gen. sg. «у тебе немає») saṅgaḥ (m. nom. sg. «зв'язок, дотик, прив'язаність») asti («є») kena api (m. instr. sg. «з чим-небудь, ні з чим»). kim («чого?») śuddhaḥ (m. nom. sg. «чистий») tyaktum (інфінітив tyaj «відкинути») icchasi (2 sg. pr. iṣ «хочеш»). saṅghātavilayam (m. acc. sg. «розчинення скупчення»: saṅghāta «скупчення, агрегат» + vilaya «розчинення») kurvan (m. nom. sg. дієприкметника теперішнього часу kṛ «роблячи») evam eva («так само») layam (m. acc. sg. «розчинення») vraja (2 sg. imper. vraj «іди»).}
\vircom{Питання kim ... icchasi («чого хочеш») викриває безглуздість самого наміру відректися: відкидають те, до чого прив'язані, а зв'язку немає (na te saṅgaḥ asti kenāpi). Слово saṅghāta («скупчення») — технічне: тіло описане не як річ, а як агрегат, суму. Тому vilaya («розчинення») і laya («розчинення») стоять в одному вірші двічі: перше — те, що відбувається з тілом само собою, друге — те, до чого слід іти. Різниця в тому, що перше трапляється, а друге обирають. РЕФРЕН ГЛАВИ. Усі чотири вірші глави 5 закінчуються однаково: evam eva layaṃ vraja («так само йди до розчинення»). Це не чотири окремі повчання, а одна літанія з незмінним приспівом: змінюється тільки перша половина — обґрунтування, а висновок лишається той самий. Тому главу слід читати як одне ціле.}
\virvars{\virvaritem{SARIT і GRETIL (перевірено scripts/fetch\_witnesses.py)}{Читання збігаються; відмінності лише в аваґрасі та поділі слів. Не різночитання.}}
\virsources{Avg\_5.1}{verse\_5.1}


\virs{5.2}{Світ постає з тебе, як бульбашка з океану. Знаючи, що Атман єдиний, — так само йди до розчинення.}
\virlit{«Світ постає з тебе, як бульбашка з океану. Знаючи, що Атман єдиний, — так само йди до розчинення.»}
\virgram{udeti (3 sg. pr. ud-i «постає, виходить угору») bhavataḥ (m. abl./gen. sg. «від тебе / твого») viśvam (n. nom. sg.). vāridheḥ (m. abl./gen. sg. «океану»: vāri «вода» + dhi «той, що містить») iva budbudaḥ (m. nom. sg. «бульбашка»). iti («так») jñātvā (герунд jñā «знаючи») ekam (m. acc. sg. «єдиний») ātmānam (m. acc. sg.). evam eva layam vraja.}
\vircom{Порівняння з бульбашкою точніше за звичне «хвиля в морі» (2.10, 2.23–2.25): хвиля є форма тієї ж води, а бульбашка — порожнеча, обтягнута плівкою з тієї ж води. Вона не має власного вмісту. bhavataḥ стоїть в ablativus/genitivus, і обидва прочитання працюють: «з тебе» (джерело) і «твого» (належність). Висновок вірша не «світ нереальний», а «Атман єдиний» — і саме з єдності випливає розчинення, а не навпаки. РЕФРЕН ГЛАВИ. Усі чотири вірші глави 5 закінчуються однаково: evam eva layaṃ vraja («так само йди до розчинення»). Це не чотири окремі повчання, а одна літанія з незмінним приспівом: змінюється тільки перша половина — обґрунтування, а висновок лишається той самий. Тому главу слід читати як одне ціле.}
\virvars{\virvaritem{SARIT verse\_5.2 і GRETIL/TextGrid Avg\_5.2}{GRETIL/TextGrid подає jñatvaikam (короткий a), SARIT — jñātvaikam (довгий ā): це різниця кількості голосного в jñātvā, а не поділу слів. Деванаґарі в Sanskrit Documents (ज्ञात्वैकमात्मानम्) підтримує довгий ā, тож короткий варіант слід вважати можливим дефектом електронної транслітерації, а не рукописним різночитанням. Картка бере jñātvā ekam. Зіставлення GRETIL-боку зафіксоване в PR \#99.}}
\virsources{Avg\_5.2}{verse\_5.2}


\virs{5.3}{Хоч світ і перед очима, він не існує в тобі, незаплямованому, — бо він не річ, як змій у мотузці, явно. Так само йди до розчинення.}
\virlit{«Хоч світ і перед очима, він не існує в тобі, незаплямованому, — через несуттєвість, як змій у мотузці, явно. Так само йди до розчинення.»}
\virgram{pratyakṣam (n. nom./acc. sg. «очевидне, те, що перед очима») api («хоч»). avastutvāt (n. abl. sg. «через несуттєвість, через відсутність речі»: a-vastu «не-річ» + -tva). viśvam (n. nom. sg.) na asti («не існує») amale (m. loc. sg. «у незаплямованому») tvayi (loc. sg. «у тобі»). rajjusarpaḥ (m. nom. sg. «змій у мотузці»: rajju «мотузка» + sarpa «змій») iva vyaktam (n. nom./acc. sg. «явно, чітко»). evam eva layam vraja.}
\vircom{Найсильніше тут — avastutvāt («через несуттєвість»). Це не «світ оманливий», а точніше: світ не є vastu, річчю, самостійною сутністю. pratyakṣam api («хоч і перед очима») визнає видимість повністю — вірш не заперечує, що ми бачимо. Заперечується лише те, що видиме має власне буття. Змій у мотузці (той самий приклад, що в 2.9) поданий тут з vyaktam («явно»): помилка не тьмяна, вона гранично чітка — і саме тому переконлива. РЕФРЕН ГЛАВИ. Усі чотири вірші глави 5 закінчуються однаково: evam eva layaṃ vraja («так само йди до розчинення»). Це не чотири окремі повчання, а одна літанія з незмінним приспівом: змінюється тільки перша половина — обґрунтування, а висновок лишається той самий. Тому главу слід читати як одне ціле.}
\virvars{\virvaritem{SARIT і GRETIL (перевірено scripts/fetch\_witnesses.py)}{Читання збігаються; відмінності лише в аваґрасі та поділі слів. Не різночитання.}}
\virsources{Avg\_5.3}{verse\_5.3}

\index{bahuvrihi@\textit{bahuvrīhi}}\index{sukha@\textit{sukha}}\index{duhkha@\textit{duḥkha}}
\virs{5.4}{Рівний у стражданні й радості, повний, рівний у сподіванні й безнадії, рівний у житті й смерті, — так само йди до розчинення.}
\virlit{«Рівний у стражданні й радості, повний, рівний у сподіванні й безнадії, рівний у житті й смерті — бувши [таким], так само йди до розчинення.»}
\virgram{samaduḥkhasukhaḥ (bahuvrīhi «рівний у стражданні й радості»: sama + duḥkha + sukha) pūrṇaḥ (m. nom. sg. «повний»). āśānairāśyayoḥ (f. loc. du. «у сподіванні й безнадії»: āśā «сподівання» + nairāśya «безнадія») samaḥ (m. nom. sg. «рівний»). samajīvitamṛtyuḥ (bahuvrīhi «рівний у житті й смерті»: sama + jīvita + mṛtyu) san (m. nom. sg. дієприкметника теперішнього часу as- «бувши») evam eva layam vraja.}
\vircom{Вірш складається з трьох пар, і всі три побудовані на sama («рівний»): страждання/радість, сподівання/безнадія, життя/смерть. Пари охоплюють три виміри — відчуття, очікування, буття. Слово pūrṇaḥ («повний») стоїть між першою і другою парою не випадково: рівність не порожнеча, а повнота, якій не бракує нічого з обох боків. Це єдиний вірш глави, де немає ні прикладу, ні заперечення — лише перелік ознак. Тим самим літанія закінчується не аргументом, а описом того, хто йде. РЕФРЕН ГЛАВИ. Усі чотири вірші глави 5 закінчуються однаково: evam eva layaṃ vraja («так само йди до розчинення»). Це не чотири окремі повчання, а одна літанія з незмінним приспівом: змінюється тільки перша половина — обґрунтування, а висновок лишається той самий. Тому главу слід читати як одне ціле.}
\virvars{\virvaritem{SARIT і GRETIL (перевірено scripts/fetch\_witnesses.py)}{Читання збігаються; відмінності лише в аваґрасі та поділі слів. Не різночитання.}}
\virsources{Avg\_5.4}{verse\_5.4}

\pilotchapter{6}{4}
\index{jagat@\textit{jagat}}\index{tatha@\textit{tathā}}
\virs{6.1}{Я безмежний, як простір, а світ — як горщик. Таке знання — і в нього немає ні відкидання, ні взяття, ні розчинення.}
\virlit{«Я безмежний, як простір; світ — як горщик. Таке знання — і в нього немає ні відкидання, ні взяття, ні розчинення.»}
\virgram{ākāśavat (прислівник «як простір, як ефір»: ākāśa + -vat) anantaḥ (m. nom. sg. «безмежний») aham (nom. sg.). ghaṭavat (прислівник «як горщик»: ghaṭa «горщик» + -vat) prakṛtam (n. nom. sg. «зроблене, створене»; від pra-kṛ) jagat (n. nom. sg. «світ»). iti jñānam (n. nom. sg. «таке знання»). tathā («так, відповідно») etasya (n. gen. sg. «у нього, у цього [знання]») na tyāgaḥ (m. nom. sg. «немає відкидання») na grahaḥ (m. nom. sg. «немає взяття») layaḥ (m. nom. sg. «[немає] розчинення»).}
\vircom{Образ простору й горщика працює точніше, ніж здається: простір не вміщується в горщик і не обмежується його стінками — горщик лише окреслює частину простору, не додаючи й не забираючи нічого. Тому рефрен і каже: немає ні tyāga (відкидання), ні graha (взяття), ні навіть laya (розчинення) — усі три дії припускають, що щось справді додалося або зникло. Слово prakṛtam («зроблене») підкреслює: горщик — виріб, а простір — ні. РЕФРЕН ГЛАВИ. Усі чотири вірші глави 6 закінчуються однаково: iti jñānaṃ tathaitasya na tyāgo na graho layaḥ («таке знання — і в нього немає ні відкидання, ні взяття, ні розчинення»). Зверни увагу: рефрен глави 5 вимагав laya («йди до розчинення»), а рефрен глави 6 заперечує навіть laya. Це не суперечність, а наступний крок: спершу розчинення як те, до чого йдуть, потім — знання, у якому немає кому йти й нічого розчиняти. Главу слід читати як одне ціле.}
\virvars{\virvaritem{SARIT verse\_6.1 і GRETIL/TextGrid Avg\_6.1}{GRETIL/TextGrid подає ghaṭavat prakṛtaṃ jagat, SARIT — ghaṭavatprākṛtaṃ jagat: різниця довготи ā, а не поділу слів, і вона може бути лексико-семантичною (prakṛta «природний, звичайний» / prākṛta «природний, первісний»). Картка бере prakṛtaṃ; питання лишається відкритим до друкованого свідка. Зіставлення GRETIL-боку зафіксоване в PR \#91.}\virvaritem{SARIT verse\_6.1}{SARIT записує ananto+ahaṃ; знак «+» передає відновлений голосний на межі сандгі. Це графіка електронного запису, а не різночитання.}}
\virsources{Avg\_6.1}{verse\_6.1}

\index{tatha@\textit{tathā}}
\virs{6.2}{Я — як великий океан, а цей плин світу — як хвилі. Таке знання — і в нього немає ні відкидання, ні взяття, ні розчинення.}
\virlit{«Я — як великий океан; цей плин світу — як хвилі. Таке знання — і в нього немає ні відкидання, ні взяття, ні розчинення.»}
\virgram{mahodadhiḥ (m. nom. sg. «великий океан»: mahā «великий» + udadhi «океан») iva aham («я — як»). saḥ (m. nom. sg.) prapañcaḥ (m. nom. sg. «плин, розгортання світу, різноманіття»: pra-pañca) vīcisaṃnibhaḥ (m. nom. sg. «подібний до хвиль»: vīci «хвиля» + saṃnibha «подібний»). iti jñānaṃ tathā etasya na tyāgaḥ na grahaḥ layaḥ.}
\vircom{prapañca — важке слово: буквально «розгортання, розпросторення», а в ужитку «різноманіття світу» з відтінком заплутаності. Тут воно вжите без засудження: хвилі не є вадою океану. Порівняй 2.23–2.25, де той самий образ розгорнутий у три кроки; тут він стиснутий до одного порівняння, бо мета інша — не показати, як хвилі виникають, а показати, що з ними не треба нічого робити. РЕФРЕН ГЛАВИ. Усі чотири вірші глави 6 закінчуються однаково: iti jñānaṃ tathaitasya na tyāgo na graho layaḥ («таке знання — і в нього немає ні відкидання, ні взяття, ні розчинення»). Зверни увагу: рефрен глави 5 вимагав laya («йди до розчинення»), а рефрен глави 6 заперечує навіть laya. Це не суперечність, а наступний крок: спершу розчинення як те, до чого йдуть, потім — знання, у якому немає кому йти й нічого розчиняти. Главу слід читати як одне ціле.}
\virvars{\virvaritem{SARIT і GRETIL (перевірено scripts/fetch\_witnesses.py)}{Читання збігаються; відмінності лише в аваґрасі та поділі слів. Не різночитання.}}
\virsources{Avg\_6.2}{verse\_6.2}

\index{visva@\textit{viśva}}\index{tatha@\textit{tathā}}
\virs{6.3}{Я — як перламутр, а уявлення світу — як срібло в ньому. Таке знання — і в нього немає ні відкидання, ні взяття, ні розчинення.}
\virlit{«Я — як перламутр, а уявлення світу — як срібло. Таке знання — і в нього немає ні відкидання, ні взяття, ні розчинення.»}
\virgram{aham (nom. sg.) saḥ (m. nom. sg. «той, те») śuktisaṅkāśaḥ (m. nom. sg. «подібний до перламутру, що має вигляд перламутру»: śukti «перламутр, устрична черепашка» + saṅkāśa «вигляд, подоба»). rūpyavat (прислівник «як срібло»: rūpya + -vat) viśvakalpanā (f. nom. sg. «уявлення світу, вигадування світу»: viśva «світ» + kalpanā «уявлення, конструкція»). iti jñānaṃ tathā etasya na tyāgaḥ na grahaḥ layaḥ.}
\vircom{Третє й останнє з трьох порівнянь глави (простір/горщик, океан/хвилі, перламутр/срібло), і воно єдине, де друга річ не існує зовсім: горщик і хвиля є, срібла немає. Тому тут і стоїть kalpanā («уявлення, конструкція») — не «світ», а саме уявлення світу. Зверни увагу: перламутр не заперечується і не змінюється — зникає тільки срібло, а разом із ним і питання, що з ним робити. Це і є сенс рефрену. РЕФРЕН ГЛАВИ. Усі чотири вірші глави 6 закінчуються однаково: iti jñānaṃ tathaitasya na tyāgo na graho layaḥ («таке знання — і в нього немає ні відкидання, ні взяття, ні розчинення»). Зверни увагу: рефрен глави 5 вимагав laya («йди до розчинення»), а рефрен глави 6 заперечує навіть laya. Це не суперечність, а наступний крок: спершу розчинення як те, до чого йдуть, потім — знання, у якому немає кому йти й нічого розчиняти. Главу слід читати як одне ціле.}
\virvars{\virvaritem{SARIT і GRETIL (перевірено scripts/fetch\_witnesses.py)}{Розбіжність śuktisaṃkāśo (SARIT) / śuktisaṅkāśo (GRETIL) — анусвара проти класового носового перед k. Однакова форма (saṃkāśa / saṅkāśa «подоба, вигляд»), різне написання. Орфографія, не різночитання.}}
\virsources{Avg\_6.3}{verse\_6.3}

\index{tatha@\textit{tathā}}
\virs{6.4}{Я — у всіх істотах, а всі істоти — у мені. Таке знання — і в нього немає ні відкидання, ні взяття, ні розчинення.}
\virlit{«Я — у всіх істотах, а всі істоти — у мені. Таке знання — і в нього немає ні відкидання, ні взяття, ні розчинення.»}
\virgram{aham (nom. sg.) vā («або ж; також») sarvabhūteṣu (n. loc. pl. «у всіх істотах»). sarvabhūtāni (n. nom. pl.) atho («а також, потім») mayi (loc. sg. «у мені»). iti jñānaṃ tathā etasya na tyāgaḥ na grahaḥ layaḥ.}
\vircom{Двозмінна формула та сама, що в 3.5, і тут вона завершує главу. Після трьох порівнянь (простір, океан, перламутр) висновок поданий без образу взагалі — лише відношення. atho («а також») з'єднує два напрямки так, що вони не зливаються в один: я в істотах і істоти в мені — не те саме твердження, прочитане двічі, а два боки одного відношення. Саме тому рефрен знову заперечує tyāga і graha: у такому відношенні немає нічого, що можна було б узяти або відкинути. РЕФРЕН ГЛАВИ. Усі чотири вірші глави 6 закінчуються однаково: iti jñānaṃ tathaitasya na tyāgo na graho layaḥ («таке знання — і в нього немає ні відкидання, ні взяття, ні розчинення»). Зверни увагу: рефрен глави 5 вимагав laya («йди до розчинення»), а рефрен глави 6 заперечує навіть laya. Це не суперечність, а наступний крок: спершу розчинення як те, до чого йдуть, потім — знання, у якому немає кому йти й нічого розчиняти. Главу слід читати як одне ціле.}
\virvars{\virvaritem{SARIT і GRETIL (перевірено scripts/fetch\_witnesses.py)}{Читання збігаються; відмінності лише в аваґрасі та поділі слів. Не різночитання.}}
\virsources{Avg\_6.4}{verse\_6.4}

\pilotchapter{7}{5}
\index{visva@\textit{viśva}}
\virs{7.1}{У мені, безмежному океані, човен світу блукає туди-сюди вітром власного нутра — і в мене немає нетерпимості.}
\virlit{«У мені, безмежному океані, човен світу блукає туди-сюди вітром власного нутра — і в мене немає нетерпимості.»}
\virgram{mayi ananta-mahā-ambhodhau (m. loc. «у безмежному великому океані — у мені») | viśva-potaḥ (m. «човен світу», pota = човен) itastataḥ («туди-сюди») bhramati («блукає») svānta-vātena (instr. «вітром власного нутра») | na mama asti asahiṣṇutā («немає в мене нетерпимості»).}
\vircom{`pota` — «човен, судно», а не «хвиля»: світ тут не хвилювання океану, а рухома річ **на** ньому. Це відрізняє 7.1 від 7.2, де світ уже `vīci` (хвиля). `svānta-vāta` («вітер власного нутра») — рушій усередині, не зовні. `asahiṣṇutā` — «нетерпимість, нездатність терпіти»: океан не страждає від човна, і саме це сказано, а не «спокій».}
\virvarnone{власних приміток про різночитання немає; звіряння свідків не позначене як виконане}
\virsources{Avg\_7.1}{verse\_7.1}

\index{jagat@\textit{jagat}}
\virs{7.2}{У мені, безмежному океані, хвиля світу — за власною природою. Нехай піднімається або спадає: у мене немає ні приросту, ні втрати.}
\virlit{«У мені, безмежному океані, хвиля світу — за власною природою. Нехай піднімається або спадає: у мене немає ні приросту, ні втрати.»}
\virgram{mayi ananta-mahā-ambhodhau | jagat-vīciḥ (f. «хвиля світу») svabhāvataḥ («за власною природою») | udetu vā astam āyātu (імперативи «нехай піднімається або нехай спадає») | na me vṛddhiḥ na ca kṣatiḥ («немає в мене ні приросту, ні втрати»).}
\vircom{Два імперативи (`udetu`, `āyātu`) — це не побажання, а **дозвіл**: текст не просить хвилю поводитися інакше. `svabhāvataḥ` («за власною природою») знімає питання причини. Пара `vṛddhi` / `kṣati` — приріст і втрата: океан не стає більшим від хвилі й не меншим без неї. Разом із 7.1 це друга метафора того самого, і різниця між човном і хвилею не випадкова.}
\virvarnone{власних приміток про різночитання немає; звіряння свідків не позначене як виконане}
\virsources{Avg\_7.2}{verse\_7.2}


\virs{7.3}{У мені, безмежному океані те, що зветься світом, є домислюванням. Вельми втихлий, без форми — саме в цьому я стою.}
\virlit{«У мені, безмежному океані, те, що зветься світом, є домислюванням. Вельми втихлий, без форми — саме в цьому я стою.»}
\virgram{mayi ananta-mahā-ambhodhau | viśvam nāma («те, що зветься світом») vikalpanā (f. «домислювання, конструкція») | atiśāntaḥ («вельми втихлий») nirākāraḥ («без форми») etat eva aham āsthitaḥ («саме в цьому я стою»).}
\vircom{`nāma` тут не «ім'я», а знецінювальна частка: «те, що **зветься** світом». `vikalpanā` — та сама «конструкція розуму», що в 2.9 і 2.17; словник карток має лишатися узгодженим. Формула `etad evāham āsthitaḥ` («саме в цьому я стою») повторюється в 7.3 і 7.4 — це не прикраса, а спосіб зчепити два вірші в одне твердження.}
\virvarnone{власних приміток про різночитання немає; звіряння свідків не позначене як виконане}
\virsources{Avg\_7.3}{verse\_7.3}

\index{atman@\textit{ātman}}\index{niranjana@\textit{nirañjana}}
\virs{7.4}{Я не в станах буття, і стан — не там, у безмежному, незаплямленому. Так — неприв'язаний, без спраги, втихлий — саме в цьому я стою.}
\virlit{«Я не в станах буття, і стан — не там, у безмежному, незаплямленому. Так — неприв'язаний, без спраги, втихлий — саме в цьому я стою.»}
\virgram{na ātmā bhāveṣu («Я не в станах [буття]») na u bhāvaḥ («і не стан — [те, що] там») tatra anante nirañjane («там, у безмежному, незаплямленому») | iti («так») asaktaḥ («неприв'язаний») aspṛhaḥ («без спраги») śāntaḥ («втихлий») etat eva aham āsthitaḥ.}
\vircom{Вірш навмисно важкий: заперечуються обидві сторони — і `ātman` як один зі `bhāva`, і `bhāva` як те, що перебуває в безмежному. Це не заклик, а **усунення альтернативи**, і спрощувати його в перекладі не можна. `nirañjana` — той самий епітет, що в 1.15 і 2.15: незаплямленість тут не метафора чистоти, а вказівка на відсутність того, до чого можна приліпитися.}
\virvarnone{власних приміток про різночитання немає; звіряння свідків не позначене як виконане}
\virsources{Avg\_7.4}{verse\_7.4}

\index{jagat@\textit{jagat}}\index{aho@\textit{aho}}\index{katham@\textit{katham}}
\virs{7.5}{О диво: я — саме лише свідомість, а світ подібний до чару. Тож звідки й де в мені домислювання про те, що відкинути і що прийняти?}
\virlit{«О диво: я — саме лише свідомість, а світ подібний до чару. Тож звідки й де в мені домислювання про те, що відкинути і що прийняти?»}
\virgram{aho cinmātram eva aham («я — саме лише свідомість») | indrajāla-upamam (n. «подібне до чару») jagat | ataḥ mama katham kutra heya-upādeya-kalpanā («звідки в мене, де — домислювання про те, що відкинути й що прийняти»).}
\vircom{`cinmātra` («саме лише свідомість») — ключове слово розділу; переклад не повинен додавати «чиста», бо `mātra` означає «не більше ніж». `indrajāla` — той самий образ чару, що в 10.2, і словник має бути спільним. Питання `kathaṃ kutra` подвійне — «як» і «де» — і саме подвійність робить його риторичним: для вибору не лишається ні способу, ні місця. Глава 7 на цьому закінчується.}
\virvarnone{власних приміток про різночитання немає; звіряння свідків не позначене як виконане}
\virsources{Avg\_7.5}{verse\_7.5}

\pilotchapter{8}{4}

\virs{8.1}{Полон — тоді, коли розум чогось бажає, сумує, щось відпускає або бере, чомусь радіє чи гнівається.}
\virlit{«Полон — тоді, коли розум чогось бажає, сумує, щось відпускає, бере, чомусь радіє, гнівається.»}
\virgram{tadā bandhaḥ yadā (кореляція «тоді… коли») cittam (n. «розум») kiñcit vāñchati («чогось бажає») śocati («сумує») | kiñcit muñcati («щось відпускає») gṛhṇāti («бере») kiñcit hṛṣyati («чомусь радіє») kupyati («гнівається»).}
\vircom{Визначення полону тут не метафізичне, а **описове**: шість дієслів, і кожне — звичайний рух розуму. Зверніть увагу, що «бере» і «відпускає» стоять в одному ряду з гнівом: вибір сам по собі є ознакою полону. Перелік не оцінює, а перелічує — і саме тому діє.}
\virvars{\virvaritem{SARIT verse\_8.1 і verse\_8.2}{У 8.1 обидва електронні свідки мають dhṛṣyati (SARIT: kiñciddhṛṣyati), тоді як hṛṣyati стоїть у 8.2 (na hṛṣyati). Поля grammar і literal\_meaning передають dhṛṣyati як «чомусь радіє», тобто за значенням hṛṣyati; це робоче перекладацьке рішення, а не засвідчене читання, і його не слід виправляти механічно за паралельним рядком 8.2. Значення dhṛṣyati лишається відкритим до друкованого свідка.}\virvaritem{SARIT verse\_8.1 і GRETIL/TextGrid Avg\_8.1}{ṣocati (GRETIL/TextGrid) проти śocati (SARIT) і kiṃcid/kiṃcin (GRETIL/TextGrid) проти kiñcid/kiñcin (SARIT) — графіка й сандгі-запис, а не встановлене рукописне різночитання. GRETIL у цій картці не звірявся (sources.gretil.checked\_on: null), тож sarit-івський бік перевірено, а GRETIL-бік наведено за зіставленням у PR \#106.}}
\virsources{Avg\_8.1}{verse\_8.1}


\virs{8.2}{Визволення — тоді, коли розум не бажає, не сумує, не відпускає й не бере, не радіє й не гнівається.}
\virlit{«Визволення — тоді, коли розум не бажає, не сумує, не відпускає, не бере, не радіє, не гнівається.»}
\virgram{tadā muktiḥ yadā cittaṃ na vāñchati na śocati | na muñcati na gṛhṇāti na hṛṣyati na kupyati — шість заперечень, дослівно ті самі дієслова, що в 8.1.}
\vircom{Вірш є **дослівним запереченням** попереднього: ті самі шість дієслів, кожне з часткою na. Це важливо для перекладу: будь-яка стилістична варіація тут зруйнує паралелізм, який і є аргументом. Визволення не додає нічого — воно знімає рівно те, що було названо полоном.}
\virvarnone{власних приміток про різночитання немає; звіряння свідків не позначене як виконане}
\virsources{Avg\_8.2}{verse\_8.2}


\virs{8.3}{Полон — тоді, коли розум прив'язаний до якихось поглядів; визволення — коли розум неприв'язаний до всіх поглядів.}
\virlit{«Полон — тоді, коли розум прив'язаний до якихось поглядів; визволення — коли розум неприв'язаний до всіх поглядів.»}
\virgram{tadā bandhaḥ yadā cittam saktaṃ (n. «прив'язаний») kāsu api dṛṣṭiṣu (f. loc. pl. «у якихось поглядах») | tadā mokṣaḥ yadā cittam asaktam (n. «неприв'язаний») sarva-dṛṣṭiṣu (loc. pl. «у всіх поглядах»).}
\vircom{`dṛṣṭi` — «погляд» у подвійному сенсі: і зір, і переконання. Асиметрія навмисна: `kāsvapi` («до якихось») проти `sarva` («до всіх»). Тобто достатньо **однієї** прив'язаності, щоб бути в полоні, але потрібна неприв'язаність до **всіх** — і це не той самий обсяг вимог, що звучить суворіше, ніж здається з першого читання.}
\virvarnone{власних приміток про різночитання немає; звіряння свідків не позначене як виконане}
\virsources{Avg\_8.3}{verse\_8.3}


\virs{8.4}{Коли немає «я» — визволення; коли є «я» — полон. Так розуміючи, легко нічого не бери й нічого не відпускай.}
\virlit{«Коли немає „я“ — визволення; коли є „я“ — полон. Так розуміючи, легко нічого не бери й нічого не відпускай.»}
\virgram{yadā na aham («коли немає „я“») tadā mokṣaḥ | yadā aham («коли є „я“») bandhanam tadā | matvā iti (герундив «так розуміючи») helayā («легко, без натуги») kiñcit mā gṛhāṇa («нічого не бери») vimuñca mā («нічого не відпускай»).}
\vircom{Тут уперше названо причину, а не симптом: попередні три вірші описували рухи розуму, цей називає їхнє джерело — `aham`. `helayā` («легко, жартома») різко змінює тон: те, що виглядало як велика вимога, подане як щось майже незначне. І `mā … mā` — подвійна заборона: не бери й не відпускай, тобто не роби й самого жесту вибору.}
\virvarnone{власних приміток про різночитання немає; звіряння свідків не позначене як виконане}
\virsources{Avg\_8.4}{verse\_8.4}

\pilotchapter{9}{8}

\virs{9.1}{Зроблене й незроблене, пари протилежностей — коли вони втихли і чиї вони? Так знаючи, будь тут, у відреченні, схильним до відкидання — і без обітниць.}
\virlit{«Зроблене й незроблене, пари протилежностей — коли вони втихли і чиї вони? Так знаючи, будь тут, у відреченні, схильним до відкидання — і без обітниць.»}
\virgram{kṛta-akṛte (n. du. «зроблене й незроблене») ca dvandvāni («пари [протилежностей]») kadā śāntāni («коли втихли») kasya vā («і чиї») | evam jñātvā (герундив «так знаючи») iha nirvedāt («тут, з відречення») bhava («будь») tyāga-paraḥ («той, хто схильний до відкидання») avratī («без обітниць»).}
\vircom{`avratī` («без обітниць») — не дрібниця: текст водночас радить `tyāga` (відкидання) і знімає `vrata` (обітницю). Тобто відкидання не повинно стати новою дисципліною — інакше воно перетвориться на те, від чого відкидають. Питання `kasya vā` («і чиї ж вони?») залишене без відповіді навмисно.}
\virvarnone{власних приміток про різночитання немає; звіряння свідків не позначене як виконане}
\virsources{Avg\_9.1}{verse\_9.1}


\virs{9.2}{У якогось щасливого, любий, від споглядання того, як поводиться світ, бажання жити, голод і жага пізнавати дійшли до втишення.}
\virlit{«У якогось щасливого, любий, від споглядання того, як поводиться світ, бажання жити, голод і жага пізнавати дійшли до втишення.»}
\virgram{kasyāpi («якогось») tāta («любий, друже» — звертання) dhanyasya («щасливого») loka-ceṣṭā-avalokanāt («від споглядання того, як поводиться світ») | jīvitecchā («бажання жити») bubhukṣā («голод, жага їсти») ca bubhutsā-upaśamam gatāḥ («дійшли до втишення бажання пізнавати»).}
\vircom{Три подібні іменники стоять поруч: `jīvitecchā` (бажання жити), `bubhukṣā` (голод) і `bubhutsā` (жага пізнавати) — останнє те саме слово, що в 10.5, і словник має бути спільним. Причина названа одна й зовнішня: **споглядання** того, як поводиться світ, а не зусилля над собою. `tāta` — ласкаве звертання, не «батько».}
\virvarnone{власних приміток про різночитання немає; звіряння свідків не позначене як виконане}
\virsources{Avg\_9.2}{verse\_9.2}


\virs{9.3}{Усе це нестале, зіпсоване троїстою мукою, без суті, осуджене, гідне відкинення — так вирішивши, людина втихає.}
\virlit{«Усе це нестале, зіпсоване троїстою мукою, без суті, осуджене, гідне відкинення — так вирішивши, [людина] втихає.»}
\virgram{anityam (n. «нестале») sarvam eva idam («усе це») tāpa-tritaya-dūṣitam («зіпсоване троїстою мукою») | asāram («без суті») ninditam («осуджене») heyam («те, що слід відкинути») iti niścitya («так вирішивши») śāmyati («втихає»).}
\vircom{`tāpa-tritaya` («троїста мука») — усталена формула: мука від себе, від інших і від природи; переклад має лишити її як формулу, а не розтлумачувати в тексті вірша. П'ять означень ідуть підряд, і кожне сильніше за попереднє: від `anitya` (нестале) до `heya` (гідне відкинення). Втишення тут — **наслідок рішення**, а не почуття.}
\virvarnone{власних приміток про різночитання немає; звіряння свідків не позначене як виконане}
\virsources{Avg\_9.3}{verse\_9.3}


\virs{9.4}{Який же той час, вік або що — коли в людей немає пар протилежностей? Знехтувавши їх, хто живе як випало, той досяг би здійснення.}
\virlit{«Який же той час, вік або що — коли в людей немає пар протилежностей? Знехтувавши їх, хто живе як випало, той досяг би здійснення.»}
\virgram{kaḥ asau kālaḥ («який же той час») vayaḥ («вік») kim vā («або що») yatra dvandvāni na u nṛṇām («коли в людей немає пар [протилежностей]») | tāni upekṣya («їх знехтувавши») yathāprāpta-vartī («той, хто живе як випало») siddhim avāpnuyāt («досяг би здійснення»).}
\vircom{Питання риторичне й знецінювальне: не існує такого віку, коли пар немає, — отже, чекати на них марно. `yathāprāpta-vartī` («той, хто живе як випало») — той самий принцип, що `yadṛcchā` у 3.14 і 4.4, і в перекладі він має звучати так само. `avāpnuyāt` — умовний спосіб: «досяг би», а не «досягне»; це не обіцянка.}
\virvarnone{власних приміток про різночитання немає; звіряння свідків не позначене як виконане}
\virsources{Avg\_9.4}{verse\_9.4}

\index{tatha@\textit{tathā}}
\virs{9.5}{Побачивши, скільки різних поглядів мають великі мудреці, праведні й йоги, — хто, дійшовши відречення, не втихне?}
\virlit{«Побачивши різні погляди великих мудреців, праведних і йогів, хто, дійшовши відречення, не втихне?»}
\virgram{nānā matam (n. «різні погляди, думки») maharṣīṇām (gen. pl. «великих мудреців») sādhūnām («праведних») yoginām tathā | dṛṣṭvā («побачивши») nirvedam āpannaḥ («той, хто дійшов до відречення») kaḥ na śāmyati mānavaḥ («яка людина не втихла б»).}
\vircom{Риторичне питання з відповіддю «ніхто». Важливий перелік: `maharṣi`, `sādhu`, `yogin` — три різні авторитети, і текст показує, що **вони не згодні між собою**. Відречення тут походить не від безладдя світу, а від розбіжності авторитетів — це сильніше й конкретніше.}
\virvarnone{власних приміток про різночитання немає; звіряння свідків не позначене як виконане}
\virsources{Avg\_9.5}{verse\_9.5}


\virs{9.6}{Хто, пізнавши образ свідомості, — чи не вчитель? — той через відречення, рівність і з'єднаність переправляє з сансари.}
\virlit{«Хто, пізнавши образ свідомості, — чи не вчитель? — через відречення, рівність і з'єднаність переправляє [інших] із сансари.»}
\virgram{kṛtvā («зробивши») mūrti-parijñānam (n. «пізнання образу») caitanyasya (gen. «свідомості») na kim guruḥ («чи не вчитель») | nirveda-samatā-yuktyā («через відречення, рівність і з'єднаність») yaḥ tārayati saṃsṛteḥ («хто переправляє з сансари»).}
\vircom{Єдине місце в книзі, де `guru` названо **позитивно** і з означенням: учитель — той, хто пізнав `caitanya`. Але й тут текст не вимагає послуху: `tārayati` — «переправляє», дієслово руху, а не влади. Три складники (`nirveda`, `samatā`, `yukti`) подані як спосіб дії, а не як чесноти вчителя.}
\virvarnone{власних приміток про різночитання немає; звіряння свідків не позначене як виконане}
\virsources{Avg\_9.6}{verse\_9.6}


\virs{9.7}{Дивись на видозміни сущого як на самі лише елементи, як є насправді — і тієї ж миті, вільний від пут, станеш тим, хто стоїть у власній формі.}
\virlit{«Дивись на видозміни сущого як на самі лише елементи, як є насправді — і тієї ж миті, вільний від пут, станеш тим, хто стоїть у власній формі.»}
\virgram{paśya («дивись») bhūta-vikārān (m. acc. pl. «видозміни [того, що] є») tvam bhūta-mātrān («самі лише елементи») yathārthataḥ («як є насправді») | tat-kṣaṇāt («тієї ж миті») bandha-nirmuktaḥ («вільний від пут») svarūpa-sthaḥ («той, хто стоїть у власній формі») bhaviṣyasi («станеш»).}
\vircom{`tatkṣaṇāt` («тієї ж миті») — найсильніше слово вірша: не поступово, не через практику, а негайно. `bhūtamātra` («самі лише елементи») знімає оцінку з речей: вони не погані й не добрі, вони — складники. `svarūpastha` («той, хто стоїть у власній формі») — стан, а не дія; і `bhaviṣyasi` — майбутній час, тому це не опис того, що вже є, а вказівка на те, що станеться.}
\virvarnone{власних приміток про різночитання немає; звіряння свідків не позначене як виконане}
\virsources{Avg\_9.7}{verse\_9.7}

\index{yatha@\textit{yathā}}\index{tatha@\textit{tathā}}
\virs{9.8}{Сансара — це саме усталені схильності, так; відпусти їх усі. Їх відкидання — через відкидання схильностей; стояння в собі — нині, так чи інакше.}
\virlit{«Сансара — це саме усталені схильності, так; відпусти їх усі. Їх відкидання — через відкидання схильностей; стояння [в собі] — нині, так чи інакше.»}
\virgram{vāsanāḥ (f. pl. «усталені схильності, просочення») eva saṃsāraḥ iti («сансара — це саме [вони], так») sarvāḥ vimuñca (imper. «відпусти всі») tāḥ | tat-tyāgaḥ («їх відкидання») vāsanā-tyāgāt («через відкидання схильностей») sthitiḥ adya yathā tathā («стояння — нині, так чи інакше»).}
\vircom{Вірш закінчує главу визначенням, а не обіцянкою: сансара — не світ, а `vāsanā`, усталені схильності. `yathā tathā` («так чи інакше») знімає вимогу до способу: результат не залежить від того, як саме це станеться. Друга половина повторює те саме слово двічі (`tyāga`), і повтор тут не недогляд — він показує, що відкидання схильностей і є відкиданням, а не засобом до нього.}
\virvarnone{власних приміток про різночитання немає; звіряння свідків не позначене як виконане}
\virsources{Avg\_9.8}{verse\_9.8}

\pilotchapter{10}{8}

\virs{10.1}{Покинувши вороже бажання і сповнену лиха користь — навіть дхарму, що є причиною обох, — усюди зневажай.}
\virlit{«Покинувши вороже бажання й користь, сповнену лиха, — навіть дхарму, що є причиною обох, — всюди зневажай.»}
\virgram{vihāya (герундив від vi-hā «покинувши») vairiṇam (m. acc. sg. «ворожого») kāmam (m. acc. sg. «бажання») artham (m. acc. sg. «користь, багатство») ca anartha-saṃkulam («сповнене лиха») | dharmam api (m. acc. sg. «навіть дхарму») etayoḥ hetum («причину обох») sarvatra anādaram kuru («всюди зневажай»).}
\vircom{Дхарма стоїть у переліку **останньою** і з підсиленням api («навіть»): спершу відкидають бажання й користь, а потім те, що їх породжує. anartha-saṃkula («сповнене лиха») стосується саме artha, і переклад не повинен переносити цю оцінку на дхарму. «Зневажай» (anādara) — сильне слово; м'якше «не зважай» змінило б ступінь.}
\virvarnone{власних приміток про різночитання немає; звіряння свідків не позначене як виконане}
\virsources{Avg\_10.1}{verse\_10.1}


\virs{10.2}{Дивись на друзів, поле, гроші, дім, дружину, спадок і решту багатств як на сон і чари — на три дні або п'ять.}
\virlit{«Дивись на друзів, поле, гроші, дім, дружину, спадок та інші багатства як на сон і чари — на три дні або на п'ять.»}
\virgram{svapna-indrajāla-vat («як сон і маґія», indrajāla — «мережа Індри», тобто чари) paśya (imper. 2 sg. «дивись») dināni trīṇi pañca vā («три дні або п'ять») | mitra-kṣetra-dhana-āgāra-dāra-dāyā-ādi-saṃpadaḥ (f. acc. pl. «багатства: друзі, поле, гроші, дім, дружина, спадок тощо»).}
\vircom{Перелік навмисно побутовий: не «світ» загалом, а конкретні речі, за які тримаються. «Три дні або п'ять» — не точний строк, а спосіб сказати «коротко»; переклад зберігає числа, бо вони й творять ефект. indrajāla — не просто «ілюзія», а фокус, свідома підміна видимого.}
\virvarnone{власних приміток про різночитання немає; звіряння свідків не позначене як виконане}
\virsources{Avg\_10.2}{verse\_10.2}


\virs{10.3}{Де б не виникла спрага — знай: саме там і сансара. Опираючись на зрілу безсторонність, стань тим, у кого спрага минула, — і будь щасливий.}
\virlit{«Де б не виникла спрага — знай: саме там сансара. Опираючись на зрілу безсторонність, стань тим, у кого спрага минула, і будь щасливий.»}
\virgram{yatra yatra («де б не») bhavet (opt. 3 sg. «виникла б») tṛṣṇā (f. nom. sg. «спрага») saṃsāram (m. acc. sg.) viddhi (imper. 2 sg. vid «знай») tatra vai («саме там») | prauḍha-vairāgyam (n. acc. sg. «зрілу безсторонність») āśritya (герундив «опираючись на») vīta-tṛṣṇaḥ («той, у кого спрага минула») sukhī bhava («будь щасливий»).}
\vircom{Тут уперше в розділі дано не заперечення, а **вказівку**: `yatra yatra … tatra` («де б … саме там»). Сансара — не місце й не стан світу, а точка, де виникла спрага. prauḍha («зріла») важливе: безсторонність має бути витриманою, а не поривчастою. Висновок «будь щасливий» — наказ, а не побажання.}
\virvarnone{власних приміток про різночитання немає; звіряння свідків не позначене як виконане}
\virsources{Avg\_10.3}{verse\_10.3}

\index{tosa@\textit{toṣa}}\index{bahuvrihi@\textit{bahuvrīhi}}
\virs{10.4}{Полон має природу самої лише спраги; її знищення називають визволенням. Від самої лише неприв'язаності до буття приходить вдоволеність досягнутим — знову й знову.}
\virlit{«Полон має природу самої лише спраги; її знищення називають визволенням. Від самої лише неприв'язаності до буття — вдоволеність досягнутим, знову й знову.»}
\virgram{tṛṣṇā-mātrātmakaḥ (bahuvrīhi «чия сутність — сама лише спрага») bandhaḥ | tat-nāśaḥ (m. «його знищення») mokṣaḥ ucyate («називається») | bhava-asaṃsakti-mātreṇa («самою лише неприв'язаністю до буття») prāpti-tuṣṭiḥ («вдоволеність досягнутим») muhur muhuḥ («знову й знову»).}
\vircom{Тут полон означено **однією** ознакою — спрага, і визволення теж однією — її знищення. Це найточніше місце глави, і воно коротке; розширювати переклад не варто. `toṣa` входить сюди як складник `prāpti-tuṣṭi` («вдоволеність досягнутим»), а не як самостійна чеснота — це варто пам'ятати при звірці з 1.2, де `toṣa` стоїть окремо. `muhur muhuḥ` — не «завжди», а «раз у раз»: вдоволеність повторюється, а не тримається одного разу назавжди.}
\virvarnone{власних приміток про різночитання немає; звіряння свідків не позначене як виконане}
\virsources{Avg\_10.4}{verse\_10.4}

\index{tatha@\textit{tathā}}
\virs{10.5}{Ти один — свідомий, чистий; світ же — неживий і несправжній. Навіть незнання — аніщо. То яке в тебе, при всьому тому, бажання пізнавати?}
\virlit{«Ти один — свідомий, чистий; світ — неживий і несправжній. Навіть незнання — аніщо. Яке ж тоді в тебе бажання пізнавати?»}
\virgram{tvam ekaḥ cetanaḥ śuddhaḥ («ти один — свідомий, чистий») | jaḍam (n. «неживе») viśvam asat (n. «несправжнє») tathā | avidyā api na kiñcit sā («і незнання — аніщо») | kā bubhutsā (f. «яке бажання пізнавати») tathāpi te (gen. «все ж у тебе»).}
\vircom{`asat` тут — «несправжнє», не «зло»; це різні речі, і переклад не повинен їх зливати. Питання `kā bubhutsā` замикає коло: якщо навіть незнання ніщо, то пізнавати нічого — і саме тому питання адресоване тому, хто вже слухає пояснення. Це найгостріше місце глави: воно ставить під сумнів користь самої розмови.}
\virvarnone{власних приміток про різночитання немає; звіряння свідків не позначене як виконане}
\virsources{Avg\_10.5}{verse\_10.5}


\virs{10.6}{Царство, сини, дружини, тіла й утіхи — усе це вже гинуло для тебе, прив'язаного, — з народження в народження.}
\virlit{«Царство, сини, дружини, тіла й утіхи — усе це гинуло для тебе, прив'язаного, з народження в народження.»}
\virgram{rājyam (n. «царство») sutāḥ (m. pl. «сини») kalatrāṇi (n. pl. «дружини») śarīrāṇi (n. pl. «тіла») sukhāni (n. pl. «утіхи») ca | saṃsaktasya api (gen. «хоч і прив'язаного») naṣṭāni (n. pl. pp. «загинули») tava (gen.) janmani janmani (loc. «в народженні за народженням»).}
\vircom{`naṣṭāni` — доконаний дієприкметник: не «загине», а «гинуло вже не раз». Переклад має зберегти цю повторюваність, інакше вірш звучить як погроза, а він — як нагадування. Звертання адресоване цареві, і царство названо **першим** у переліку: порядок не випадковий.}
\virvarnone{власних приміток про різночитання немає; звіряння свідків не позначене як виконане}
\virsources{Avg\_10.6}{verse\_10.6}

\index{samsara@\textit{saṃsāra}}
\virs{10.7}{Досить уже й багатства, і бажання, і навіть добрих учинків: у цьому лісі сансари твій розум так і не знав відпочинку.}
\virlit{«Досить уже й багатства, і бажання, і навіть добрих учинків: у цьому лісі сансари твій розум не знав відпочинку.»}
\virgram{alam (indecl. «досить») arthena (instr. «багатством») kāmena sukṛtena api karmaṇā («і навіть добрим учинком») | ebhyaḥ (abl. pl. «від них») saṃsāra-kāntāre (m. loc. «у лісі сансари») na viśrāntam (n. nom. «не відпочив») abhūt (3 sg. bhū) manaḥ (n. «розум»).}
\vircom{`sukṛta` тут — «заслуга, добре зроблене», а не «добро» взагалі; і текст відкидає саме заслугу, ставлячи її в один ряд із багатством і бажанням. Це найнезручніше місце глави для читача, який звик, що добрі вчинки беззастережно рахуються. `kāntāra` — «ліс», образ блукання, а не місце. `alam` — «годі вже», з інтонацією втоми, а не заборони.}
\virvarnone{власних приміток про різночитання немає; звіряння свідків не позначене як виконане}
\virsources{Avg\_10.7}{verse\_10.7}

\index{duhkha@\textit{duḥkha}}
\virs{10.8}{Скільки ж за ці народження зроблено тілом, розумом і мовою діл, що дають біль і втому! Тож нехай це припиниться хоч тепер.}
\virlit{«Скільки ж за ці народження зроблено тілом, розумом і мовою діл, що дають біль і втому! Тож нехай це припиниться хоч тепер.»}
\virgram{kṛtam (n. pp. «зроблено») na kati («чи мало, скільки ж») janmāni (n. pl. «народження») kāyena manasā girā (instr. «тілом, розумом, мовою») | duḥkha-āyāsa-dam («те, що дає біль і втому») karma tat adya api uparamyatām (пасивний імператив «нехай буде припинене»).}
\vircom{Потрійний інструментал — тіло, розум, мова — це класична формула трьох дверей дії, і переклад має зберегти всі три, а не згортати їх у «вчинки». `uparamyatām` — пасивний імператив: не «припини», а «нехай припиниться»; дія знімається не зусиллям того, хто діє. `adya api` — «хоч тепер», з відтінком докору, а не нейтрального «зараз». Цим віршем глава 10 закінчується.}
\virvarnone{власних приміток про різночитання немає; звіряння свідків не позначене як виконане}
\virsources{Avg\_10.8}{verse\_10.8}

\pilotchapter{11}{8}

\virs{11.1}{«Видозміна буття й небуття — за власною природою» — хто має таке переконання, той без видозмін, з минулими турбаціями, втихає легко.}
\virlit{«Видозміна буття й небуття — за власною природою» — хто має таке переконання, той без видозмін, з минулими турбаціями, втихає легко.}
\virgram{bhāva-abhāva-vikāraḥ (m. «видозміна буття й небуття») ca svabhāvāt («за власною природою») iti niścayī («той, хто має таке переконання») | nirvikāraḥ («без видозмін») gata-kleśaḥ («той, у кого минули турбації») sukhena eva upaśāmyati («втихає саме з легкістю»).}
\vircom{`niścayī` («той, хто має переконання») повторюється на початку 11.1–11.4 — це структурна формула всієї глави, і переклад має тримати її однаково. `nirvikāraḥ` уточнює: не «без змін» у сенсі нерухомості, а без **власних** видозмін від того, що змінюється назовні. `sukhena eva` («саме з легкістю») — втишення не коштує зусиль, і це важливо.}
\virvarnone{власних приміток про різночитання немає; звіряння свідків не позначене як виконане}
\virsources{Avg\_11.1}{verse\_11.1}


\virs{11.2}{«Творець усього — Ішвара, і немає тут іншого» — хто має таке переконання, той із розчиненими всередині сподіваннями, втихлий, ніде не застрягає.}
\virlit{«Творець усього — Ішвара, і немає тут іншого» — хто має таке переконання, той із розчиненими всередині сподіваннями, втихлий, ніде не застрягає.}
\virgram{īśvaraḥ (m. nom. «Ішвара, володар») sarva-nirmātā («творець усього») na iha anyaḥ («немає тут іншого») iti niścayī | antargalita-sarvāśaḥ («той, у кого всі сподівання розчинилися всередині») śāntaḥ kvāpi na sajjate («ніде не застрягає»).}
\vircom{Тут уперше й востаннє в книзі з'являється `īśvara` як творець — і одразу з запереченням іншого (`na iha anyaḥ`). Це не теологія, а спосіб зняти питання причини. `sajjate` — «застрягає, чіпляється»: наслідок не в тому, що людина стала всемогутньою, а в тому, що вона більше ні за що не тримається.}
\virvarnone{власних приміток про різночитання немає; звіряння свідків не позначене як виконане}
\virsources{Avg\_11.2}{verse\_11.2}


\virs{11.3}{«Нещастя й достаток у свій час — саме від долі» — хто має таке переконання, той вдоволений, із чуттями в порядку, повсякчас не бажає й не сумує.}
\virlit{«Нещастя й достаток у свій час — саме від долі» — хто має таке переконання, той вдоволений, із чуттями в порядку, повсякчас не бажає й не сумує.}
\virgram{āpadaḥ (f. pl. «нещастя») saṃpadaḥ («достаток») kāle («у свій час») daivāt eva («саме від долі») iti niścayī | tṛptaḥ («вдоволений») svasthendriyaḥ («той, чиї чуття в порядку») nityam na vāñchati na śocati.}
\vircom{`daiva` — «доля» як те, що розподіляє, а не як примха. Формула «не бажає й не сумує» — та сама пара, що в 8.1–8.2, і переклад має її впізнавати. `svasthendriya` («із чуттями в порядку») — рідкісний складник: ідеться не про приборкання чуттів, а про їхню справність.}
\virvarnone{власних приміток про різночитання немає; звіряння свідків не позначене як виконане}
\virsources{Avg\_11.3}{verse\_11.3}

\index{sukha@\textit{sukha}}
\virs{11.4}{«Щастя й страждання, народження й смерть — саме від долі» — хто має таке переконання, той, не бачачи нічого досяжного, без натуги, навіть діючи, не забруднюється.}
\virlit{«Щастя й страждання, народження й смерть — саме від долі» — хто має таке переконання, той, не бачачи нічого досяжного, без натуги, навіть діючи, не забруднюється.}
\virgram{sukha-duḥkhe (n. du. «щастя й страждання») janma-mṛtyū (m. du. «народження й смерть») daivāt eva iti niścayī | sādhya-adarśī («той, хто не бачить [нічого] досяжного») nirāyāsaḥ («без натуги») kurvan api na lipyate («навіть діючи, не забруднюється»).}
\vircom{`sādhyādarśī` («той, хто не бачить досяжного») — найтонше слово глави: дія не припиняється, але зникає те, заради чого вона нібито робиться. `kurvann api na lipyate` — «навіть діючи, не забруднюється»: це не відмова від дії, а її розпредметнення. Перелік розширюється від 11.3 до 11.4 (додано народження й смерть), і це наростання навмисне.}
\virvarnone{власних приміток про різночитання немає; звіряння свідків не позначене як виконане}
\virsources{Avg\_11.4}{verse\_11.4}


\virs{11.5}{«Страждання народжується з думання, не інакше» — хто має таке переконання, той, позбавлений його, щасливий, втихлий, скрізь із опалою спрагою.}
\virlit{«Страждання народжується з думання, не інакше» — хто має таке переконання, той, позбавлений його, щасливий, втихлий, скрізь із опалою спрагою.}
\virgram{cintayā (instr. «думанням, турботою») jāyate duḥkham («народжується страждання») na anyathā («не інакше») iti niścayī | tayā hīnaḥ («позбавлений її») sukhī śāntaḥ sarvatra galita-spṛhaḥ («скрізь із опалою спрагою»).}
\vircom{`cintā` — не «думка» взагалі, а **турботливе обертання думки**, і саме воно назване джерелом страждання. Переклад «думання» тут кращий за «думка», бо остання звучить нейтрально. `galitaspṛha` («із опалою спрагою») — образ, який текст уже вживав у 11.2 (`antargalita`); словник лишається спільним.}
\virvarnone{власних приміток про різночитання немає; звіряння свідків не позначене як виконане}
\virsources{Avg\_11.5}{verse\_11.5}


\virs{11.6}{«Я не тіло, і тіло не моє; я — усвідомлення» — хто має таке переконання, той, наче досягши окремішності, не пам'ятає незробленого як зробленого.}
\virlit{«Я не тіло, і тіло не моє; я — усвідомлення» — хто має таке переконання, той, наче досягши окремішності, не пам'ятає незробленого як зробленого.}
\virgram{na aham dehaḥ na me dehaḥ («я не тіло, і тіло не моє») bodhaḥ aham («я — усвідомлення») iti niścayī | kaivalyam iva saṃprāptaḥ («наче досяг окремішності») na smarati («не пам'ятає») akṛtam kṛtam («незробленого як зробленого»).}
\vircom{Формула «не тіло, і тіло не моє» та сама, що в 8.4 і 2.22 — тут вона подана як переконання, а не як висновок. `kaivalya` («окремішність») з'являється один раз і з `iva` («наче»): текст не каже, що стан досягнуто. `akṛtaṃ kṛtam` — «незроблене як зроблене»: ідеться не про забудькуватість, а про те, що минуле більше не рахується.}
\virvarnone{власних приміток про різночитання немає; звіряння свідків не позначене як виконане}
\virsources{Avg\_11.6}{verse\_11.6}


\virs{11.7}{«Від Брахми до стеблини — це саме я» — хто має таке переконання, той без домислювання, чистий, втихлий, вдоволений і досягнутим, і недосягнутим.}
\virlit{«Від Брахми до стеблини — це саме я» — хто має таке переконання, той без домислювання, чистий, втихлий, вдоволений і досягнутим, і недосягнутим.}
\virgram{ā-brahma-stamba-paryantam («від Брахми до стеблини») aham eva («саме я») iti niścayī | nirvikalpaḥ («без домислювання») śuciḥ («чистий») śāntaḥ prāpta-aprāpta-vinirvṛtaḥ («вдоволений і досягнутим, і недосягнутим»).}
\vircom{`ābrahmastambaparyantam` — та сама формула межі, що в 2.11 (`brahmādistambaparyantam`), і переклад має бути однаковий. `prāptāprāpta-vinirvṛta` — «вдоволений і досягнутим, і недосягнутим»: це не байдужість, а відсутність різниці між двома станами. `śuci` тут — не ритуальна чистота, а відсутність домішки, бо стоїть поруч із `nirvikalpa`.}
\virvarnone{власних приміток про різночитання немає; звіряння свідків не позначене як виконане}
\virsources{Avg\_11.7}{verse\_11.7}


\virs{11.8}{«Цей світ, різноманітно дивовижний, — аніщо» — хто має таке переконання, той без усталених схильностей, саме лише сяйнування, втихає, наче аніщо.}
\virlit{«Цей світ, різноманітно дивовижний, — аніщо» — хто має таке переконання, той без усталених схильностей, саме лише сяйнування, втихає, наче аніщо.}
\virgram{nānā-āścaryam («різноманітно дивовижне») idam viśvam na kiñcit («аніщо») iti niścayī | nirvāsanaḥ («без усталених схильностей») sphūrti-mātraḥ («саме лише сяйнування») na kiñcit iva śāmyati («втихає, наче аніщо»).}
\vircom{Глава закінчується тим самим словом, що й почалася в 11.1: `niścayī`, але тепер воно стоїть при «аніщо». `nirvāsana` («без усталених схильностей») — те саме `vāsanā`, що в 9.8; словник спільний. `sphūrtimātra` («саме лише сяйнування») і `na kiñcid iva` («наче аніщо») разом дають обережність, яку легко втратити: текст не каже, що людина зникає, — він каже «наче».}
\virvarnone{власних приміток про різночитання немає; звіряння свідків не позначене як виконане}
\virsources{Avg\_11.8}{verse\_11.8}

\pilotchapter{12}{8}

\virs{12.1}{Спершу нездатний до тілесних справ, потім нездатний до поширення мови, далі нездатний до думання — тому саме в цьому я стою.}
\virlit{«Спершу нездатний до тілесних справ, потім нездатний до поширення мови, далі нездатний до думання — тому саме в цьому я стою.»}
\virgram{kāya-kṛtya-asahaḥ («нездатний до тілесних справ») pūrvam tataḥ vāk-vistara-asahaḥ («потім нездатний до поширення мови») | atha cintā-asahaḥ («далі нездатний до думання») tasmāt evam eva aham āsthitaḥ («тому саме в цьому я стою»).}
\vircom{Тут відмова йде **не за бажанням, а через неспроможність**: `asahaḥ` — «нездатний витримати», а не «той, хто відмовився». Це рідкісний для тексту випадок, коли стан описано знизу — через те, чого людина вже не може. Формула `evam evāham āsthitaḥ` повторюється в 12.1–12.4, як `niścayī` у главі 11: переклад мусить тримати її однаково.}
\virvarnone{власних приміток про різночитання немає; звіряння свідків не позначене як виконане}
\virsources{Avg\_12.1}{verse\_12.1}

\index{atman@\textit{ātman}}
\virs{12.2}{Через відсутність прихильності до звуку й подібного та через невидимість його для самого себе — серце, для якого розсіяність і зосередженість одне — саме в цьому я стою.}
\virlit{«Через відсутність прихильності до звуку й подібного та через невидимість [його] для самого себе — серце, для якого розсіяність і зосередженість одне — саме в цьому я стою.»}
\virgram{prīti-abhāvena («через відсутність прихильності») śabda-ādeḥ («до звуку й подібного») adṛśyatvena ca ātmanaḥ («і через невидимість для себе самого») | vikṣepa-aikāgrya-hṛdayaḥ («той, чиє серце — і розсіяність, і зосередженість разом») evam eva aham āsthitaḥ.}
\vircom{`vikṣepa-aikāgrya` в одному складнику — найважливіше місце вірша: розсіяність і зосередженість названі **одним** словом, і це не поетична вільність, а твердження, що для того, хто стоїть у собі, різниці між ними немає. `adṛśyatva` («невидимість») стосується `ātman` — себе не видно собі як об'єкт.}
\virvarnone{власних приміток про різночитання немає; звіряння свідків не позначене як виконане}
\virsources{Avg\_12.2}{verse\_12.2}


\virs{12.3}{Коли розсіяність походить від прагнення самадгі та подібного, то й звичайна діяльність — задля самадгі. Так побачивши це правило, саме в цьому я стою.}
\virlit{«Коли розсіяність походить від прагнення самадгі та подібного, то й звичайна діяльність — задля самадгі. Так побачивши це правило, саме в цьому я стою.»}
\virgram{samādhi-āsādi-vikṣiptau («у розсіяності, [викликаній] прагненням самадгі та подібного») vyavahāraḥ samādhaye («[звичайна] діяльність — для самадгі») | evam vilokya niyamam («так побачивши правило») evam eva aham āsthitaḥ.}
\vircom{Найдошкульніший вірш глави: **прагнення самадгі саме стає розсіяністю** (`vikṣipti`). І далі — іронія, яку легко проґавити: якщо навіть буденна діяльність робиться «задля самадгі», то правило (`niyama`) виявляється всеохопним, і виходу з нього немає. Переклад мусить зберегти цю замкненість, а не зробити з неї пораду.}
\virvarnone{власних приміток про різночитання немає; звіряння свідків не позначене як виконане}
\virsources{Avg\_12.3}{verse\_12.3}


\virs{12.4}{Через відсутність того, що відкидати й приймати, а отже — через відсутність радості й смутку, нині, о Брахмане, саме в цьому я стою.}
\virlit{«Через відсутність того, що відкидати й приймати, а отже — через відсутність радості й смутку, нині, о Брахмане, саме в цьому я стою.»}
\virgram{heya-upādeya-virahāt («через відсутність того, що відкидати й приймати») evam harṣa-viṣādayoḥ («так — радості й смутку») abhāvāt («через відсутність») adya he brahman («нині, о Брахмане») | наприкінці в SARIT стоїть `nevamevāhamāsthitaḥ` — sandhi, яке можна розкласти двома способами (див. variants).}
\vircom{`he brahman` — звертання до Джанаки як до брахмана, єдине таке в розділі. Ланцюг причинний: немає пари «відкидати/приймати» → немає радості й смутку → нема чого досягати. **Відкрите питання:** у SARIT наприкінці стоїть `nevamevāhamāsthitaḥ` (sandhi `na + evameva`), і я прочитав це як ствердне `evameva` за паралелізмом. Перевірка решти глави це підтверджує лише частково: 12.5, 12.6 і 12.7 усі закінчуються на `evamevāhamāsthitaḥ` **без** `na`, тобто 12.4 — єдиний виняток у главі. Отже `na` може бути або справжнім (і тоді сенс вірша протилежний), або помилкою транскрипції. **GRETIL недоступний, другого свідка немає** — картка лишається чернеткою, а не «перевірено».}
\virvars{\virvaritem{SARIT, решта глави 12 (12.5–12.7)}{У 12.4 формула записана як `nevamevāhamāsthitaḥ`, тоді як 12.5, 12.6 і 12.7 закінчуються на `evamevāhamāsthitaḥ` без `na`. Отже 12.4 — єдиний виняток у главі: `na` може бути справжнім (і тоді сенс вірша протилежний) або помилкою транскрипції. З GRETIL не звірено, тому це зафіксоване спостереження, а не підтверджене різночитання.}}
\virsources{Avg\_12.4}{verse\_12.4}

\index{asrama@\textit{āśrama}}
\virs{12.5}{Медитація зі стадією життя і без неї, прийняття й відкидання, зроблені умом, — усе це домислювання. Побачивши це в собі, через саме це я так стою.}
\virlit{«Медитація зі стадією життя і без неї, прийняття й відкидання, зроблені умом, — усе це домислювання. Побачивши це у собі, через саме це я так стою.»}
\virgram{āśrama-anāśramaṃ dhyānam («медитація зі стадією життя і без неї») citta-svīkṛta-varjanam («прийняття й відкидання, зроблені умом») | vikalpam («домислювання») mama vīkṣya («моє побачивши») etaiḥ evam eva aham āsthitaḥ («через саме це я так стою»).}
\vircom{Єдине місце в книзі, де `āśrama` (стадія життя) названо **поряд з її відсутністю** — і обидва варіанти відкинуті як `vikalpa`. Це важливо для нашого словника: текст не обирає між «життям у стадії» і «поза нею», він знімає саму пару. `citta-svīkṛta-varjana` — прийняття й відкидання не речей, а **вироблені умом**; переклад мусить зберегти, що йдеться про акт ума, а не про оцінку світу.}
\virvarnone{власних приміток про різночитання немає; звіряння свідків не позначене як виконане}
\virsources{Avg\_12.5}{verse\_12.5}

\index{ajnana@\textit{ajñāna}}\index{yatha@\textit{yathā}}\index{tatha@\textit{tathā}}
\virs{12.6}{Виконання обрядів через незнання і припинення їх — однаково. Зрозумівши правильно цю суть, я так стою.}
\virlit{«Виконання обрядів через незнання і припинення їх — однаково. Зрозумівши правильно цю суть, я так стою.»}
\virgram{karma-anuṣṭhānam («виконання обрядів») ajñānāt («через незнання») yathā eva uparamaḥ tathā («припинення — так само») | buddhvā samyak («зрозумівши правильно») idam tattvam evam eva aham āsthitaḥ.}
\vircom{`yathaiva … tathā` («так само») ставить дію й відмову від дії на один рівень — і обидві названі `ajñāna`, коли виходять із незнання. Це не заклик припинити обряди: текст каже, що припинення не краще за виконання. Переклад мусить звучати рівно, без інтонації «краще не робити».}
\virvarnone{власних приміток про різночитання немає; звіряння свідків не позначене як виконане}
\virsources{Avg\_12.6}{verse\_12.6}


\virs{12.7}{Незбагненне, навіть коли його осмислюють, набирає форми думки. Тому, відкинувши те плекання, я так стою.}
\virlit{«Незбагненне, навіть коли його осмислюють, набирає форми думки. Тому, відкинувши те плекання, я так стою.»}
\virgram{acintyam («незбагненне») cintyamānaḥ api («навіть коли його осмислюють») cintā-rūpam bhajati asau («воно набирає форми думки») | tyaktvā tad-bhāvanam («відкинувши те плекання») tasmāt evam eva aham āsthitaḥ.}
\vircom{Механізм названий точно: незбагненне не залишається собою, коли його осмислюють — воно **набирає форми думки** (`cintārūpa`). Тому відкидається не предмет, а `bhāvana` — плекання, повторюване обертання. Це та сама `cintā`, що в 11.5, і словник має бути спільним.}
\virvarnone{власних приміток про різночитання немає; звіряння свідків не позначене як виконане}
\virsources{Avg\_12.7}{verse\_12.7}


\virs{12.8}{У кого це так зроблено — той стає тим, чию мету здійснено; у кого така природа — той стає тим, чию мету здійснено.}
\virlit{«У кого це так зроблено — той стає тим, чию мету здійснено; у кого така природа — той стає тим, чию мету здійснено.»}
\virgram{evam eva kṛtam yena («у кого це так зроблено») saḥ kṛta-arthaḥ bhavet asau | evam eva svabhāvaḥ yaḥ («у кого така природа») saḥ kṛta-arthaḥ bhavet asau.}
\vircom{Вірш повторює рядок двічі й міняє одне слово: `kṛtam` (зроблено) → `svabhāvaḥ` (природа). Це не тавтологія, а розрізнення: спершу те, що людина зробила, потім те, чим вона стала. `kṛtārtha` — «той, чию мету здійснено», усталене слово; переклад тримає його однаково в обох половинах. Глава закінчується не настановою, а констатацією.}
\virvarnone{власних приміток про різночитання немає; звіряння свідків не позначене як виконане}
\virsources{Avg\_12.8}{verse\_12.8}

\pilotchapter{13}{7}

\virs{13.1}{Стан неволодіння — це бути в собі; навіть у жебрацькому поясі він рідкісний. Відкинувши віддавання й приймання, я сиджу, як мені завгодно.}
\virlit{«Стан неволодіння — це бути в собі; навіть у [носінні] жебрацького пояса він рідкісний. Відкинувши віддавання й приймання, тому я сиджу, як мені завгодно.»}
\virgram{akiṃcanabhavam (n. nom. sg. «стан неволодіння, буття-без-нічого»: a-kiṃcana «той, у кого нічого немає» + bhava «стан») svāsthyam (n. nom. sg. «перебування в собі, здоров'я, лад»: sva «своє» + stha «той, що стоїть»). kaupīnatve (n. loc. sg. «у [носінні] каупіни»: kaupīna «жебрацький пояс, стегнова пов'язка» — знак відречення) api («навіть») durlabham (n. nom. sg. «важкодосяжне, рідкісне»). tyāgādāne (n. nom./acc. du. «віддавання й приймання») vihāya (герунд vi-hā «відкинувши») asmāt («тому») aham āse (1 sg. pr. ās «я сиджу, перебуваю») yathāsukham («як мені завгодно, як зручно»).}
\vircom{Найважливіше тут — svāsthya. Слово означає буквально «перебування в собі» і в сучасній мові значить «здоров'я»; у вірші воно вжите в першому сенсі, і це не метафора, а визначення: бути здоровим — значить бути в собі. Друга точність — kaupīna: жебрацький пояс названий не як приклад аскези, а як те, у чому svāsthya все одно рідкісне. Тобто зовнішня ознака відречення нічого не гарантує. Пара tyāga / ādāna («віддавання» / «приймання») — та сама, що heya / upādeya у 7.5, і тут її теж відкидають цілою, а не по частинах. РЕФРЕН ГЛАВИ. Шість із семи віршів (13.1, 13.3–13.7) закінчуються однаково: aham āse yathāsukham («я сиджу — як мені завгодно»). ās — не «бути», а саме СИДІТИ, перебувати на місці; yathāsukham — «як зручно, за бажанням». Разом це формула, яка скасовує будь-яку поставу як вимогу: не треба ні стояти, ні лежати правильно. 13.2 — єдиний вірш глави без цього рефрену, і його треба читати окремо.}
\virvars{\virvaritem{SARIT і GRETIL (перевірено scripts/fetch\_witnesses.py)}{Читання збігаються; відмінності лише в аваґрасі та поділі слів. Не різночитання.}}
\virsources{Avg\_13.1}{verse\_13.1}

\index{duhkha@\textit{duḥkha}}
\virs{13.2}{Десь втомлюється тіло, десь втомлюється язик, десь розум — відкинувши це, я стою в меті людини, щасливо.}
\virlit{«Десь [є] втома тіла, десь втомлюється язик, десь розум — відкинувши це, [я] стою в меті людини, щасливо.»}
\virgram{kutrāpi («десь, у чомусь») khedaḥ (m. nom. sg. «втома, стомленість, докука») kāyasya (m. gen. sg. «тіла»). jihvā (f. nom. sg. «язик, мова») kutrāpi khidyate (3 sg. pr. pass. khid «втомлюється, знемагає»). manaḥ (n. nom. sg. «розум») kutrāpi. tat (n. acc. sg. «те») tyaktvā (герунд tyaj «відкинувши») puruṣārthe (m. loc. sg. «у меті людини»: puruṣa «людина» + artha «мета»; чотири мети — дгарма, користь, бажання, визволення) sthitaḥ (m. nom. sg. «той, хто стоїть») sukham (n. nom./acc. sg. «щасливо, зручно»).}
\vircom{ЄДИНИЙ вірш глави без рефрену aham āse yathāsukham, і це не випадково: тут ідеться не про те, як сидіти, а про те, ЩО втомлюється. Три рівні названі поспіль — тіло, язик, розум — і кожен втомлюється «десь», тобто не завжди й не однаково. Слово kheda («втома») м'якше за duḥkha («страждання») з глави 8: це не біль, а зношування. puruṣārtha («мета людини») вжите в однині, хоч традиційно їх чотири, — тобто йдеться не про вибір між ними, а про те, що всі вони вже не потрібні.}
\virvars{\virvaritem{SARIT і GRETIL (перевірено scripts/fetch\_witnesses.py)}{Читання збігаються; відмінності лише в аваґрасі та поділі слів. Не різночитання.}}
\virsources{Avg\_13.2}{verse\_13.2}


\virs{13.3}{Подумавши по-справжньому: «зробленого не буде нічого» — що б не випало зробити, зробивши те, я сиджу, як мені завгодно.}
\virlit{«Подумавши істинно: „зробленого не буде анічого“ — що б не випало зробити, зробивши те, я сиджу, як мені завгодно.»}
\virgram{kṛtam (n. nom. sg. «зроблене») kim api («що-небудь, анічого») na eva syāt («нехай не буде») iti («так»). saṃcintya (герунд sam-cint «подумавши, зваживши») tattvataḥ («істинно, по суті»). yadā yat («коли що») kartum (інфінітив kṛ «зробити») āyāti (3 sg. pr. ā-i «приходить, випадає») tat (n. acc. sg.) kṛtvā (герунд «зробивши») āse (1 sg. pr. ās «я сиджу») yathāsukham.}
\vircom{Парадокс у тому, що дія не скасовується: yadā yat kartum āyāti («коли що випадає зробити») — дієслово āyāti означає саме те, що приходить само, без запрошення. Отже, діють, але без того, хто міг би зарахувати її собі. kṛtvāse — злитна форма kṛtvā + āse, і в ній дві дії стоять впритул: зробивши — сиджу. Умова naiva syāt («нехай не буде») стоїть у optativus, тобто це не констатація, а спосіб бачення, прийнятий заздалегідь. РЕФРЕН ГЛАВИ. Шість із семи віршів (13.1, 13.3–13.7) закінчуються однаково: aham āse yathāsukham («я сиджу — як мені завгодно»). ās — не «бути», а саме СИДІТИ, перебувати на місці; yathāsukham — «як зручно, за бажанням». Разом це формула, яка скасовує будь-яку поставу як вимогу: не треба ні стояти, ні лежати правильно. 13.2 — єдиний вірш глави без цього рефрену, і його треба читати окремо.}
\virvars{\virvaritem{SARIT і GRETIL (перевірено scripts/fetch\_witnesses.py)}{Читання збігаються; відмінності лише в аваґрасі та поділі слів. Не різночитання.}}
\virsources{Avg\_13.3}{verse\_13.3}


\virs{13.4}{У йога, що перебуває в тілі, є стани, до яких його змушують дія й недіяння; та через відсутність з'єднання й роз'єднання я сиджу, як мені завгодно.}
\virlit{«У йога, що перебуває в тілі, [є] стани, змушені [до] дії-й-недії; через відсутність з'єднання й роз'єднання я сиджу, як мені завгодно.»}
\virgram{karma-naiskarmya-nirbandha-bhāvāḥ (m. nom. pl. «стани, змушені до дії-й-недії»: karma «дія» + naiskarmya «недіяння» + nirbandha «змушування, наполягання» + bhāva «стан») dehasthayoginaḥ (m. gen. sg. «у йога, що перебуває в тілі»: deha «тіло» + stha «той, що стоїть» + yogin). saṃyoga-ayoga-virahāt (m. abl. sg. «через відсутність з'єднання й роз'єднання»: saṃyoga «з'єднання» + ayoga «роз'єднання» + viraha «розлука, відсутність») aham āse yathāsukham.}
\vircom{Слово nirbandha («змушування») — ключове: доки є тіло, є й примус то до дії, то до недіяння, і жоден вибір не вільний. saṃyoga / ayoga («з'єднання» / «роз'єднання») — пара з мови йоги, де з'єднання зазвичай означає зв'язок з об'єктом; тут відкидається сама можливість такого зв'язку. Зверни увагу на порядок: спершу визнано примус (він є), потім названо те, чого немає (з'єднання). Вірш не обіцяє свободи від тіла — лише від рахунку. РЕФРЕН ГЛАВИ. Шість із семи віршів (13.1, 13.3–13.7) закінчуються однаково: aham āse yathāsukham («я сиджу — як мені завгодно»). ās — не «бути», а саме СИДІТИ, перебувати на місці; yathāsukham — «як зручно, за бажанням». Разом це формула, яка скасовує будь-яку поставу як вимогу: не треба ні стояти, ні лежати правильно. 13.2 — єдиний вірш глави без цього рефрену, і його треба читати окремо.}
\virvars{\virvaritem{SARIT і GRETIL (перевірено scripts/fetch\_witnesses.py)}{Читання збігаються; відмінності лише в аваґрасі та поділі слів. Не різночитання.}}
\virsources{Avg\_13.4}{verse\_13.4}


\virs{13.5}{Ні користі, ні шкоди немає мені від стояння, ходіння чи лежання — тому, стоячи, ходячи чи сплячи, я сиджу, як мені завгодно.}
\virlit{«Ні користі, ні шкоди немає в мене від стояння, ходіння чи лежання — тому, стоячи, ходячи, сплячи, я сиджу, як мені завгодно.»}
\virgram{artha-anarthau (m. nom. du. «користь і шкода»: artha «користь» + anartha «лихо, шкода») na me («немає в мене») sthityā (f. instr. sg. «стоянням»: sthiti) gatyā (f. instr. sg. «ходінням, рухом»: gati) na śayanena (n. instr. sg. «лежанням, сном»: śayana) vā («або»). tiṣṭhan (m. nom. sg. дієприкметника теперішнього часу sthā «стоячи») gacchan (m. nom. sg. дієприкметника gam «ходячи») svapan (m. nom. sg. дієприкметника svap «сплячи») tasmāt («тому») aham āse yathāsukham.}
\vircom{Найгостріший парадокс глави, і він у самій граматиці: три дієприкметники (tiṣṭhan, gacchan, svapan — «стоячи, ходячи, сплячи») стоять поруч із дієсловом āse («сиджу»). Формально це суперечність, і вона навмисна: ās тут уже не про поставу тіла, а про перебування взагалі. Пара artha / anartha («користь» / «шкода») — та сама, що в 7.2 vṛddhi / kṣati, і знову відкидається цілою: не «немає шкоди», а немає ні того, ні того. РЕФРЕН ГЛАВИ. Шість із семи віршів (13.1, 13.3–13.7) закінчуються однаково: aham āse yathāsukham («я сиджу — як мені завгодно»). ās — не «бути», а саме СИДІТИ, перебувати на місці; yathāsukham — «як зручно, за бажанням». Разом це формула, яка скасовує будь-яку поставу як вимогу: не треба ні стояти, ні лежати правильно. 13.2 — єдиний вірш глави без цього рефрену, і його треба читати окремо.}
\virvars{\virvaritem{SARIT і GRETIL (перевірено scripts/fetch\_witnesses.py)}{Читання збігаються; відмінності лише в аваґрасі та поділі слів. Не різночитання.}}
\virsources{Avg\_13.5}{verse\_13.5}


\virs{13.6}{У сплячого немає мені втрати, як і немає досягнення в того, хто докладає зусиль; відкинувши загибель і піднесення, я сиджу, як мені завгодно.}
\virlit{«У сплячого немає в мене втрати, [як] і немає досягнення в того, хто докладає зусиль; відкинувши загибель і піднесення, я сиджу, як мені завгодно.»}
\virgram{svapataḥ (m. gen. sg. дієприкметника svap «у сплячого») na asti me hāniḥ (f. nom. sg. «немає в мене втрати»). siddhiḥ (f. nom. sg. «досягнення, успіх») yatnavataḥ (m. gen. sg. «у того, хто докладає зусиль»: yatna «зусилля») na vā («і немає»). nāśa-ullāsau (m. nom. du. «загибель і піднесення»: nāśa «загибель» + ullāsa «піднесення, радісний підйом») vihāya (герунд «відкинувши») asmāt aham āse yathāsukham.}
\vircom{Вірш будується на запереченні двох протилежних чеснот: сон не шкодить, зусилля не дає. Тим самим знімається сама вісь, на якій тримається будь-яка практика, — «старанність». ullāsa («піднесення») — слово з відтінком радості, не просто зростання; у парі з nāśa («загибель») воно описує не події, а те, як людина їх переживає. Порівняй 3.9, де похвала й мука були описані однією граматичною формою: тут той самий хід, але на рівні оцінки. РЕФРЕН ГЛАВИ. Шість із семи віршів (13.1, 13.3–13.7) закінчуються однаково: aham āse yathāsukham («я сиджу — як мені завгодно»). ās — не «бути», а саме СИДІТИ, перебувати на місці; yathāsukham — «як зручно, за бажанням». Разом це формула, яка скасовує будь-яку поставу як вимогу: не треба ні стояти, ні лежати правильно. 13.2 — єдиний вірш глави без цього рефрену, і його треба читати окремо.}
\virvars{\virvaritem{SARIT і GRETIL (перевірено scripts/fetch\_witnesses.py)}{Читання збігаються; відмінності лише в аваґрасі та поділі слів. Не різночитання.}}
\virsources{Avg\_13.6}{verse\_13.6}

\index{sukha@\textit{sukha}}
\virs{13.7}{Побачивши не раз у всіх станах те правило, що має вигляд задоволення, — відкинувши добре й погане, я сиджу, як мені завгодно.}
\virlit{«Побачивши не раз у [всіх] станах правило, що має вигляд задоволення, — відкинувши добре й погане, я сиджу, як мені завгодно.»}
\virgram{sukhādirūpā (f. nom. sg. «та, чия подоба — задоволення тощо»: sukha «задоволення» + ādi «тощо» + rūpa «подоба»; узгодження проблемне, див. коментар) niyamam (m./n. acc. sg. «правило, припис») bhāveṣu (m. loc. pl. «у станах») ālokya (герунд ā-lok «побачивши, оглянувши») bhūriśaḥ («багаторазово, щедро, рясно»). śubhāśubhe (n. nom./acc. du. «добре й погане») vihāya (герунд «відкинувши») asmāt aham āse yathāsukham.}
\vircom{Найважче місце глави, і не через сенс, а через текст. Свідки розходяться в одному слові: SARIT подає aniyamaṃ («безладдя, відсутність правила»), GRETIL — niyamaṃ («правило»). Сенс при цьому стає ПРОТИЛЕЖНИМ: побачити, що все, що має вигляд задоволення, насправді безладдя — або що воно підлягає правилу. Крім того, форма sukhādirūpā (жіночий рід, однина) не узгоджується ні з niyamaṃ/aniyamaṃ (середній рід), ні з bhāveṣu (множина, чоловічий/середній). Найімовірніше, обидва свідки успадкували зіпсоване читання, і правильним було sukhādirūpāṇi (середній рід, множина) — «речі, чия подоба — задоволення тощо». У перекладі взято читання SARIT, але це ВИБІР, а не розв'язання: потрібен третій свідок. РЕФРЕН ГЛАВИ. Шість із семи віршів (13.1, 13.3–13.7) закінчуються однаково: aham āse yathāsukham («я сиджу — як мені завгодно»). ās — не «бути», а саме СИДІТИ, перебувати на місці; yathāsukham — «як зручно, за бажанням». Разом це формула, яка скасовує будь-яку поставу як вимогу: не треба ні стояти, ні лежати правильно. 13.2 — єдиний вірш глави без цього рефрену, і його треба читати окремо.}
\virvars{\virvaritem{SARIT і GRETIL (перевірено scripts/fetch\_witnesses.py)}{СПРАВЖНЄ РІЗНОЧИТАННЯ З ПРОТИЛЕЖНИМ ЗМІСТОМ: SARIT — sukhādirūpā+aniyamaṃ («безладдя»), GRETIL — sukhādirūpā niyamaṃ («правило»). Це не орфографія: ani-yama і ni-yama — різні слова з протилежним значенням. Додатково: форма sukhādirūpā (f. sg.) не узгоджується з жодним із двох читань (обидва n./m. sg. або n. pl.), що вказує на зіпсоване читання в обох свідках; ймовірне початкове — sukhādirūpāṇi. НЕ РОЗВ'ЯЗАНО: третього свідка для 13.7 не перевірено. У перекладі взято SARIT і це позначено.}}
\virsources{Avg\_13.7}{verse\_13.7}

\pilotchapter{14}{4}
\index{bahuvrihi@\textit{bahuvrīhi}}
\virs{14.1}{Хто за природою має порожній розум — навіть якщо з недбальства плекає [в собі] стани, — той наче заснув і прокинувся: блукання його вичерпане.}
\virlit{«Хто за природою має порожній розум — хоч би з недбальства він і плекав стани — той наче заснулий і розбуджений: у нього вичерпане блукання.»}
\virgram{prakṛtyā (f. instr. sg. «за природою, від природи») śūnyacittaḥ (bahuvrīhi «той, чий розум порожній»: śūnya «порожній» + citta) yaḥ (m. nom. sg.). pramādāt (m. abl. sg. «через недбальство, з безпечності»: pra-māda) bhāvabhāvanaḥ (m. nom. sg. «той, хто плекає стани, вирощує почуття»: bhāva «стан» + bhāvana «вирощування, плекання»). nidritaḥ (m. nom. sg. «заснулий») bodhitaḥ (m. nom. sg. «розбуджений») iva («наче»). kṣīṇasaṃsaraṇaḥ (bahuvrīhi «той, у кого вичерпане блукання [сансарою]»: kṣīṇa «вичерпаний» + saṃsaraṇa «блукання, перехід») hi («адже») saḥ (m. nom. sg.).}
\vircom{Вірш навмисно знімає вимогу старанності. pramādāt («через недбальство») — це не похвала, але й не докір: недбальство нічого не змінює, бо порожнеча розуму тут не досягнута, а природна (prakṛtyā). Складене bhāvabhāvana («плекання станів») описує звичну роботу розуму — і вона названа марною. nidritaḥ / bodhitaḥ iva («наче заснулий і розбуджений») — не два стани, а один, побачений із двох боків: людина прокинулась, але ще не ворушиться. БУДОВА ГЛАВИ. Глава 14 не має спільного рефрену — на відміну від 5, 6 і 13. Натомість усі чотири вірші тримаються однієї думки, поданої з чотирьох боків: хто такий цей чоловік зсередини (14.1), що для нього зникає (14.2), про що він уже не турбується (14.3) і як його бачать ззовні (14.4).}
\virvars{\virvaritem{SARIT і GRETIL (перевірено scripts/fetch\_witnesses.py)}{СПРАВЖНЯ РОЗБІЖНІСТЬ, і вона змінює сенс: SARIT — kṣīṇasaṃsaraṇo («той, у кого вичерпане БЛУКАННЯ»), GRETIL — kṣīṇasaṃsmaraṇo («той, у кого вичерпане ПАМ'ЯТАННЯ»). saṃsaraṇa і saṃsmaraṇa — різні слова: перше про сансару й перехід між народженнями, друге про згадування. У контексті глави, де йдеться про вичерпаність шляху, читається SARIT. Яке читання правильне, НЕ ВСТАНОВЛЕНО: третього свідка для 14.1 не перевірено. У перекладі взято SARIT і це позначено як вибір.}}
\virsources{Avg\_14.1}{verse\_14.1}


\virs{14.2}{Де багатства, де друзі, де мої слуги — чуттєві предмети? Де писання, де знання — коли в мене розтанула спрага?}
\virlit{«Де багатства, де друзі, де мої раби-предмети? Де писання, де знання — коли в мене розтанула спрага?»}
\virgram{kva ... kva ... («де... де...»; питальне прислівник місця, повторене як фігура зникнення). dhanāni (n. nom. pl. «багатства») mitrāṇi (n. nom. pl. «друзі»). me (gen. sg.) viṣayadasyavaḥ (m. nom. pl. «раби чуттєвих предметів»: viṣaya «предмет чуття» + dasyu/dāsya «раб, слуга»). śāstram (n. nom. sg. «писання, наука») ca vijñānam (n. nom. sg. «знання, розпізнавання»). yadā («коли») me galitā (f. nom. sg. дієприкметника минулого часу gal «розтанула, вислизнула, опала») spṛhā (f. nom. sg. «спрага, жадоба»).}
\vircom{П'ять разів kva («де») — і жодної відповіді: питання не шукає місця, а показує, що місця немає. Зверни увагу на viṣayadasyavaḥ («раби предметів»): предмети названі не панами, а слугами, і це перевертає звичну картину — не вони тримають людину, а вона їх утримує. galitā («розтанула, опала») — слово про те, що сходить само, як пух або як лушпиння; спрага не переможена, вона зійшла. БУДОВА ГЛАВИ. Глава 14 не має спільного рефрену — на відміну від 5, 6 і 13. Натомість усі чотири вірші тримаються однієї думки, поданої з чотирьох боків: хто такий цей чоловік зсередини (14.1), що для нього зникає (14.2), про що він уже не турбується (14.3) і як його бачать ззовні (14.4).}
\virvars{\virvaritem{SARIT і GRETIL (перевірено scripts/fetch\_witnesses.py)}{Читання збігаються; відмінності лише в аваґрасі та поділі слів. Не різночитання.}}
\virsources{Avg\_14.2}{verse\_14.2}

\index{atman@\textit{ātman}}\index{saksin@\textit{sākṣin}}
\virs{14.3}{Коли пізнано пурушу-свідка, вищого Атмана й Ішвару — а також безнадію, узу й визволення — у мене немає турботи про визволення.}
\virlit{«Коли пізнано пурушу-свідка, вищого Атмана й Ішвару, [а також] безнадію, узу й визволення — у мене немає турботи про визволення.»}
\virgram{vijñāte (m. loc. sg. дієприкметника минулого часу vi-jñā «коли пізнано») sākṣipuruṣe (m. loc. sg. «пуруша-свідок»: sākṣin «свідок» + puruṣa «пуруша») paramātmani (m. loc. sg. «вищий Атман») ca īśvare (m. loc. sg. «Ішвара, володар»). nairāśye (n. loc. sg. «безнадія, відсутність сподівань») bandhamokṣe (m. loc. du./sg. «уза й визволення») ca. na cintā (f. nom. sg. «немає турботи, думки») muktaye (f. dat. sg. «про визволення») mama (gen. sg. «у мене»).}
\vircom{Тут уперше в книзі з'являється складене sākṣipuruṣa — «пуруша-свідок». Це не те саме, що просто ātman: sākṣin додає саме те, що свідок не бере участі, а лише присутній при. Найтонше — перелік: у ньому стоять поруч Ішвара і nairāśya («безнадія»), уза й визволення, тобто протилежності названі в одному ряді й пізнані разом. Тому й na cintā muktaye («немає турботи про визволення»): коли визволення пізнане як одне з пари, турбуватися про нього нема про що. БУДОВА ГЛАВИ. Глава 14 не має спільного рефрену — на відміну від 5, 6 і 13. Натомість усі чотири вірші тримаються однієї думки, поданої з чотирьох боків: хто такий цей чоловік зсередини (14.1), що для нього зникає (14.2), про що він уже не турбується (14.3) і як його бачать ззовні (14.4).}
\virvars{\virvaritem{SARIT і GRETIL (перевірено scripts/fetch\_witnesses.py)}{Читання збігаються; відмінності лише в аваґрасі та поділі слів. Не різночитання.}}
\virsources{Avg\_14.3}{verse\_14.3}

\index{bahuvrihi@\textit{bahuvrīhi}}
\virs{14.4}{У того, хто всередині порожній від розрізнень, а зовні ходить, як сам захоче, стани бувають усякі, наче в божевільного; і знають його тільки такі самі, як він.}
\virlit{«У того, хто всередині порожній від розрізнень, а зовні ходить за власною волею, — стани [бувають] різні, наче в божевільного; і знають його лише такі ж, як він.»}
\virgram{antaḥ-vikalpa-śūnyasya (bahuvrīhi «той, хто всередині порожній від розрізнень»: antaḥ «усередині» + vikalpa «розрізнення, побудова» + śūnya «порожній») bahiḥ («зовні») svacchandacāriṇaḥ (m. gen. sg. «того, хто ходить за власною волею»: sva-cchanda «власна воля» + cārin «той, хто ходить»). bhrāntasya (m. gen. sg. «блукаючого, божевільного, збентеженого») iva («наче»). daśāḥ (f. nom. pl. «стани, положення») tāḥ tāḥ («ті й ті, різні») tādṛśāḥ (m. nom. pl. «такі ж, подібні») eva («лише») jānate (3 pl. pr. jñā «знають»).}
\vircom{bhrānta означає і «заблукалий», і «божевільний» — двозначність тут не прикраса, а точність: ззовні така людина не відрізняється від божевільної, бо різниця всередині, а всередину не видно. tādṛśāḥ eva («лише такі ж») закриває тему пізнання: зрозуміти його може тільки той, хто в такому самому стані, тобто свідчення ззовні неможливе. Це прямий виклик будь-якому зовнішньому описові, і він поставлений останнім віршем глави. БУДОВА ГЛАВИ. Глава 14 не має спільного рефрену — на відміну від 5, 6 і 13. Натомість усі чотири вірші тримаються однієї думки, поданої з чотирьох боків: хто такий цей чоловік зсередини (14.1), що для нього зникає (14.2), про що він уже не турбується (14.3) і як його бачать ззовні (14.4).}
\virvars{\virvaritem{SARIT і GRETIL (перевірено scripts/fetch\_witnesses.py)}{Розбіжність у написанні: SARIT — daśāstāstāstādṛśā (злитно, з подвоєнням tās), GRETIL — daśās tās tās tādṛśā (окремо). Слова ті самі: daśāḥ + tāḥ + tāḥ + tādṛśāḥ. Форма SARIT виникла від злитного запису, сенс не змінюється. Не різночитання, але SARIT-написання важче читати; у граматиці розібрано за GRETIL.}}
\virsources{Avg\_14.4}{verse\_14.4}

\pilotchapter{19}{8}

\virs{19.1}{Взявши кліщі знання істини, я витяг із нутра серця колючки різноманітних хибних уявлень.}
\virlit{«Взявши кліщі знання істини, я витяг з нутра серця колючки різноманітних хибних уявлень.»}
\virgram{tattvavijñānasandaṃśam (m. acc. sg. «кліщі знання істини»: tattva «істина, суть» + vijñāna «знання» + sandaṃśa «кліщі, щипці»; санскритське sandaṃśa буквально «те, що кусає разом» — так само, як українське «щипці» від «щипати») ādāya (герунд ā-dā «взявши») hṛdayodarāt (n. abl. sg. «з нутра серця»: hṛdaya «серце» + udara «нутро, черево»). nānāvidhaparāmarśaśalyoddhāraḥ (m. nom. sg. «витягання колючок різноманітних хибних схоплень»: nānāvidha «різноманітний» + parāmarśa «схоплювання, упереджене уявлення» + śalya «колючка, стріла» + uddhāra «витягання») kṛtaḥ (m. nom. sg. «зроблено») mayā (instr. sg. «мною»).}
\vircom{Вірш — не декларація, а опис процедури, і метафора тут хірургічна: sandaṃśa — кліщі, якими виймають стороннє тіло з рани. Три слова працюють разом: śalya («колючка») — те, що застрягло; uddhāra («витягання») — дія; hṛdayodara («нутро серця») — місце, і воно назване нутром, а не поверхнею. parāmarśa тут не «думка», а хибне схоплювання — розумове прийняття чогось за інше. Тобто йдеться не про прибирання думок, а про виймання того, що застрягло й болить. РЕФРЕН ГЛАВИ. 19.2–19.7 закінчуються однаково: svamahimni sthitasya me («у мене, що стоїть у власній величі»). Кожен із цих шести віршів — низка питань kva («де»), і жодне не дістає відповіді. 19.1 і 19.8 стоять поза серією: перший описує, ЯК це сталося, останній — ЩО з цього випливає для мови.}
\virvars{\virvaritem{SARIT і GRETIL (перевірено scripts/fetch\_witnesses.py)}{Читання збігаються; відмінності лише в аваґрасі та поділі слів. Не різночитання.}}
\virsources{Avg\_19.1}{verse\_19.1}


\virs{19.2}{Де дгарма, де бажання, де користь, де вміння розрізняти? Де двоїстість, де недвоїстість — у мене, що стоїть у власній величі?}
\virlit{«Де дгарма, де бажання, де користь, де розрізнювальність? Де двоїстість, де недвоїстість — у мене, що стоїть у власній величі?»}
\virgram{kva ... kva («де... де»; повторюване питання). dharmaḥ (m. nom. sg. «дгарма, обов'язок») ca vā kāmaḥ (m. nom. sg. «бажання, насолода») ca arthaḥ (m. nom. sg. «користь») vivekatā (f. nom. sg. «розрізнювальність, здатність розрізняти»). dvaitam (n. nom. sg. «двоїстість») ca vā advaitam (n. nom. sg. «недвоїстість»). svamahimni (m./n. loc. sg. «у власній величі»: sva «власний» + mahiman «велич») sthitasya (m. gen. sg. «того, хто стоїть») me (gen. sg. «у мене»).}
\vircom{Перший із шести віршів-переліків. Три з перших чотирьох понять — це три з чотирьох пурушартг (дгарма, користь, бажання); четверте, визволення, не назване, бо його місце займає сама позиція того, хто питає. Найважливіше — остання пара: питається не лише «де двоїстість», а й «де недвоїстість». Тобто заперечується не світ на користь єдності, а сама вісь, на якій обидва поняття мають сенс. svamahimni («у власній величі») — не гордість, а міра: те, що не має ні межі, ні місця. РЕФРЕН ГЛАВИ. 19.2–19.7 закінчуються однаково: svamahimni sthitasya me («у мене, що стоїть у власній величі»). Кожен із цих шести віршів — низка питань kva («де»), і жодне не дістає відповіді. 19.1 і 19.8 стоять поза серією: перший описує, ЯК це сталося, останній — ЩО з цього випливає для мови.}
\virvars{\virvaritem{SARIT і GRETIL (перевірено scripts/fetch\_witnesses.py)}{Читання збігаються; відмінності лише в аваґрасі та поділі слів. Не різночитання.}}
\virsources{Avg\_19.2}{verse\_19.2}


\virs{19.3}{Де минуле, де майбутнє, де навіть теперішнє? Де сторона світу, де вічність — у мене, що стоїть у власній величі?}
\virlit{«Де минуле, де майбутнє, де навіть теперішнє? Де місце, де вічність — у мене, що стоїть у власній величі?»}
\virgram{kva bhūtam (n. nom. sg. «минуле, те, що було») kva bhaviṣyat (n. nom. sg. «майбутнє») vā vartamānam (n. nom. sg. «теперішнє») api kva vā. kva deśaḥ (m. nom. sg. «місце, сторона») kva ca vā nityam (n. nom. sg. «вічне, постійне»). svamahimni sthitasya me.}
\vircom{Три часи названі поспіль, і слово api («навіть») стоїть саме при теперішньому: минуле й майбутнє відсутні легко, а теперішнє — те, що здається найбільш незаперечним. Друга пара — deśa («місце») і nitya («вічне») — зіставляє простір із часом; разом обидві пари охоплюють усе, у чому можна було б розміститися. Вірш показує, що «стояти у власній величі» означає не займати іншого місця, а не мати місця взагалі. РЕФРЕН ГЛАВИ. 19.2–19.7 закінчуються однаково: svamahimni sthitasya me («у мене, що стоїть у власній величі»). Кожен із цих шести віршів — низка питань kva («де»), і жодне не дістає відповіді. 19.1 і 19.8 стоять поза серією: перший описує, ЯК це сталося, останній — ЩО з цього випливає для мови.}
\virvars{\virvaritem{SARIT і GRETIL (перевірено scripts/fetch\_witnesses.py)}{Читання збігаються; відмінності лише в аваґрасі та поділі слів. Не різночитання.}}
\virsources{Avg\_19.3}{verse\_19.3}

\index{tatha@\textit{tathā}}
\virs{19.4}{Де Атман, де не-Атман, де добре, де недобре? Де думка, де відсутність думки — у мене, що стоїть у власній величі?}
\virlit{«Де Атман, де не-Атман, де добре, де недобре? Де думка, де не-думка — у мене, що стоїть у власній величі?»}
\virgram{kva ca ātmā (m. nom. sg. «Атман») kva ca vā anātmā (m. nom. sg. «не-Атман, те, що не є Атманом»). kva śubham (n. nom. sg. «добре, сприятливе») kva aśubham (n. nom. sg. «недобре») tathā («так само»). kva cintā (f. nom. sg. «думка, турбота») kva ca vā acintā (f. nom. sg. «не-думка, відсутність думки»). svamahimni sthitasya me.}
\vircom{Три пари, і в кожній заперечується не один бік, а обидва. Найгостріша — перша: питати «де Атман» і «де не-Атман» однаково означає не знайти місця ні для того, ні для того, тобто сама пара втрачає ґрунт. Друга пара (śubha / aśubha) — та сама, що в 13.7, де вона відкидалася цілою. Третя (cintā / acintā) важлива тим, що сюди потрапляє й «відсутність думки»: навіть спокій розуму не стає позицією, з якої можна стояти. РЕФРЕН ГЛАВИ. 19.2–19.7 закінчуються однаково: svamahimni sthitasya me («у мене, що стоїть у власній величі»). Кожен із цих шести віршів — низка питань kva («де»), і жодне не дістає відповіді. 19.1 і 19.8 стоять поза серією: перший описує, ЯК це сталося, останній — ЩО з цього випливає для мови.}
\virvars{\virvaritem{SARIT і GRETIL (перевірено scripts/fetch\_witnesses.py)}{СПРАВЖНЯ РОЗБІЖНІСТЬ у двох місцях. Перше: SARIT — kvāśubhaṃ («де НЕДОБРЕ»), GRETIL — kvāsubhaṃ («де НЕЖИВЕ, позбавлене життя»); aśubha і asubha — різні слова. Друге: SARIT — tathā («так само»), GRETIL — yathā («як»). У контексті пари «добре / недобре» читається SARIT; asubha («неживе») випадає з низки понять глави. Яке читання правильне, НЕ ВСТАНОВЛЕНО: третього свідка для 19.4 не перевірено. У перекладі взято SARIT і це позначено як вибір.}}
\virsources{Avg\_19.4}{verse\_19.4}

\index{tatha@\textit{tathā}}
\virs{19.5}{Де сон зі сновидіннями, де глибокий сон, де неспання? Де четвертий стан, де навіть страх — у мене, що стоїть у власній величі?}
\virlit{«Де сон, де глибокий сон, де неспання? Де четвертий [стан], де навіть страх — у мене, що стоїть у власній величі?»}
\virgram{kva svapnaḥ (m. nom. sg. «сон зі сновидіннями») kva suṣuptiḥ (f. nom. sg. «глибокий сон без снів») vā kva ca jāgaraṇam (n. nom. sg. «неспання, бадьорість») tathā. kva turīyam (n. nom. sg. «четвертий [стан]», турия — те, що поза трьома) bhayam (n. nom. sg. «страх») vā api. svamahimni sthitasya me.}
\vircom{Три стани — неспання, сон зі сновидіннями, глибокий сон — це стандартна схема упанішад, і четвертий (turīya) стоїть поза нею. Тут відкидається й він: навіть те, що мало б бути виходом, не має місця. Страх (bhaya) доданий останнім, і це не випадково: страх — те, що тримає в межах трьох станів, і якщо його немає, то й виходити нікуди. РЕФРЕН ГЛАВИ. 19.2–19.7 закінчуються однаково: svamahimni sthitasya me («у мене, що стоїть у власній величі»). Кожен із цих шести віршів — низка питань kva («де»), і жодне не дістає відповіді. 19.1 і 19.8 стоять поза серією: перший описує, ЯК це сталося, останній — ЩО з цього випливає для мови.}
\virvars{\virvaritem{SARIT і GRETIL (перевірено scripts/fetch\_witnesses.py)}{Розбіжність turīyaṃ (SARIT) / turiyaṃ (GRETIL) — та сама форма, різне позначення довгого ī, орфографія. Не різночитання.}}
\virsources{Avg\_19.5}{verse\_19.5}


\virs{19.6}{Де далеке, де близьке, де зовнішнє, де внутрішнє? Де грубе, де тонке — у мене, що стоїть у власній величі?}
\virlit{«Де далеке, де близьке, де зовнішнє, де внутрішнє? Де грубе, де тонке — у мене, що стоїть у власній величі?»}
\virgram{kva dūram (n. nom. sg. «далеке») kva samīpam (n. nom. sg. «близьке») vā bāhyam (n. nom. sg. «зовнішнє») kva abhyantaram (n. nom. sg. «внутрішнє») kva vā. kva sthūlam (n. nom. sg. «грубе, велике») kva ca vā sūkṣmam (n. nom. sg. «тонке, дрібне»). svamahimni sthitasya me.}
\vircom{Три пари, і кожна описує простір по-своєму: відстань (далеке / близьке), напрямок (зовнішнє / внутрішнє), щільність (грубе / тонке). Разом вони мали б вичерпати всі способи бути десь — і саме тому всі три відкидаються разом. Пара sthūla / sūkṣma особлива: це не лише «велике й мале», а й «тілесне й тонке», тобто остання межа між видимим і невидимим. РЕФРЕН ГЛАВИ. 19.2–19.7 закінчуються однаково: svamahimni sthitasya me («у мене, що стоїть у власній величі»). Кожен із цих шести віршів — низка питань kva («де»), і жодне не дістає відповіді. 19.1 і 19.8 стоять поза серією: перший описує, ЯК це сталося, останній — ЩО з цього випливає для мови.}
\virvars{\virvaritem{SARIT і GRETIL (перевірено scripts/fetch\_witnesses.py)}{Розбіжність bāhyaṃ (SARIT) / bahyaṃ (GRETIL) — та сама форма, різне позначення довгого ā, орфографія. Не різночитання.}}
\virsources{Avg\_19.6}{verse\_19.6}


\virs{19.7}{Де смерть, де життя, де світи, де мирське? Де розчинення, де самадгі — у мене, що стоїть у власній величі?}
\virlit{«Де смерть, де життя, де світи, де мирське? Де розчинення, де самадгі — у мене, що стоїть у власній величі?»}
\virgram{kva mṛtyuḥ (m. nom. sg. «смерть») jīvitam (n. nom. sg. «життя») vā kva lokāḥ (m. nom. pl. «світи») kva asya («де цього»; форма темна, див. коментар) kva laukikam (n. nom. sg. «мирське, світське»). kva layaḥ (m. nom. sg. «розчинення») kva samādhiḥ (m. nom. sg. «самадгі, зосередження») vā. svamahimni sthitasya me.}
\vircom{Останній із шести переліків, і він забирає найвище: laya і samādhi — те, до чого глави 5 і 13 кликали, тут теж не мають місця. Це замикає книгу на собі: навіть мета попередніх глав не стає притулком. Форма kvāsya (kva + asya, «де цього») темна: asya — родового від займенника, але без опори в реченні; можливо, це зіпсоване читання. У перекладі цей член пропущено, і це позначено, а не замасковано. РЕФРЕН ГЛАВИ. 19.2–19.7 закінчуються однаково: svamahimni sthitasya me («у мене, що стоїть у власній величі»). Кожен із цих шести віршів — низка питань kva («де»), і жодне не дістає відповіді. 19.1 і 19.8 стоять поза серією: перший описує, ЯК це сталося, останній — ЩО з цього випливає для мови.}
\virvars{\virvaritem{SARIT і GRETIL (перевірено scripts/fetch\_witnesses.py)}{Читання збігаються в усьому, крім поділу слів (samādhirvā / samādhir vā). Окремо: форма kvāsya (kva + asya) у САМОМУ тексті темна за змістом — це не розбіжність між свідками, а місце, яке потребує третього свідка або консультації. Позначено в граматиці й у перекладі.}}
\virsources{Avg\_19.7}{verse\_19.7}


\virs{19.8}{Досить розмов про трицінність, досить і розмов про йогу, досить розмов про знання — у мене, що спочив в Атмані.}
\virlit{«Досить розмов про трицінність, досить і розмов про йогу, досить розмов про знання — у мене, що спочив в Атмані.»}
\virgram{alam («досить, годі»; з instr. «досить про») trivargakathayā (f. instr. sg. «розмовою про трицінність»: tri-varga «трійця [дгарма, користь, бажання]» + kathā «розмова»). yogasya (m. gen. sg. «йоги») kathayā (f. instr. sg. «розмовою») api alam. alam vijñānakathayā (f. instr. sg. «розмовою про знання»). viśrāntasya (m. gen. sg. дієприкметника минулого часу vi-śram «того, хто спочив, відпочив») mama (gen. sg.) ātmani (m. loc. sg. «в Атмані»).}
\vircom{Вірш не має рефрену попередніх шести й замінює його іншим: не svamahimni sthitasya («що стоїть у власній величі»), а viśrāntasya mamātmani («що спочив в Атмані»). Різниця точна: стояння — це позиція, спочинок — її відсутність потреби. Тричі повторене alam («досить») знищує три теми поспіль: етику (трицінність), практику (йога) і теорію (знання). Це останній вірш книги перед главой 20, і він закриває розмову — буквально: kathā означає «розмова», і саме розмову тут оголошено зайвою.}
\virvars{\virvaritem{SARIT і GRETIL (перевірено scripts/fetch\_witnesses.py)}{Читання збігаються; відмінності лише в аваґрасі та поділі слів. Не різночитання.}}
\virsources{Avg\_19.8}{verse\_19.8}



% =====================================================================
%  АПАРАТ
% =====================================================================
\backmatter
\chapter*{Покажчик термінів}
\addcontentsline{toc}{chapter}{Покажчик термінів}
\markboth{Покажчик термінів}{Покажчик термінів}
\printindex

\chapter*{Словник термінів}
\addcontentsline{toc}{chapter}{Словник термінів}
\markboth{Словник термінів}{Словник термінів}
\noindent{\sffamily\small Робочі відповідники з \texttt{editorial/}\allowbreak\texttt{GLOSSARY.md}.
Вони підлягають санскритологічній і літературній перевірці; це не
нормативний словник.}

\vspace{0.8em}
% ЗГЕНЕРОВАНО scripts/build_pilot_book.py — не редагувати вручну.
\begin{longtable}{@{}p{0.20\textwidth}p{0.30\textwidth}p{0.40\textwidth}@{}}
\textbf{\sffamily Санскрит} & \textbf{\sffamily Робочий відповідник} & \textbf{\sffamily Застереження} \\
\hline
\endhead
\textit{ātman} & Я / істинне Я & Не тотожне повсякденному его \\
\textit{cit} & свідомість / усвідомлення & Значення визначає контекст \\
\textit{sākṣin} & свідок & Філософський, а не судовий образ \\
\textit{mukti} & звільнення / визволення & Обрати послідовний термін у виданні \\
\textit{vairāgya} & неприв'язаність & Не підміняти байдужістю \\
\textit{toṣa} & вдоволеність & Важлива пропущена чеснота у посередньому перекладі 1.2 \\
\textit{āśrama} & стадія / стан життя & Не зводити до «обов'язків» \\
\textit{varṇa} & суспільна категорія & Відрізняти від життєвого стану \\
\end{longtable}


\chapter*{Джерела й межі цього видання}
\addcontentsline{toc}{chapter}{Джерела й межі цього видання}
\markboth{Джерела й межі}{Джерела й межі}

\noindent\textbf{\sffamily Що є джерелом нашого тексту.}
Власні редакторські картки \texttt{editorial/}\allowbreak\texttt{cards/*.json}: поля
\texttt{ukrainian\_translation}, \texttt{literal\_meaning}, \texttt{grammar},
\texttt{commentary} і \texttt{variants}. Це наш авторський внесок під
ліцензією «ВОЛЬНІСТЬ».

\vspace{0.6em}
\noindent\textbf{\sffamily Що ми не відтворюємо.}
Транскрипції GRETIL і SARIT, їхню розмітку та редакторський апарат. Локатори
(\texttt{Avg\_2.1}, \texttt{verse\_2.1}) подано як покажчик для звірки.

\vspace{0.6em}
\noindent\textbf{\sffamily Чого ми не стверджуємо.}
Що це критичне видання. Що переклад вивірений. Що коментар рецензований. Що
права на всі використані матеріали остаточно з'ясовані.

\vspace{0.6em}
\noindent\textbf{\sffamily Статус.}
Робоче пілотне видання. Картки — «чернетка». Незалежної рецензії немає.

\end{document}
